% Travail et puissance d'une force — Physique 1ère TI — Cours signé OnBuch+
% Programme officiel MINESEC. Compiler avec : tectonic N03-travail-puissance-force.tex
\newcommand{\DOCMATIERE}{Physique}
\newcommand{\DOCNIVEAU}{1ère TI}
\newcommand{\DOCMODULE}{Module 2 : Mouvements et interactions mécaniques}
\newcommand{\DOCLECON}{N03}
\newcommand{\DOCTITRE}{Travail et puissance d'une force}
% OnBuch+ — préambule « cours signés » ENRICHI (compilé avec tectonic / XeLaTeX)
% Système visuel « Pop » d'OnBuch+ : fond crème (jamais blanc), bordures encre
% épaisses, ombres « dures » décalées, titres Archivo Black en capitales.
% Version MATHÉMATIQUES : graphes (pgfplots/tikz), boîtes pédagogiques
% (définition, propriété, méthode, attention, le savais-tu, exercice résolu,
% exercice, TP, bilan), page de garde et signature OnBuch+ ancrée en bas.
%
% Macros attendues dans main.tex AVANT \input{preamble} :
%   \DOCMATIERE  \DOCNIVEAU  \DOCMODULE  \DOCLECON (numéro)  \DOCTITRE
\documentclass[11pt]{article}

\usepackage{fontspec}
\usepackage[french]{babel}
\usepackage[a4paper,margin=2cm,top=3.4cm,bottom=2.8cm,headheight=1.4cm,footskip=1.3cm]{geometry}
\usepackage{amsmath,amssymb,amsfonts}
\usepackage{mathtools} % \xmapsto, \coloneqq... souvent utilisés par les modèles
\usepackage{cancel} % simplifications barrées (bilans de forces)
% \abs{z} (module, valeur absolue) et \norm{u} (norme), utilisés par les modèles
\providecommand{\abs}{}\let\abs\relax\DeclarePairedDelimiter{\abs}{\lvert}{\rvert}
\providecommand{\norm}{}\let\norm\relax\DeclarePairedDelimiter{\norm}{\lVert}{\rVert}
\usepackage[table]{xcolor} % [table] : fournit aussi \rowcolors (alternance de lignes)
\usepackage{pagecolor}
\usepackage{tikz}
\usetikzlibrary{arrows.meta,positioning,calc,shapes.geometric,shapes.misc,shapes.arrows,decorations.pathmorphing,decorations.pathreplacing,decorations.markings,patterns,shadows,mindmap,trees,fit,backgrounds,matrix,intersections,angles,quotes}
\usepackage{pgfplots}
\usepackage[version=4]{mhchem} % équations nucléaires \ce{^{235}_{92}U}
\usepackage[european,siunitx]{circuitikz} % schémas électriques
\pgfplotsset{compat=1.18}
\usepgfplotslibrary{fillbetween,groupplots} % aires entre courbes (calcul intégral)
\usepackage{enumitem}
\usepackage{fancyhdr}
% [most] charge toutes les bibliothèques tcolorbox (listings, minted, raster,
% documentation...) et gonfle inutilement la mémoire TeX sur un document long
% (~300 boîtes) : on ne charge que ce qui est réellement utilisé.
\usepackage[skins,breakable]{tcolorbox}
\usepackage{graphicx}
\usepackage{colortbl}
\usepackage{booktabs}
\usepackage{tabularx}
\usepackage{diagbox} % case d'en-tête diagonale des tableaux à double entrée
% Colonnes extensibles centrée / gauche / droite (C, L, R), souvent utilisées par les modèles
\newcolumntype{C}{>{\centering\arraybackslash}X}
\newcolumntype{L}{>{\raggedright\arraybackslash}X}
\newcolumntype{R}{>{\raggedleft\arraybackslash}X}
\usepackage{array}
\usepackage{multicol}
\usepackage{siunitx}
% parse-numbers=false : accepte aussi \SI{2,5 \times 10^{-3}}{...}
\sisetup{locale=FR,output-decimal-marker={,},group-minimum-digits=4,parse-numbers=false}
\DeclareSIUnit\ohm{\ensuremath{\symup{\Omega}}} % U+2126 absent de la police
\DeclareSIUnit\division{div} % réglages d'oscilloscope (ms/div, V/div)
\DeclareSIUnit\tr{tr} % tours (vitesse de rotation, ex. \SI{1500}{\tr\per\minute})
\DeclareSIUnit\ch{ch} % cheval-vapeur (puissance moteur, unité usuelle hors SI)
\DeclareSIUnit\franccfa{FCFA} % franc CFA (coûts, contexte camerounais)
\DeclareSIUnit\dioptre{\ensuremath{\delta}} % vergence d'une lentille (dioptrie)
\usepackage{qrcode}
% \degree : commande fréquemment utilisée par les modèles pour un angle ou une
% température en dehors de \SI{}{} (non fournie par les paquets chargés).
\providecommand{\degree}{^{\circ}}
% \qed (fin de démonstration), utilisé par les modèles sans amsthm
\providecommand{\qed}{\hfill$\square$}
% \barre : double barre des valeurs interdites dans un tableau de variations (tkz-tab)
\providecommand{\barre}{\ensuremath{\|}}
% environnement proof (amsthm non chargé) utilisé par les modèles
\ifdefined\proof\else\newenvironment{proof}[1][Démonstration]{\par\noindent\textit{#1.}\ }{\qed\par}\fi
\usepackage{lastpage}
\usepackage{hyperref}
\hypersetup{hidelinks}

% ============================================================
% Palette « Pop » OnBuch+ — mêmes hex que la classe P (plus_ui.dart)
% ============================================================
\definecolor{poporange}{HTML}{F59321}
\definecolor{popdark}{HTML}{E07A0C}
\definecolor{popcream}{HTML}{FFF6E8}
\definecolor{popink}{HTML}{1C1714}
\definecolor{poppurple}{HTML}{7A5AE0}
\definecolor{popgreen}{HTML}{2E9E6A}
\definecolor{poppink}{HTML}{E0457B}
\definecolor{popblue}{HTML}{2F7DE1}
\definecolor{popgold}{HTML}{C98A12}
\definecolor{popmuted}{HTML}{8A7F76}
\definecolor{popline}{HTML}{EADFD2}
\definecolor{popgray}{HTML}{8A7F76}
\colorlet{popgrey}{popgray}
% Alias tolérants (le modèle nomme parfois les couleurs différemment)
\colorlet{popwhite}{white}
\colorlet{popblack}{popink}
\colorlet{popred}{poppink}
\colorlet{popyellow}{popgold}
% Teintes claires pour fonds de tableaux / zones
\colorlet{popblueL}{popblue!10!white}
\colorlet{poporangeL}{poporange!14!white}
\colorlet{popgreenL}{popgreen!12!white}
\colorlet{poppurpleL}{poppurple!12!white}
\colorlet{poppinkL}{poppink!12!white}
\colorlet{popgoldL}{popgold!14!white}

\pagecolor{popcream}

% ============================================================
% Typographie
% ============================================================
\defaultfontfeatures{Ligatures=TeX}
\setmainfont{PlusJakartaSans-Medium.ttf}[Path=../../../fonts/,BoldFont=PlusJakartaSans-Bold.ttf]
% Maths en sans-serif (Fira Math) avec chiffres et lettres droites de Plus
% Jakarta, pour que les formules se fondent dans le texte.
\usepackage[math-style=ISO,bold-style=ISO]{unicode-math}
\setmathfont{FiraMath-Regular.otf}[Path=../../../fonts/]
\setmathfont{PlusJakartaSans-Medium.ttf}[Path=../../../fonts/,range={up/num,up/{Latin,latin},\mathup}]
\usepackage{icomma} % virgule décimale sans espace en mode math
% Caractères mathématiques Unicode absents des polices de texte (Plus Jakarta,
% Archivo Black) : rendus par leur équivalent mathématique, en texte comme en
% formule — sinon ils s'affichent en carré vide (ex. « Arithmétique dans ℤ »).
\usepackage{newunicodechar}
\newunicodechar{ℝ}{\ensuremath{\mathbb{R}}}
\newunicodechar{ℤ}{\ensuremath{\mathbb{Z}}}
\newunicodechar{ℂ}{\ensuremath{\mathbb{C}}}
\newunicodechar{ℕ}{\ensuremath{\mathbb{N}}}
\newunicodechar{ℚ}{\ensuremath{\mathbb{Q}}}
\newunicodechar{𝔻}{\ensuremath{\mathbb{D}}}
\newunicodechar{α}{\ensuremath{\alpha}}
\newunicodechar{β}{\ensuremath{\beta}}
\newunicodechar{γ}{\ensuremath{\gamma}}
\newunicodechar{δ}{\ensuremath{\delta}}
\newunicodechar{ε}{\ensuremath{\varepsilon}}
\newunicodechar{θ}{\ensuremath{\theta}}
\newunicodechar{λ}{\ensuremath{\lambda}}
\newunicodechar{μ}{\ensuremath{\mu}}
\newunicodechar{σ}{\ensuremath{\sigma}}
\newunicodechar{τ}{\ensuremath{\tau}}
\newunicodechar{φ}{\ensuremath{\varphi}}
\newunicodechar{ψ}{\ensuremath{\psi}}
\newunicodechar{ω}{\ensuremath{\omega}}
\newunicodechar{Σ}{\ensuremath{\Sigma}}
% majuscules produites par \MakeUppercase dans les titres (α-aminés -> Α)
\newunicodechar{Α}{\ensuremath{\alpha}}
\newunicodechar{Β}{\ensuremath{\beta}}
\newunicodechar{Γ}{\ensuremath{\Gamma}}
\newunicodechar{Δ}{\ensuremath{\Delta}}
\newunicodechar{Ε}{\ensuremath{\varepsilon}}
\newunicodechar{Θ}{\ensuremath{\Theta}}
\newunicodechar{Λ}{\ensuremath{\Lambda}}
\newunicodechar{Μ}{\ensuremath{\mu}}
\newunicodechar{Τ}{\ensuremath{\tau}}
\newunicodechar{Φ}{\ensuremath{\Phi}}
\newunicodechar{Ψ}{\ensuremath{\Psi}}
\newunicodechar{Ω}{\ensuremath{\Omega}}
\newunicodechar{↦}{\ensuremath{\mapsto}}
\newunicodechar{∈}{\ensuremath{\in}}
\newunicodechar{∉}{\ensuremath{\notin}}
\newunicodechar{∧}{\ensuremath{\wedge}}
\newunicodechar{⇔}{\ensuremath{\Leftrightarrow}}
\newunicodechar{⟺}{\ensuremath{\Longleftrightarrow}}
\newunicodechar{⇒}{\ensuremath{\Rightarrow}}
\newunicodechar{⟹}{\ensuremath{\Longrightarrow}}
\newunicodechar{∘}{\ensuremath{\circ}}
\newunicodechar{∪}{\ensuremath{\cup}}
\newunicodechar{∩}{\ensuremath{\cap}}
\newunicodechar{≡}{\ensuremath{\equiv}}
\newunicodechar{⊂}{\ensuremath{\subset}}
\newunicodechar{⊥}{\ensuremath{\perp}}
\newunicodechar{∃}{\ensuremath{\exists}}
\newunicodechar{∀}{\ensuremath{\forall}}
\newunicodechar{∛}{\ensuremath{\sqrt[3]{\,}}}
\newunicodechar{∜}{\ensuremath{\sqrt[4]{\,}}}
\newunicodechar{⃗}{\ensuremath{{}^{\to}}}
\newunicodechar{ⁿ}{\ifmmode^{n}\else\textsuperscript{\ensuremath{n}}\fi}
\newunicodechar{ˣ}{\ifmmode^{x}\else\textsuperscript{\ensuremath{x}}\fi}
\newunicodechar{ᵇ}{\ifmmode^{b}\else\textsuperscript{\ensuremath{b}}\fi}
\newunicodechar{ʳ}{\ifmmode^{r}\else\textsuperscript{\ensuremath{r}}\fi}
\newunicodechar{ᵘ}{\ifmmode^{u}\else\textsuperscript{\ensuremath{u}}\fi}
\newunicodechar{ʸ}{\ifmmode^{y}\else\textsuperscript{\ensuremath{y}}\fi}
\newunicodechar{ᵃ}{\ifmmode^{a}\else\textsuperscript{\ensuremath{a}}\fi}
\newunicodechar{ᵉ}{\ifmmode^{e}\else\textsuperscript{\ensuremath{e}}\fi}
\newunicodechar{ᵏ}{\ifmmode^{k}\else\textsuperscript{\ensuremath{k}}\fi}
\newunicodechar{ᵖ}{\ifmmode^{p}\else\textsuperscript{\ensuremath{p}}\fi}
\newunicodechar{ᴺ}{\ifmmode^{N}\else\textsuperscript{\ensuremath{N}}\fi}
\newunicodechar{ᶜ}{\ifmmode^{c}\else\textsuperscript{\ensuremath{c}}\fi}
\newunicodechar{ʰ}{\ifmmode^{h}\else\textsuperscript{\ensuremath{h}}\fi}
\newunicodechar{ᵠ}{\ifmmode^{q}\else\textsuperscript{\ensuremath{q}}\fi}
\newunicodechar{ₙ}{\ifmmode_{n}\else\textsubscript{\ensuremath{n}}\fi}
\newunicodechar{ₐ}{\ifmmode_{a}\else\textsubscript{\ensuremath{a}}\fi}
\newunicodechar{ᵢ}{\ifmmode_{i}\else\textsubscript{\ensuremath{i}}\fi}
\newunicodechar{ᵣ}{\ifmmode_{r}\else\textsubscript{\ensuremath{r}}\fi}
\newunicodechar{ₖ}{\ifmmode_{k}\else\textsubscript{\ensuremath{k}}\fi}
\newunicodechar{ᵦ}{\ifmmode_{\beta}\else\textsubscript{\ensuremath{\beta}}\fi}
\newunicodechar{ₓ}{\ifmmode_{x}\else\textsubscript{\ensuremath{x}}\fi}
\newfontfamily\popdisplay{ArchivoBlack-Regular.ttf}[Path=../../../fonts/]
\newfontfamily\poplogo{SpaceGrotesk-Bold.ttf}[Path=../../../fonts/]
\newfontfamily\popsemi{PlusJakartaSans-SemiBold.ttf}[Path=../../../fonts/]
\newfontfamily\popxbold{PlusJakartaSans-ExtraBold.ttf}[Path=../../../fonts/]
\newfontfamily\popxboldls{PlusJakartaSans-ExtraBold.ttf}[Path=../../../fonts/,LetterSpace=3.0]

\setlist[enumerate]{leftmargin=1.7em,itemsep=5pt,topsep=4pt}
\setlist[itemize]{leftmargin=1.7em,itemsep=4pt,topsep=4pt,label={\color{poporange}\textbullet}}
\setlength{\parindent}{0pt}
\setlength{\parskip}{5pt}
\renewcommand{\arraystretch}{1.25}

% pgfplots : style Pop par défaut
\pgfplotsset{
  every axis/.append style={axis line style={popink,line width=1pt},tick style={popink},
    grid style={popline,line width=0.5pt},label style={font=\small},tick label style={font=\footnotesize}},
  popcurve/.style={poporange,line width=1.8pt,no markers,smooth},
}
% Même style disponible dans un \draw TikZ simple (hors pgfplots)
\tikzset{popcurve/.style={poporange,line width=1.8pt}}
\tikzset{popgrid/.style={popline,very thin}}
\colorlet{popcurve}{poporange} % le nom du style est parfois utilisé comme couleur

% ============================================================
% Pill / squiggle
% ============================================================
\newcommand{\poppill}[3]{%
  \tikz[baseline=-0.55ex]\node[draw=popink,line width=1.2pt,rounded corners=2.6pt,%
    fill=#2,inner xsep=6.5pt,inner ysep=3pt]{%
    \popxboldls\fontsize{7}{7}\selectfont\color{#3}\MakeUppercase{#1}};%
}
\newcommand{\popsquiggle}[1]{%
  \tikz[baseline=-0.4ex]{
    \draw[#1,line width=2.6pt,line cap=round]
      plot [smooth] coordinates {(0,0.09) (0.34,0.01) (0.68,0.21) (1.02,0.01) (1.36,0.10)};
  }%
}
% Mot-clé surligné (marqueur orange sous le texte)
\newcommand{\cle}[1]{{\popxbold\color{popdark}#1}}

% ============================================================
% En-tête / pied
% ============================================================
\pagestyle{fancy}
\fancyhf{}
\renewcommand{\headrulewidth}{0pt}
\renewcommand{\footrulewidth}{0pt}
\lhead{\raisebox{-0.15ex}{\poplogo\fontsize{13}{13}\selectfont\color{popink}OnBuch\color{poporange}+}}
\rhead{\poppill{\DOCMATIERE\ \textbullet\ \DOCNIVEAU\ \textbullet\ Leçon \DOCLECON}{white}{popink}}
\lfoot{\popxbold\fontsize{7.6}{7.6}\selectfont\color{popmuted}\MakeUppercase{Cours signé OnBuch+}\ \textbullet\ {\popsemi\DOCTITRE}}
\rfoot{\tikz[baseline=-0.6ex]\node[rounded corners=8.5pt,fill=popink,minimum height=17pt,minimum width=17pt,inner xsep=5pt,inner sep=0]{\popxbold\fontsize{8}{8}\selectfont\color{popcream}\thepage\,/\,\pageref*{LastPage}};}

% ============================================================
% Titres
% ============================================================
\newcommand{\chaptitre}[2]{%
  \begin{tcolorbox}[enhanced,colback=poporange,colframe=popink,boxrule=2.6pt,arc=11pt,
      drop shadow=popink,left=15pt,right=15pt,top=12pt,bottom=11pt,before skip=3pt,after skip=18pt]
    \poppill{#2}{popink}{popcream}\par\vspace{9pt}
    {\raggedright\hyphenpenalty=10000\exhyphenpenalty=10000\popdisplay\fontsize{22}{24}\selectfont\color{popcream}\MakeUppercase{#1}\par}
  \end{tcolorbox}
}
% Section numérotée : pastille orange avec le numéro + titre Archivo
\newcounter{popsec}
\newcommand{\coursec}[1]{%
  \stepcounter{popsec}\setcounter{popsub}{0}%
  \par\vspace{16pt}\noindent
  \begin{minipage}{\linewidth}
  \tikz[baseline=-0.7ex]\node[draw=popink,line width=1.6pt,fill=poporange,rounded corners=4pt,minimum size=22pt,inner sep=0]{\popdisplay\fontsize{12}{12}\selectfont\color{popcream}\thepopsec};%
  \hspace{8pt}{\popdisplay\fontsize{15}{17}\selectfont\color{popink}\MakeUppercase{#1}}\par
  \vspace{4pt}\noindent{\color{poporange}\rule{36pt}{3.6pt}}
  \end{minipage}\par\nopagebreak\vspace{8pt}
}
\newcounter{popsub}
\newcommand{\courssub}[1]{%
  \stepcounter{popsub}%
  \par\vspace{9pt}\noindent{\popxbold\fontsize{11.5}{13}\selectfont\color{popink}{\color{poporange}\thepopsec.\thepopsub}\hspace{6pt}#1}\par\nopagebreak\vspace{4pt}
}

% ============================================================
% Boîtes pédagogiques — même recette : fond blanc, bordure épaisse colorée,
% ombre dure encre, badge plein. Titre optionnel [#1].
% ============================================================
\tcbset{popbase/.style={breakable,enhanced,colback=white,boxrule=2.3pt,arc=9pt,
  shadow={3pt}{-3pt}{0pt}{popink},left=14pt,right=14pt,top=11pt,bottom=11pt,before skip=11pt,after skip=15pt,
  fontupper=\popsemi\fontsize{10.6}{14.5}\selectfont\color{popink},
  fontlower=\popsemi\fontsize{10.6}{14.5}\selectfont\color{popink}}}
% Coche dessinée (le glyphe ✓ n'existe pas dans les polices du document)
\newcommand{\popcheck}[1]{\tikz[baseline=-0.5ex]\draw[#1,line width=1.6pt,line cap=round,line join=round](-0.13,0.0)--(-0.04,-0.09)--(0.13,0.1);}
\providecommand{\coche}{\popcheck{popgreen}}
\providecommand{\resultat}[1]{\par\smallskip\noindent\fcolorbox{popgreen}{popgreenL}{\parbox{\dimexpr\linewidth-2\fboxsep-2\fboxrule\relax}{\textbf{Résultat.} #1}}\par\smallskip}
\newcommand{\popboxhead}[3]{\poppill{#1}{#2}{white}\ifx\relax#3\relax\else\hspace{6pt}{\popxbold\fontsize{10.6}{12}\selectfont\color{popink}#3}\fi\par\vspace{7pt}}

\newtcolorbox{aretenir}[1][]{popbase,colframe=popgreen,before upper={\popboxhead{À retenir}{popgreen}{#1}}}
\newtcolorbox{exemplebox}[1][]{popbase,colframe=poporange,before upper={\popboxhead{Exemple}{poporange}{#1}}}
\newtcolorbox{definition}[1][]{popbase,colframe=popblue,before upper={\popboxhead{Définition}{popblue}{#1}}}
\newtcolorbox{propriete}[1][]{popbase,colframe=poppurple,before upper={\popboxhead{Propriété}{poppurple}{#1}}}
\newtcolorbox{methode}[1][]{popbase,colframe=popgold,before upper={\popboxhead{Méthode}{popgold}{#1}}}
\newtcolorbox{attention}[1][]{popbase,colframe=poppink,before upper={\popboxhead{Attention}{poppink}{#1}}}
\newtcolorbox{experience}[1][]{popbase,colframe=popblue,colback=popblueL,before upper={\popboxhead{Expérience}{popblue}{#1}}}
% « Le savais-tu ? » : carte violette pleine (contexte, culture, Cameroun)
\newtcolorbox{savaistu}[1][]{popbase,colback=poppurple,colframe=popink,
  fontupper=\popsemi\fontsize{10.4}{14.2}\selectfont\color{white},
  before upper={\poppill{Le savais-tu\,?}{white}{poppurple}\ifx\relax#1\relax\else\hspace{6pt}{\popxbold\color{white}#1}\fi\par\vspace{7pt}}}
% Exercice résolu : énoncé en haut, \tcblower puis la solution sur fond crème
\newcounter{popexo}\newcounter{popexr}
\newtcolorbox{exoresolu}[1][]{popbase,colframe=poporange,colbacklower=popcream,
  segmentation style={popink,line width=1.2pt,dashed},
  before upper={\stepcounter{popexr}\popboxhead{Exercice résolu \thepopexr}{poporange}{#1}},
  before lower={\poppill{Solution}{popink}{popcream}\par\vspace{6pt}}}
% Exercice (non résolu en place) — la correction est dans la partie Corrigés
\newtcolorbox{exercice}[1][]{popbase,colframe=popink,
  before upper={\stepcounter{popexo}\popboxhead{Exercice \thepopexo}{popink}{#1}}}
% Corrigé
\newtcolorbox{corrige}[1][]{popbase,colframe=popgreen,colback=popgreenL,
  before upper={\popboxhead{Corrigé}{popgreen}{#1}}}
% Objectifs / prérequis / bilan
\newtcolorbox{objectifs}{popbase,colframe=popink,colback=poporangeL,
  before upper={\popboxhead{Objectifs de la leçon}{popink}{}}}
\newtcolorbox{prerequis}{popbase,colframe=popink,colback=popblueL,
  before upper={\popboxhead{Prérequis}{popblue}{}}}
\newtcolorbox{bilan}{popbase,colframe=popink,colback=white,boxrule=2.8pt,
  before upper={\popboxhead{Fiche bilan}{poporange}{L'essentiel à connaître pour l'examen}}}
% Figure encadrée avec légende
\newtcolorbox{popfigure}[1][]{enhanced,breakable=false,colback=white,colframe=popline,boxrule=1.4pt,arc=8pt,
  left=10pt,right=10pt,top=10pt,bottom=8pt,before skip=10pt,after skip=12pt,halign=center,
  lower separated=false,halign lower=center,
  fontlower=\popsemi\fontsize{9}{11.5}\selectfont\color{popmuted},
  #1}
\newcommand{\legende}[1]{\par\vspace{6pt}{\popsemi\fontsize{9}{11.5}\selectfont\color{popmuted}#1}}

% ============================================================
% Signature OnBuch+
% ============================================================
\newcommand{\poplogoinline}[1]{{\poplogo\fontsize{#1}{#1}\selectfont\color{popink}OnBuch\color{poporange}+}}
\newcommand{\signatureonbuch}{%
  \begin{tcolorbox}[enhanced,colback=popink,colframe=popink,boxrule=2.2pt,arc=10pt,
      drop shadow=poporange,left=14pt,right=14pt,top=10pt,bottom=10pt,before skip=0pt,after skip=0pt]
    \tikz[baseline=-0.7ex]\node[circle,fill=popgreen,draw=popcream,line width=1.3pt,minimum size=22pt,inner sep=0]{\popcheck{white}};%
    \hspace{9pt}\begin{minipage}[c]{0.62\linewidth}
      {\popxbold\fontsize{12}{13}\selectfont\color{popcream}Signé\ \textbullet\ {\poplogo OnBuch\color{poporange}+}}\par\vspace{2pt}
      {\raggedright\popsemi\fontsize{8}{10.5}\selectfont\color{popline}\MakeUppercase{Équipe pédagogique OnBuch+}\\\MakeUppercase{Conforme au programme officiel MINESEC}\par}
    \end{minipage}\hfill
    {\poplogo\fontsize{22}{22}\selectfont\color{popcream}OnBuch\color{poporange}+}
  \end{tcolorbox}%
}

% ============================================================
% Page de garde
% #1 titre · #2 matière · #3 niveau · #4 module · #5 n° leçon · #6 durée ·
% #7 description / objectifs (texte ou itemize)
% ============================================================
\newcommand{\pagegarde}[7]{%
  \thispagestyle{empty}
  \newgeometry{a4paper,margin=1.7cm,top=1.6cm,bottom=1.4cm}%
  \begin{tikzpicture}[remember picture,overlay]
    \fill[poporange!18] ([xshift=-3.2cm,yshift=-3.6cm]current page.north east) circle (4.6cm);
    \fill[poppurple!14] ([xshift=2.4cm,yshift=7.6cm]current page.south west) circle (3.8cm);
  \end{tikzpicture}%
  {\poplogo\fontsize{30}{30}\selectfont\color{popink}OnBuch\color{poporange}+}\hfill
  \raisebox{0.6ex}{\poppill{Cours signé \textbullet\ Premium}{popink}{popcream}}\par
  \vspace{1pt}\popsquiggle{poppurple}\par
  \vspace{8pt}
  \begin{tcolorbox}[enhanced,colback=poporange,colframe=popink,boxrule=2.8pt,arc=13pt,
      drop shadow=popink,left=18pt,right=18pt,top=12pt,bottom=13pt,after skip=10pt]
    \poppill{#2 \textbullet\ #3}{popink}{popcream}\hspace{5pt}\poppill{#4}{white}{popink}\par\vspace{8pt}
    {\popdisplay\fontsize{14}{14}\selectfont\color{popink}LEÇON #5}\par\vspace{4pt}
    {\raggedright\hyphenpenalty=10000\exhyphenpenalty=10000\popdisplay\fontsize{25}{27}\selectfont\color{popcream}\MakeUppercase{#1}\par}
    \vspace{8pt}
    {\popxbold\fontsize{9.5}{9.5}\selectfont\color{popink}\MakeUppercase{Durée conseillée : #6}}
  \end{tcolorbox}
  \begin{tcolorbox}[enhanced,colback=white,colframe=popink,boxrule=2.2pt,arc=10pt,
      drop shadow=popink,left=16pt,right=16pt,top=10pt,bottom=10pt,after skip=8pt]
    \poppill{Dans cette leçon}{poppurple}{white}\par\vspace{5pt}
    \popsemi\fontsize{9.3}{12.6}\selectfont\color{popink}#7
  \end{tcolorbox}
  \noindent\begin{minipage}[t]{0.487\linewidth}
    \begin{tcolorbox}[enhanced,colback=popgreen,colframe=popink,boxrule=2.2pt,arc=10pt,
        drop shadow=popink,left=11pt,right=11pt,top=10pt,bottom=10pt,height=3.1cm,valign=center]
      \tcbox[colback=white,colframe=popink,boxrule=1.3pt,arc=5pt,boxsep=0pt,nobeforeafter,
        left=3pt,right=3pt,top=3pt,bottom=3pt,valign=center]{\qrcode[height=1.9cm]{https://play.google.com/store/apps/details?id=com.luvvix.onbuch&hl=fr}}%
      \hspace{8pt}\begin{minipage}[c]{0.5\linewidth}
        {\popdisplay\fontsize{11}{12.5}\selectfont\color{white}\MakeUppercase{Télécharge OnBuch}}\par\vspace{4pt}
        {\raggedright\popsemi\fontsize{8.4}{11}\selectfont\color{white}Gratuit \textbullet\ Google Play\\Tuteur IA, annales, concours, résultats\par}
      \end{minipage}
    \end{tcolorbox}
  \end{minipage}\hfill
  \begin{minipage}[t]{0.487\linewidth}
    \begin{tcolorbox}[enhanced,colback=poppurple,colframe=popink,boxrule=2.2pt,arc=10pt,
        drop shadow=popink,left=11pt,right=11pt,top=10pt,bottom=10pt,height=3.1cm,valign=center]
      \tikz[baseline=-0.7ex]\node[circle,fill=white,draw=popink,line width=1.3pt,minimum size=1.6cm,inner sep=0]{\popdisplay\fontsize{24}{24}\selectfont\color{poppurple}+};%
      \hspace{8pt}\begin{minipage}[c]{0.6\linewidth}
        {\popdisplay\fontsize{11}{12.5}\selectfont\color{white}\MakeUppercase{Passe à OnBuch+}}\par\vspace{4pt}
        {\raggedright\popsemi\fontsize{8.4}{11}\selectfont\color{white}Tous les cours signés, Tuteur IA illimité, fiches PDF — onglet OnBuch+ de l'app\par}
      \end{minipage}
    \end{tcolorbox}
  \end{minipage}
  \vspace{8pt}
  \popcontenu
  \vfill
  \signatureonbuch
  \restoregeometry
  \newpage
}

% Rangée « Ce cours contient » (page de garde)
\newcommand{\poptile}[3]{% #1 couleur #2 gros texte #3 légende
  \begin{tcolorbox}[enhanced,colback=#1,colframe=popink,boxrule=2pt,arc=9pt,shadow={2.5pt}{-2.5pt}{0pt}{popink},
      left=6pt,right=6pt,top=9pt,bottom=9pt,halign=center,valign=center,height=1.75cm,width=\linewidth,nobeforeafter]
    {\popdisplay\fontsize{12.5}{13}\selectfont\color{white}\MakeUppercase{#2}}\par\vspace{3pt}
    {\centering\popsemi\fontsize{7.8}{9.5}\selectfont\color{white}#3\par}
  \end{tcolorbox}}
\newcommand{\popcontenu}{%
  \par\noindent{\popxboldls\fontsize{7.5}{7.5}\selectfont\color{popmuted}CE COURS SIGNÉ CONTIENT}\par\vspace{7pt}
  \noindent\begin{minipage}[t]{0.235\linewidth}\poptile{popblue}{Cours}{complet, illustré, pas à pas}\end{minipage}\hfill
  \begin{minipage}[t]{0.235\linewidth}\poptile{poporange}{Méthodes}{et exercices résolus}\end{minipage}\hfill
  \begin{minipage}[t]{0.235\linewidth}\poptile{poppink}{Exercices}{type examen + corrigés}\end{minipage}\hfill
  \begin{minipage}[t]{0.235\linewidth}\poptile{popgreen}{Fiche bilan}{l'essentiel à retenir}\end{minipage}\par}

% Sommaire
\newtcolorbox{sommaire}{enhanced,colback=white,colframe=popink,boxrule=2.2pt,arc=10pt,
  drop shadow=popink,left=18pt,right=18pt,top=13pt,bottom=13pt,before skip=6pt,after skip=20pt,
  before upper={\poppill{Sommaire}{poppurple}{white}\par\vspace{9pt}\popsemi\fontsize{10.6}{14.5}\selectfont\color{popink}}}

% Signature de fin de cours — poussée tout en bas de la dernière page
\newcommand{\findecours}{%
  \par\vfill\nopagebreak
  \begin{center}\popsquiggle{poporange}\end{center}\vspace{4pt}
  \signatureonbuch
}
\AtBeginDocument{\def\llbracket{\mathopen{[\mkern-3.5mu[}}\def\rrbracket{\mathclose{]\mkern-3.5mu]}}} % intervalles d'entiers (glyphe ⟦ absent de la police)
% Glyphes absents de Fira Math : redéfinis à partir de glyphes disponibles
\AtBeginDocument{%
  \def\setminus{\mathbin{\backslash}}\let\smallsetminus\setminus
  \def\vdots{\mathord{\vbox{\baselineskip3.2pt\lineskiplimit0pt\kern4pt\hbox{.}\hbox{.}\hbox{.}}}}%
  \def\ddots{\mathinner{\mkern1mu\raise6.5pt\hbox{.}\mkern2mu\raise3.6pt\hbox{.}\mkern2mu\raise0.7pt\hbox{.}\mkern1mu}}%
  \def\triangle{\mathord{\scalebox{0.9}{$\symup{\Delta}$}}}%
  \def\nabla{\mathord{\raisebox{\depth}{\scalebox{1}[-1]{$\symup{\Delta}$}}}}%
  \def\checkmark{\popcheck{popdark}}%
  \def\star{\mathbin{\ast}}\def\bigstar{\mathord{\ast}}%
  \def\diamond{\mathbin{\scalebox{0.6}{$\Diamond$}}}%
  \def\bigcup{\mathop{\mathchoice{\raisebox{-0.3em}{\scalebox{1.7}{$\cup$}}}{\scalebox{1.25}{$\cup$}}{\cup}{\cup}}\displaylimits}%
  \def\bigcap{\mathop{\mathchoice{\raisebox{-0.3em}{\scalebox{1.7}{$\cap$}}}{\scalebox{1.25}{$\cap$}}{\cap}{\cap}}\displaylimits}%
  \def\lesseqqgtr{\mathrel{\lessgtr}}\def\bigoplus{\mathop{\oplus}\displaylimits}%
}
\providecommand{\ssi}{si et seulement si} % abréviation fréquente
\AtBeginDocument{\providecommand{\wideparen}{\overparen}} % arc de cercle

%% FIGFIX-BEGIN v5 — correctifs globaux des figures (docs/onbuch-generation/tools/figfix)
\usepackage{adjustbox}\usepackage{etoolbox}\tcbuselibrary{raster}
% toute figure TikZ/pgfplots plus large que la ligne est réduite à la largeur disponible
\BeforeBeginEnvironment{tikzpicture}{\begin{adjustbox}{max width=\linewidth}}
\AfterEndEnvironment{tikzpicture}{\end{adjustbox}}
% légendes pgfplots lisibles par-dessus les courbes
\pgfplotsset{every axis/.append style={legend style={fill=white,fill opacity=0.92,text opacity=1,draw=black!15,inner sep=3pt}}}
% étiquettes d'arêtes (syntaxe edge["..."]) avec fond blanc
\tikzset{every edge quotes/.append style={fill=white,inner sep=1.2pt}}
%% FIGFIX-END
\begin{document}

\pagegarde{Travail et puissance d'une force}{Physique}{1ère TI}{Module 2 : Mouvements et interactions mécaniques}{N03}{8 h}{Cette leçon introduit le travail d'une force, grandeur qui mesure l'effet d'une force le long d'un déplacement, puis la puissance qui mesure la rapidité à fournir ce travail. Elle traite le cas des forces constantes, du poids (force conservative), de la tension d'un ressort et des forces de moment constant appliquées aux solides en rotation.
\vspace{4pt}
\begin{multicols}{2}\small
\begin{itemize}
\item À la fin de cette leçon, je sais définir le travail d'une force constante et donner son unité
\item À la fin de cette leçon, je sais exprimer W(F) = F·AB·cos α et exploiter cette relation dans des situations concrètes
\item À la fin de cette leçon, je sais distinguer travail moteur, travail résistant et travail nul
\item À la fin de cette leçon, je sais définir une force conservative et utiliser le fait que le poids est conservative
\item À la fin de cette leçon, je sais exprimer le travail de la tension d'un ressort à réponse linéaire et l'exploiter
\item À la fin de cette leçon, je sais exploiter la relation W = M·Δθ entre travail et moment d'une force
\end{itemize}\end{multicols}}

\begin{sommaire}
\begin{multicols}{2}
\begin{enumerate}
\item Travail d'une force constante
\item Travail du poids et forces conservatives
\item Travail de la tension d'un ressort
\item Travail et moment d'une force (rotation)
\item Puissance d'une force
\item Puissance d'une force de moment constant
\item Activité d'intégration
\item Méthodes et exercices résolus
\item Exercices
\item Corrigés des exercices
\item Fiche bilan
\end{enumerate}
\end{multicols}
\end{sommaire}

\begin{objectifs}
\begin{itemize}
\item À la fin de cette leçon, je sais définir le travail d'une force constante et donner son unité
\item À la fin de cette leçon, je sais exprimer W(F) = F·AB·cos α et exploiter cette relation dans des situations concrètes
\item À la fin de cette leçon, je sais distinguer travail moteur, travail résistant et travail nul
\item À la fin de cette leçon, je sais définir une force conservative et utiliser le fait que le poids est conservative
\item À la fin de cette leçon, je sais exprimer le travail de la tension d'un ressort à réponse linéaire et l'exploiter
\item À la fin de cette leçon, je sais exploiter la relation W = M·Δθ entre travail et moment d'une force
\item À la fin de cette leçon, je sais définir la puissance moyenne P = W/Δt et l'exprimer en watts
\item À la fin de cette leçon, je sais exprimer la puissance d'une force de moment constant P = M·ω
\end{itemize}
\end{objectifs}

\begin{prerequis}
\begin{itemize}
\item Notion de force, représentation vectorielle et projection (Seconde)
\item Produit scalaire de deux vecteurs (mathématiques, 1ère C)
\item Moment d'une force par rapport à un axe et équilibre de rotation (leçon précédente du module)
\item Vitesse angulaire ω et mouvement de rotation uniforme
\item Unités du système international et conversions (cm→m, min→s)
\end{itemize}
\end{prerequis}

\newpage

\coursec{Travail d'une force constante}

À la carrière de sable de Nkolbisson, deux porteurs soulèvent des sacs de ciment : l'un monte les sacs par un escalier raide, l'autre utilise une rampe inclinée plus longue. Lequel des deux « travaille » le plus au sens de la physique ? Et le moteur de la grue du chantier voisin, qui hisse la même charge en quelques secondes, développe-t-il plus de travail ou plus de puissance ? Pour répondre, il faut définir rigoureusement deux grandeurs que le langage courant confond : le \cle{travail} et la \cle{puissance} d'une force. Dans le langage de tous les jours, « travailler » signifie fournir un effort ; en physique, le mot a un sens précis, mesurable, et parfois surprenant : un porteur immobile qui maintient un sac de \SI{50}{kg} sur la tête se fatigue, mais au sens de la physique, sa force ne travaille pas ! Cette première section construit la notion de travail d'une force constante ; la puissance viendra ensuite.

\courssub{Quand dit-on qu'une force « travaille » ?}

\begin{experience}[Trois situations pour deviner la définition]
Observons trois situations de la vie quotidienne :
\begin{enumerate}
\item Un enfant pousse un chariot de courses : le chariot avance dans le sens de la force. On sent que la force « sert à quelque chose ».
\item Un voyageur tire sa valise à roulettes à l'aéroport de Douala : la force de traction est inclinée vers le haut, mais seule sa partie « horizontale » fait avancer la valise ; la partie verticale soulage simplement les roulettes.
\item Un athlète maintient une haltère immobile au-dessus de sa tête : il transpire, il tremble... mais l'haltère ne bouge pas.
\end{enumerate}
\textbf{Observations.} Dans les cas 1 et 2, le point d'application de la force se \emph{déplace} ; dans le cas 3, il n'y a aucun déplacement. Dans le cas 2, seule la composante de la force \emph{dans la direction du mouvement} semble utile.

\textbf{Interprétation.} Une force travaille si deux conditions sont réunies : son point d'application se déplace, et la force possède une composante non nulle dans la direction du déplacement. C'est exactement ce que traduit le produit $F \times AB \times \cos\alpha$.
\end{experience}

\courssub{Définition du travail d'une force constante}

Considérons un chariot tiré sur un sol horizontal par une force $\vec{F}$ constante (en direction, en sens et en norme), inclinée d'un angle $\alpha$ au-dessus de l'horizontale. Le point d'application de la force se déplace en ligne droite d'un point $A$ vers un point $B$.

\begin{popfigure}
\begin{center}
\begin{tikzpicture}[scale=1.05, >=Stealth]
% sol
\fill[popcream] (-0.5,-0.35) rectangle (8.5,0);
\draw[popline] (-0.5,0) -- (8.5,0);
\foreach \x in {-0.3,0.3,...,8.3} \draw[popline] (\x,0) -- (\x-0.25,-0.3);
% chariot
\fill[popblueL] (1.2,0.25) rectangle (3.0,1.15);
\draw[popblue, thick] (1.2,0.25) rectangle (3.0,1.15);
\fill[popdark] (1.6,0.12) circle (0.13);
\fill[popdark] (2.6,0.12) circle (0.13);
\node[popdark] at (2.1,0.72) {\small chariot};
% points A et B
\fill[popink] (1.2,0) circle (0.05) node[below left=1pt] {$A$};
\fill[popink] (7.2,0) circle (0.05) node[below right=1pt] {$B$};
\draw[->, popgreen, very thick] (1.2,-0.02) -- (7.2,-0.02) node[midway, below=4pt] {$\overrightarrow{AB}$};
% force F
\coordinate (P) at (3.0,0.9);
\draw[->, poporange, very thick] (P) -- ++(40:2.6) node[above right] {$\vec{F}$};
% composante utile
\draw[->, poppink, very thick, dashed] (P) -- ++(2.0,0) node[midway, below=2pt] {\small $F\cos\alpha$};
\draw[->, poppurple, thick, dashed] (P) ++(2.0,0) -- ++(0,1.68) node[midway, right=1pt] {\small $F\sin\alpha$};
% angle
\draw[popink] (P) ++(0.9,0) arc (0:40:0.9);
\node[popink] at ($(P)+(20:1.25)$) {$\alpha$};
\end{tikzpicture}
\end{center}
\legende{Chariot tiré par une force constante $\vec{F}$ inclinée d'un angle $\alpha$ sur l'horizontale. Seule la composante $F\cos\alpha$, colinéaire au déplacement $\overrightarrow{AB}$, contribue au travail.}
\end{popfigure}

\begin{definition}[Travail d'une force constante]
Le \cle{travail} d'une force \textbf{constante} $\vec{F}$ dont le point d'application se déplace en ligne droite d'un point $A$ à un point $B$ est le produit scalaire du vecteur force par le vecteur déplacement :
\[ W_{AB}(\vec{F}) = \vec{F}\cdot\overrightarrow{AB} = F \times AB \times \cos\alpha \]
où $F = \|\vec{F}\|$ est la norme de la force (en newtons), $AB$ la distance parcourue (en mètres) et $\alpha$ l'angle entre les vecteurs $\vec{F}$ et $\overrightarrow{AB}$.
\end{definition}

Le lien avec le \cle{produit scalaire} (vu en mathématiques) est précieux au Probatoire : il permet d'utiliser toutes ses propriétés. En particulier, si l'on décompose $\vec{F}$ en une composante $F\cos\alpha$ parallèle au déplacement et une composante $F\sin\alpha$ perpendiculaire, seule la composante parallèle travaille :
\[ W_{AB}(\vec{F}) = \underbrace{F\cos\alpha}_{\text{composante utile}} \times AB. \]
C'est exactement ce que l'expérience de la valise à roulettes nous avait suggéré.

\begin{attention}[Conditions d'application de la formule]
La formule $W_{AB}(\vec{F}) = F \cdot AB \cdot \cos\alpha$ n'est valable que si la force est \textbf{constante} (même norme, même direction, même sens pendant tout le déplacement) \textbf{et} si le déplacement est \textbf{rectiligne}. Si l'une de ces conditions n'est pas respectée (force variable, trajet courbe), il faut découper le trajet en petits morceaux — voir le paragraphe sur le trajet quelconque.
\end{attention}

\courssub{Unité de travail : le joule}

\begin{definition}[Le joule]
L'unité de travail dans le Système international est le \cle{joule} (symbole J), du nom du physicien anglais James Prescott Joule (1818--1889). Un joule est le travail d'une force de \SI{1}{N} dont le point d'application se déplace de \SI{1}{m} dans la direction de la force :
\[ \SI{1}{J} = \SI{1}{N} \times \SI{1}{m} = \SI{1}{kg.m^2.s^{-2}}. \]
\end{definition}

Vérifions l'homogénéité : $[W] = [F]\times[L] = \mathrm{kg.m.s^{-2}}\times\mathrm{m} = \mathrm{kg.m^2.s^{-2}}$. C'est bien la dimension d'une énergie — nous verrons au chapitre suivant que travail et énergie sont intimement liés.

Ordres de grandeur : soulever un sac de ciment de \SI{50}{kg} de \SI{1}{m} demande un travail d'environ \SI{500}{J} ; gravir un étage coûte à un élève de \SI{60}{kg} environ \SI{1800}{J}.

\begin{center}
\begin{tabularx}{\linewidth}{|l|X|X|}\hline
\rowcolor{popblueL} \textbf{Unité} & \textbf{Équivalence} & \textbf{Usage typique} \\ \hline
joule (J) & unité SI & physique, exercices \\ \hline
kilojoule (kJ) & $\SI{1}{kJ} = 10^{3}\ \mathrm{J}$ & travaux mécaniques importants \\ \hline
wattheure (Wh) & $\SI{1}{Wh} = \SI{3600}{J}$ & petites consommations électriques \\ \hline
kilowattheure (kWh) & $\SI{1}{kWh} = \num{3,6}\times 10^{6}\ \mathrm{J}$ & facture d'électricité ENEO \\ \hline
\end{tabularx}
\end{center}

\begin{savaistu}[Le kilowattheure d'ENEO]
Le \cle{kilowattheure} (kWh) que facture ENEO sur ta facture d'électricité est une unité de \textbf{travail} (c'est-à-dire d'énergie), et non de puissance ! Un kWh correspond au travail fourni par un appareil de puissance \SI{1000}{W} fonctionnant pendant une heure :
\[ \SI{1}{kWh} = \SI{1000}{W} \times \SI{3600}{s} = \num{3,6}\times 10^{6}\ \mathrm{J}. \]
Quand ta famille paie « 100 kWh » à ENEO, elle paie donc $\num{3,6}\times 10^{8}\ \mathrm{J}$ d'énergie électrique. Retiens bien : le watt mesure une puissance, le wattheure mesure un travail.
\end{savaistu}

\courssub{Le signe du travail : moteur, résistant ou nul}

Le travail est une grandeur \emph{algébrique} : il peut être positif, négatif ou nul, selon la valeur de l'angle $\alpha$, car c'est le signe de $\cos\alpha$ qui décide.

\begin{propriete}[Signe du travail — admise]
Le travail d'une force constante est une \cle{grandeur algébrique} dont le signe dépend de l'angle $\alpha$ entre la force et le déplacement :
\begin{itemize}
\item si $0 \leq \alpha < 90°$ : $\cos\alpha > 0$, donc $W > 0$ : le travail est \cle{moteur} — la force favorise le mouvement ;
\item si $\alpha = 90°$ : $\cos\alpha = 0$, donc $W = 0$ : le travail est \cle{nul} — la force ne favorise ni ne s'oppose au mouvement ;
\item si $90° < \alpha \leq 180°$ : $\cos\alpha < 0$, donc $W < 0$ : le travail est \cle{résistant} — la force s'oppose au mouvement.
\end{itemize}
Le travail est également nul si la force est nulle ($F = 0$) ou si le déplacement est nul ($A = B$).
\end{propriete}

\begin{popfigure}
\begin{center}
\begin{tikzpicture}[scale=0.95, >=Stealth]
% Cas 1 : moteur
\begin{scope}
\draw[popline] (-1.4,0) -- (1.4,0);
\draw[->, popgreen, very thick] (-0.9,0.12) -- (0.9,0.12) node[midway, below=2pt] {\small $\overrightarrow{AB}$};
\draw[->, poporange, very thick] (0,0.12) -- (35:1.3) node[above] {\small $\vec{F}$};
\draw[popink] (0.55,0.12) arc (0:35:0.55);
\node[popink] at (18:0.85) {\scriptsize $\alpha$};
\node[popdark, align=center] at (0,-0.9) {\small $\alpha < 90°$\\ \small \textbf{moteur} : $W>0$};
\end{scope}
% Cas 2 : nul
\begin{scope}[xshift=5.2cm]
\draw[popline] (-1.4,0) -- (1.4,0);
\draw[->, popgreen, very thick] (-0.9,0.12) -- (0.9,0.12) node[midway, below=2pt] {\small $\overrightarrow{AB}$};
\draw[->, poppurple, very thick] (0,0.12) -- (90:1.3) node[above] {\small $\vec{F}$};
\draw[popink] (0.45,0.12) arc (0:90:0.45);
\node[popink] at (45:0.72) {\scriptsize $90°$};
\node[popdark, align=center] at (0,-0.9) {\small $\alpha = 90°$\\ \small \textbf{nul} : $W=0$};
\end{scope}
% Cas 3 : résistant
\begin{scope}[xshift=10.4cm]
\draw[popline] (-1.4,0) -- (1.4,0);
\draw[->, popgreen, very thick] (-0.9,0.12) -- (0.9,0.12) node[midway, below=2pt] {\small $\overrightarrow{AB}$};
\draw[->, poppink, very thick] (0,0.12) -- (145:1.3) node[above] {\small $\vec{F}$};
\draw[popink] (0.55,0.12) arc (0:145:0.55);
\node[popink] at (72:0.8) {\scriptsize $\alpha$};
\node[popdark, align=center] at (0,-0.9) {\small $\alpha > 90°$\\ \small \textbf{résistant} : $W<0$};
\end{scope}
\end{tikzpicture}
\end{center}
\legende{Les trois cas possibles pour le signe du travail, selon l'angle $\alpha$ entre la force $\vec{F}$ et le déplacement $\overrightarrow{AB}$.}
\end{popfigure}

\begin{exemplebox}[Des forces qui ne travaillent pas]
Deux exemples fondamentaux de forces de travail nul, à connaître par cœur :
\begin{itemize}
\item La \textbf{réaction normale} $\vec{R}_N$ d'un support horizontal ou incliné est perpendiculaire au déplacement du solide : $\alpha = 90°$, donc $W(\vec{R}_N) = 0$.
\item La \textbf{tension d'un fil inextensible} qui retient un objet en mouvement circulaire (une pierre tournoyant au bout d'une ficelle, un satellite autour de la Terre) est à chaque instant dirigée vers le centre, donc perpendiculaire à la vitesse : son travail est nul. C'est pourquoi un satellite peut tourner « indéfiniment » sans moteur.
\end{itemize}
De même, le porteur immobile qui soutient un sac exerce bien une force, mais sans déplacement : $W = 0$. Sa fatigue vient de phénomènes musculaires internes, pas d'un travail mécanique au sens de la physique.
\end{exemplebox}

\courssub{Méthode de calcul et exemple chiffré}

\begin{methode}[Calculer le travail d'une force constante]
\begin{enumerate}
\item \textbf{Identifier} la force $\vec{F}$ étudiée et le déplacement $\overrightarrow{AB}$ de son point d'application ; vérifier que la force est constante et le trajet rectiligne.
\item \textbf{Mesurer l'angle} $\alpha$ entre le vecteur force et le vecteur déplacement (jamais avec la verticale par défaut !).
\item \textbf{Convertir} toutes les grandeurs en unités SI : force en newtons (N), distance en mètres (m).
\item \textbf{Appliquer} la formule $W_{AB}(\vec{F}) = F \times AB \times \cos\alpha$, donner le résultat en joules avec un nombre raisonnable de chiffres significatifs, puis \textbf{conclure sur le signe} : travail moteur, résistant ou nul.
\end{enumerate}
\end{methode}

\begin{exoresolu}[Le porteur de la gare routière]
\textbf{Énoncé.} À la gare routière de Mokolo, un porteur tire une charge de masse $m = \SI{20}{kg}$ posée sur un chariot, avec une force de norme $F = \SI{150}{N}$ faisant un angle $\alpha = 30°$ avec l'horizontale. Le chariot parcourt $AB = \SI{50}{m}$ en ligne droite sur un sol horizontal. Calculer le travail de la force de traction et préciser sa nature.
\tcblower
\textbf{Solution détaillée.}
\begin{enumerate}
\item Force étudiée : la traction $\vec{F}$, constante ; déplacement rectiligne $\overrightarrow{AB}$ horizontal.
\item L'angle entre $\vec{F}$ et $\overrightarrow{AB}$ est $\alpha = 30°$ (donné par rapport à l'horizontale, qui est la direction du déplacement).
\item Conversions : $F = \SI{150}{N}$, $AB = \SI{50}{m}$ : tout est déjà en unités SI.
\item Application numérique :
\begin{align*}
W_{AB}(\vec{F}) &= F \times AB \times \cos\alpha \\
&= 150 \times 50 \times \cos 30° \\
&= 7500 \times 0{,}866 \\
&\approx \SI{6495}{J} \approx \SI{6,5}{kJ}.
\end{align*}
\end{enumerate}
Comme $\alpha = 30° < 90°$, le travail est \textbf{moteur} : la traction favorise l'avancée du chariot. Remarquons que la composante verticale $F\sin 30° = \SI{75}{N}$ ne travaille pas : elle allège simplement l'appui du chariot sur le sol.
\end{exoresolu}

\begin{attention}[Les deux erreurs classiques au Probatoire]
\begin{itemize}
\item \textbf{Mauvais angle.} L'angle $\alpha$ de la formule est toujours l'angle entre la force et le \emph{déplacement}. Si l'énoncé donne l'angle avec la verticale, il faut d'abord calculer $\alpha = 90° - \beta$. Lis l'énoncé deux fois et fais un schéma !
\item \textbf{Conversions oubliées.} Une distance donnée en centimètres ou en kilomètres doit être convertie en mètres : $AB = \SI{250}{cm} = \SI{2,50}{m}$. Un résultat en joules calculé avec des centimètres est 100 fois trop grand — et faux.
\end{itemize}
\end{attention}

\courssub{Travail d'une force sur un trajet quelconque}

Que faire si le trajet n'est pas rectiligne, ou si la force change de direction en cours de route ? L'idée, que nous admettrons, consiste à \textbf{découper le trajet} en une multitude de très petits déplacements $\delta\vec{\ell}$, suffisamment petits pour être assimilés à des segments de droite et pour que la force y soit pratiquement constante. Sur chaque petit tronçon, le travail élémentaire est
\[ \delta W = \vec{F}\cdot\delta\vec{\ell} = F\,\delta\ell\cos\alpha, \]
et le travail total est la \textbf{somme} de tous ces travaux élémentaires le long du trajet de $A$ à $B$ :
\[ W_{AB}(\vec{F}) = \sum \vec{F}\cdot\delta\vec{\ell}. \]

\begin{aretenir}[L'essentiel de la section]
\begin{itemize}
\item Le travail d'une force constante sur un déplacement rectiligne est $W_{AB}(\vec{F}) = \vec{F}\cdot\overrightarrow{AB} = F \times AB \times \cos\alpha$.
\item Unité SI : le \cle{joule} ; $\SI{1}{J} = \SI{1}{N}\times\SI{1}{m}$ ; $\SI{1}{kWh} = \num{3,6}\times 10^{6}\ \mathrm{J}$.
\item Le signe dépend de $\alpha$ : travail \cle{moteur} ($\alpha < 90°$), \cle{nul} ($\alpha = 90°$), \cle{résistant} ($\alpha > 90°$).
\item Une force perpendiculaire au déplacement (réaction normale, tension d'un fil en mouvement circulaire) ne travaille pas ; une force sans déplacement ne travaille pas non plus.
\item Sur un trajet quelconque, on découpe en petits déplacements et on additionne les travaux élémentaires.
\end{itemize}
\end{aretenir}

\begin{exercice}[Pour s'entraîner]
Un élève pousse une cantine (grande marmite) sur le sol de la cour du lycée avec une force horizontale de norme $F = \SI{80}{N}$ sur une distance $d = \SI{12}{m}$.
\begin{enumerate}
\item Calculer le travail de cette force et préciser sa nature.
\item Les forces de frottement du sol s'opposent au mouvement avec une norme constante $f = \SI{30}{N}$. Calculer le travail des frottements et préciser sa nature.
\item Quel est le travail de la réaction normale du sol ? Justifier.
\end{enumerate}
\end{exercice}

\begin{corrige}[Exercice : la cantine poussée]
\begin{enumerate}
\item La force de poussée est horizontale, dans le sens du déplacement : $\alpha = 0°$.
\[ W(\vec{F}) = F \times d \times \cos 0° = 80 \times 12 \times 1 = \SI{960}{J}. \]
$\alpha < 90°$ : travail \textbf{moteur}.
\item Les frottements sont opposés au déplacement : $\alpha = 180°$, donc $\cos 180° = -1$.
\[ W(\vec{f}) = f \times d \times \cos 180° = 30 \times 12 \times (-1) = \SI{-360}{J}. \]
$\alpha > 90°$ : travail \textbf{résistant}.
\item La réaction normale $\vec{R}_N$ est verticale, perpendiculaire au déplacement horizontal : $\alpha = 90°$, donc $W(\vec{R}_N) = \SI{0}{J}$. Une force perpendiculaire au déplacement ne travaille pas.
\end{enumerate}
\end{corrige}

\coursec{Travail du poids et forces conservatives}

Dans la section précédente, nous avons défini le travail d'une force constante $\vec{F}$ dont le point d'application se déplace de $A$ vers $B$ :
\[ W_{AB}(\vec{F}) = \vec{F}\cdot\overrightarrow{AB} = F\times AB \times \cos\alpha, \]
où $\alpha$ est l'angle entre la force et le déplacement. Nous allons maintenant appliquer cette relation à la force la plus familière de notre quotidien : le \cle{poids}. Le résultat que nous obtiendrons est remarquable et constitue l'un des piliers de la mécanique : le travail du poids possède une propriété que ne partagent pas toutes les forces, celle d'être \emph{indépendant du chemin suivi}. Cette propriété définit ce qu'on appelle une \cle{force conservative}.

\courssub{Travail du poids entre deux points}

\textbf{Mise en situation.} Un porteur au marché Mokolo monte un sac d'ignames du rez-de-chaussée au premier étage. Il peut emprunter l'escalier droit, très raide, ou bien une longue rampe inclinée douce. Dans les deux cas, le sac part du même point $A$ et arrive au même point $B$. Le travail du poids du sac est-il le même ? C'est la question à laquelle nous allons répondre rigoureusement.

\textbf{Expression générale du travail du poids.}\par

Considérons un solide de masse $m$ dont le centre d'inertie $G$ se déplace d'un point $A$ d'altitude $z_A$ vers un point $B$ d'altitude $z_B$. On choisit un axe vertical $(Oz)$ orienté \emph{vers le haut}, muni du vecteur unitaire $\vec{k}$. Le poids s'écrit alors :
\[ \vec{P} = m\vec{g} = -mg\,\vec{k}, \]
car le vecteur champ de pesanteur $\vec{g}$ est vertical, dirigé vers le bas.

\begin{experience}[Démonstration guidée : travail du poids]
\textbf{Objectif :} établir l'expression du travail du poids entre deux points $A$ et $B$ quelconques.
\begin{enumerate}
\item \textbf{Choisir un repère.} On prend un repère orthonormé $(O\,;\vec{i},\vec{j},\vec{k})$ avec $\vec{k}$ vertical ascendant. Les points $A$ et $B$ ont pour coordonnées $A(x_A\,;y_A\,;z_A)$ et $B(x_B\,;y_B\,;z_B)$.
\item \textbf{Exprimer les vecteurs.} Le poids est $\vec{P} = -mg\,\vec{k}$ et le déplacement est
\[ \overrightarrow{AB} = (x_B - x_A)\,\vec{i} + (y_B - y_A)\,\vec{j} + (z_B - z_A)\,\vec{k}. \]
\item \textbf{Calculer le produit scalaire.} Comme $\vec{i}\cdot\vec{k} = 0$ et $\vec{j}\cdot\vec{k} = 0$ (vecteurs perpendiculaires), seule la composante verticale du déplacement contribue :
\[ W_{AB}(\vec{P}) = \vec{P}\cdot\overrightarrow{AB} = -mg\,\vec{k}\cdot\overrightarrow{AB} = -mg\,(z_B - z_A). \]
\item \textbf{Conclure.}
\[ W_{AB}(\vec{P}) = mg\,(z_A - z_B). \]
\end{enumerate}
\textbf{Interprétation physique :} le poids étant vertical, seule la \emph{projection verticale} du déplacement compte. Tout déplacement horizontal se fait perpendiculairement au poids : le travail du poids y est nul ($\cos 90\degree = 0$).
\end{experience}

\begin{propriete}[Travail du poids — démonstration exigible]
Le travail du poids d'un solide de masse $m$ dont le centre d'inertie passe d'un point $A$ d'altitude $z_A$ à un point $B$ d'altitude $z_B$ est :
\[ \boxed{W_{AB}(\vec{P}) = mg\,(z_A - z_B)} \]
avec $m$ en kilogrammes (kg), $g$ en newtons par kilogramme (N/kg), $z_A$ et $z_B$ en mètres (m), et $W$ en joules (J). L'axe des altitudes est orienté vers le haut. \emph{Cette démonstration est exigible au Probatoire.}
\end{propriete}

Vérifions l'homogénéité : $[mg(z_A-z_B)] = \mathrm{kg}\times\mathrm{m\,s^{-2}}\times\mathrm{m} = \mathrm{kg\,m^2\,s^{-2}} = \mathrm{J}$. La relation est bien dimensionnellement correcte.

\textbf{Analyse du signe.} Posons $h = z_A - z_B$, la \cle{dénivellation} entre le point de départ et le point d'arrivée.
\begin{itemize}
\item Si le corps \textbf{descend} ($z_A > z_B$, donc $h>0$) : $W_{AB}(\vec{P}) = mgh > 0$. Le travail du poids est \cle{moteur} : le poids favorise le mouvement. C'est ce qui fait descendre l'eau des barrages de Lom Pangar ou d'Edéa sans aucun moteur.
\item Si le corps \textbf{monte} ($z_A < z_B$, donc $h<0$) : $W_{AB}(\vec{P}) < 0$. Le travail du poids est \cle{résistant} : le poids s'oppose au mouvement. C'est pourquoi monter une côte à moto exige du moteur un effort supplémentaire.
\item Si le déplacement est \textbf{horizontal} ($z_A = z_B$) : $W_{AB}(\vec{P}) = 0$. Le poids ne travaille pas.
\end{itemize}

\courssub{Le travail du poids ne dépend pas du chemin suivi}

Regardons attentivement la formule $W_{AB}(\vec{P}) = mg(z_A - z_B)$ : elle ne fait intervenir \emph{que} les altitudes de départ et d'arrivée. Ni la longueur du trajet, ni sa forme (droite, brisée, courbe) n'apparaissent. C'est un constat fondamental.

\begin{popfigure}
\begin{tikzpicture}[scale=1.0]
% Axe vertical des altitudes
\draw[->,thick,popdark] (-0.6,0) -- (-0.6,4.6) node[above]{$z$};
\draw[popdark] (-0.75,0.6) -- (-0.45,0.6) node[left=2pt]{$z_B$};
\draw[popdark] (-0.75,3.8) -- (-0.45,3.8) node[left=2pt]{$z_A$};
% Points A et B
\coordinate (A) at (0.5,3.8);
\coordinate (B) at (6.5,0.6);
\fill[popblue] (A) circle (2.2pt) node[above left]{$A$};
\fill[poporange] (B) circle (2.2pt) node[below right]{$B$};
% Chemin 1 : vertical puis horizontal
\draw[very thick,popblue] (A) -- (0.5,0.6) -- (B);
\node[popblue, font=\small] at (0.2,2.2) {chemin 1};
% Chemin 2 : droite inclinée
\draw[very thick,popgreen] (A) -- (B);
\node[popgreen, font=\small, rotate=-28] at (3.4,2.55) {chemin 2};
% Chemin 3 : courbe
\draw[very thick,poppurple] (A) .. controls (2.2,3.9) and (3.2,1.0) .. (B);
\node[poppurple, font=\small] at (3.1,3.75) {chemin 3};
% Poids en un point du chemin 2
\coordinate (M) at (3.5,2.2);
\fill[poppurple] (M) circle (2pt) node[above right=1pt]{\small $G$};
\draw[->,very thick,popink] (M) -- ++(0,-1.3) node[below]{$\vec{P}$};
% Dénivellation
\draw[<->,thick,poporange] (7.1,3.8) -- (7.1,0.6) node[midway,right]{$h = z_A - z_B$};
\draw[dashed,popmuted] (A) -- (7.1,3.8);
\draw[dashed,popmuted] (B) -- (7.1,0.6);
\end{tikzpicture}
\legende{Un corps descend de $A$ vers $B$ par trois chemins différents : vertical puis horizontal (1), droit incliné (2), courbe (3). Dans les trois cas, le travail du poids est le même : $W_{AB}(\vec{P}) = mg(z_A - z_B) = mgh$.}
\end{popfigure}

Pour s'en convaincre sur le chemin 1 (vertical puis horizontal) : sur la portion verticale de longueur $h$, le poids et le déplacement sont colinéaires et de même sens, donc $W_1 = mgh$ ; sur la portion horizontale, le poids est perpendiculaire au déplacement, donc $W_2 = 0$. Au total : $W = mgh + 0 = mgh$, exactement comme pour les chemins 2 et 3. On peut découper n'importe quel chemin courbe en une infinité de petits segments que l'on projette sur la verticale : la somme des projections verticales vaut toujours $z_A - z_B$.

\begin{definition}[Force conservative]
Une force est dite \cle{conservative} si son travail entre deux points $A$ et $B$ est \textbf{indépendant du chemin suivi} pour aller de $A$ à $B$ : il ne dépend que des positions de $A$ et de $B$.
\end{definition}

\begin{propriete}[Le poids est une force conservative]
Puisque $W_{AB}(\vec{P}) = mg(z_A - z_B)$ ne dépend que des altitudes de $A$ et de $B$, le \cle{poids est une force conservative}.
\textbf{Conséquence immédiate :} sur tout \emph{trajet fermé} (le corps revient à son point de départ, $z_B = z_A$), le travail du poids est nul :
\[ W_{\text{trajet fermé}}(\vec{P}) = 0. \]
\end{propriete}

\begin{attention}[Erreur fréquente : $mg\times L$ au lieu de $mg\times h$]
L'erreur classique au Probatoire consiste à multiplier le poids par la \emph{longueur du trajet} $L$ parcouru : $W = mgL$. C'est \textbf{faux} en général ! Le travail du poids fait intervenir uniquement la \cle{dénivellation} $h = z_A - z_B$, c'est-à-dire la différence d'altitude, et non la distance parcourue. Sur un plan incliné d'angle $\alpha$ et de longueur $L$, on a $h = L\sin\alpha$, donc $W(\vec{P}) = mgL\sin\alpha$, et non $mgL$. Toujours se poser la question : « de combien le corps est-il monté ou descendu verticalement ? »
\end{attention}

\begin{methode}[Calculer le travail du poids en trois étapes]
\begin{enumerate}
\item \textbf{Repérer les altitudes} des points de départ $A$ et d'arrivée $B$ sur un axe vertical orienté vers le haut (choisir une origine quelconque, le résultat n'en dépend pas).
\item \textbf{Calculer la dénivellation} $h = z_A - z_B$ (attention au signe !).
\item \textbf{Appliquer} $W_{AB}(\vec{P}) = mg(z_A - z_B)$ et \textbf{conclure sur le signe} : travail moteur ($W>0$) si le corps descend, résistant ($W<0$) s'il monte, nul si le déplacement est horizontal ou si le trajet est fermé.
\end{enumerate}
\end{methode}

\courssub{Application : escalier raide ou rampe douce ?}

Revenons à notre porteur du marché. Le schéma ci-dessous montre deux façons de franchir la même dénivellation $h$ : un plan incliné très pentu (angle $\alpha_1$ grand, longueur $L_1$ courte) et un plan incliné doux (angle $\alpha_2$ petit, longueur $L_2$ grande).

\begin{popfigure}
\begin{tikzpicture}[scale=0.85]
% Sol
\draw[thick,popdark] (0,0) -- (9.5,0);
\fill[pattern=north east lines,pattern color=popmuted] (0,-0.3) rectangle (9.5,0);
% Points
\coordinate (A) at (0,0);
\coordinate (B) at (8.24, 3.0); % h=3.0, alpha2=20 -> x=h/tan(20)=8.24
\coordinate (C) at (1.73, 3.0); % h=3.0, alpha1=60 -> x=h/tan(60)=1.73
\fill[popblue] (A) circle (2.2pt) node[below left]{$A$};
\fill[poporange] (B) circle (2.2pt) node[above right]{$B$};
\fill[popblue] (C) circle (2.2pt) node[above left]{$C$};
% Rampe douce (A -> B)
\draw[very thick,popgreen] (A) -- (B);
\node[popgreen, font=\small, rotate=20] at (4.1,1.0) {rampe douce $\alpha_2 = 20\degree$};
% Pente raide (A -> C -> B)
\draw[very thick,popblue] (A) -- (C) -- (B);
\node[popblue, font=\small, rotate=60] at (0.8,1.5) {pente raide $\alpha_1 = 60\degree$};
% Vertical reference
\draw[dashed,popmuted] (B) -- (8.24,0);
\draw[<->,thick,poporange] (8.8,0) -- (8.8,3.0) node[midway,right]{$h$};
% Poids sur la rampe douce
\coordinate (M) at (4.12, 1.5);
\fill[poppurple] (M) circle (2pt);
\draw[->,very thick,popink] (M) -- ++(0,-1.0) node[below]{$\vec{P}$};
% Angles
\draw[popblue] (0.5,0) arc (0:60:0.5) node[midway,right=2pt]{\scriptsize $60\degree$};
\draw[popgreen] (1.5,0) arc (0:20:1.5) node[midway,right=2pt]{\scriptsize $20\degree$};
\end{tikzpicture}
\legende{Deux chemins pour monter de $A$ à $B$ ($h=3,0$ m) : une pente raide ($\alpha_1=60\degree$) suivie d'un horizontal, et une rampe douce ($\alpha_2=20\degree$). La dénivellation $h$ est la même, donc le travail du poids est le même : $W_{AB}(\vec{P}) = -mgh$ (résistant).}
\end{popfigure}

\begin{exemplebox}[Sac d'ignames hissé de 10 m]
Un porteur hisse un sac d'ignames de masse $m = \num{50}$ kg d'une hauteur $h = \num{10}$ m, soit par un escalier raide, soit par une longue rampe inclinée. On prend $g = \num{9,8}$ N/kg.
\textbf{Travail du poids :} le sac monte, donc $z_A - z_B = -h = -\num{10}$ m :
\[ W_{AB}(\vec{P}) = mg(z_A - z_B) = -mgh = -\num{50}\times \num{9,8}\times \num{10} = -\num{4900}\ \mathrm{J}. \]
Le travail du poids est \textbf{résistant} et vaut $-\num{4900}$ J, \emph{identique} que l'on emprunte l'escalier ou la rampe. Pour monter le sac à vitesse constante, le porteur doit fournir un travail moteur d'au moins $+\num{4900}$ J dans les deux cas. La rampe ne réduit pas le travail total : elle réduit seulement la \emph{force} à exercer à chaque instant, au prix d'un déplacement plus long.
\end{exemplebox}

\begin{exoresolu}[Comparaison escalier / rampe]
\textbf{Énoncé.} Une charge de masse $m = \num{80}$ kg doit être montée du sol (point $A$) à une terrasse (point $B$) située à $h = \num{6,0}$ m de hauteur. Chemin 1 : un escalier dont la pente fait $\alpha_1 = 60\degree$ avec l'horizontale. Chemin 2 : une rampe faisant $\alpha_2 = 20\degree$ avec l'horizontale. On donne $g = \num{9,8}$ N/kg.
\begin{enumerate}
\item Calculer la longueur $L$ de chaque chemin.
\item Calculer le travail du poids pour chaque chemin. Conclure.
\end{enumerate}
\tcblower
\textbf{Solution.}
\begin{enumerate}
\item La dénivellation est $h = L\sin\alpha$, donc $L = \dfrac{h}{\sin\alpha}$.
\begin{align*}
L_1 &= \frac{\num{6,0}}{\sin 60\degree} = \frac{\num{6,0}}{\num{0,866}} \approx \num{6,9}\ \mathrm{m},\\
L_2 &= \frac{\num{6,0}}{\sin 20\degree} = \frac{\num{6,0}}{\num{0,342}} \approx \num{17,5}\ \mathrm{m}.
\end{align*}
La rampe est environ 2,5 fois plus longue que l'escalier.
\item Le travail du poids ne dépend que de la dénivellation :
\[ W_{AB}(\vec{P}) = -mgh = -\num{80}\times \num{9,8}\times \num{6,0} = -\num{4704}\ \mathrm{J} \approx -\num{4,7e3}\ \mathrm{J}. \]
\textbf{Conclusion :} le travail du poids est le même ($\approx -\num{4,7}$ kJ, résistant) sur l'escalier et sur la rampe, bien que les longueurs soient très différentes. C'est la signature d'une force conservative. La rampe présente un avantage mécanique : la composante du poids à vaincre le long du déplacement vaut $mg\sin\alpha$, soit $\approx \num{679}$ N sur l'escalier contre $\approx \num{268}$ N sur la rampe — l'effort est moindre, mais il s'exerce plus longtemps.
\end{enumerate}
\end{exoresolu}

\courssub{Forces non conservatives : l'exemple des frottements}

Toutes les forces ne partagent pas cette belle propriété. Considérons les \cle{forces de frottement} $\vec{f}$ qui s'exercent sur un meuble que l'on traîne sur le sol : elles sont constamment dirigées \emph{en sens opposé} au déplacement. Leur travail est donc \emph{toujours résistant}, et il est proportionnel à la \emph{longueur du trajet} effectivement parcouru : $W(\vec{f}) = -f\times L$ pour une force de frottement de norme constante.

Traîner le meuble directement de $A$ à $B$ sur \num{5} m, ou faire un détour de \num{20} m, ne donne pas du tout le même travail : plus le chemin est long, plus le travail des frottements est grand en valeur absolue. De plus, sur un trajet fermé, le travail des frottements n'est \emph{pas nul} : il est strictement négatif (l'énergie « perdue » échauffe les surfaces en contact).

\begin{aretenir}[Conservatif ou non ?]
\begin{itemize}
\item \textbf{Force conservative} : travail indépendant du chemin suivi ; travail nul sur tout trajet fermé. Exemple : \textbf{le poids} (et, nous le verrons, la tension d'un ressort à réponse linéaire).
\item \textbf{Force non conservative} : travail dépendant du chemin suivi ; travail non nul sur un trajet fermé. Exemple : \textbf{les forces de frottement}, dont le travail est toujours résistant et d'autant plus grand que le trajet est long.
\end{itemize}
\end{aretenir}

\begin{savaistu}[Les routes de montagne en lacets]
Pour gravir les reliefs abrupts du Cameroun — le col de la Falèse sur l'axe Yaoundé--Bafoussam ou les côtes de la route de Dschang — les ingénieurs tracent des routes en lacets plutôt qu'une rampe directe. Un élève pressé pourrait croire que ces détours « économisent du travail » : c'est faux ! Le travail à fournir contre le poids pour amener un camion en haut vaut toujours $mgh$, quelle que soit la longueur de la route. L'intérêt des lacets est ailleurs : en augmentant la longueur du trajet, on \emph{diminue la pente}, donc la composante du poids $mg\sin\alpha$ que le moteur doit vaincre à chaque instant. Le moteur, de puissance limitée, peut alors gravir la côte — plus lentement, mais sans caler. Même principe que la rampe du marché : on échange de la force contre de la distance, le travail, lui, reste inchangé.
\end{savaistu}

\begin{exercice}[Pour s'entraîner]
Un enfant de masse $m = \num{30}$ kg descend un toboggan de hauteur $h = \num{2,5}$ m. On donne $g = \num{9,8}$ N/kg.
\begin{enumerate}
\item Calculer le travail du poids pendant la descente. Préciser sa nature.
\item Le même enfant remonte ensuite à pied au sommet par l'escalier du toboggan. Que vaut le travail du poids sur cette montée ? Que vaut le travail total du poids sur l'aller-retour ? Commenter.
\end{enumerate}
\end{exercice}

\begin{corrige}[Exercice : toboggan]
\begin{enumerate}
\item L'enfant descend : $W(\vec{P}) = mgh = \num{30}\times \num{9,8}\times \num{2,5} = \num{735}$ J. Travail \textbf{moteur}.
\item À la montée : $W(\vec{P}) = -mgh = -\num{735}$ J, travail résistant. Sur l'aller-retour (trajet fermé) : $W_{\text{total}} = \num{735} - \num{735} = 0$ J. On retrouve la propriété des forces conservatives : travail nul sur tout trajet fermé.
\end{enumerate}
\end{corrige}

En résumé, le poids est une force conservative : son travail $W_{AB}(\vec{P}) = mg(z_A - z_B)$ ne dépend que de la dénivellation. Cette propriété, que ne possèdent pas les frottements, ouvrira la porte, en classe de Terminale, à la notion d'\emph{énergie potentielle de pesanteur}. Mais auparavant, intéressons-nous à une autre force conservative fondamentale : la tension d'un ressort.

\coursec{Travail de la tension d'un ressort}

Dans la section précédente, nous avons établi que le poids est une force \cle{conservative} : son travail entre deux points ne dépend pas du chemin suivi, mais uniquement des altitudes de départ et d'arrivée. Nous avons pu le démontrer parce que le poids est une force \emph{constante} (en norme, direction et sens) au voisinage du sol. Mais toutes les forces ne sont pas constantes ! Pense à un ressort que tu étires progressivement : plus tu tires, plus il faut forcer. La tension du ressort \emph{varie} au cours du déplacement. Comment calculer alors son travail ? La formule $W = F \cdot d \cdot \cos\alpha$, valable uniquement pour une force constante, ne s'applique plus directement. C'est tout l'objet de cette section : définir la tension d'un ressort, puis calculer son travail par une méthode graphique élégante, et enfin montrer que cette force, comme le poids, est conservative.

\courssub{Rappel : la loi de Hooke, une force variable}

\begin{experience}[Étude expérimentale d'un ressort à réponse linéaire]
\textbf{Matériel :} un ressort à spires non jointives, un support vertical, une règle graduée, des masses marquées (50 g, 100 g, 150 g\ldots), un crochet.

\textbf{Protocole :} on suspend le ressort verticalement et on mesure sa longueur à vide $\ell_0$. On accroche ensuite successivement des masses de plus en plus grandes et on mesure à chaque fois la nouvelle longueur $\ell$ du ressort, donc son allongement $x = \ell - \ell_0$.

\textbf{Observations :} tant que l'on ne dépasse pas la limite d'élasticité du ressort, l'allongement $x$ est proportionnel à la masse suspendue, donc au poids, donc à la force exercée sur le ressort. La courbe $T = f(x)$ est une droite passant par l'origine.

\textbf{Interprétation :} à l'équilibre d'une masse $m$ accrochée au ressort, le bilan des forces donne $\vec{T} + \vec{P} = \vec{0}$, soit $T = mg$. La proportionnalité observée entre $T$ et $x$ traduit une loi fondamentale : la loi de Hooke.
\end{experience}

\begin{definition}[Tension d'un ressort à réponse linéaire — loi de Hooke]
Un ressort à spires non jointives, étiré ou comprimé dans son domaine d'élasticité, exerce une force de rappel appelée \cle{tension}, dirigée le long de son axe et de sens opposé à sa déformation. Sa norme est proportionnelle à la valeur absolue de l'\cle{allongement algébrique} $x = \ell - \ell_0$ :
\[
T = k\,|x|
\]
où $k$ est la \cle{constante de raideur} du ressort, exprimée en newtons par mètre (\SI{}{N.m^{-1}}), et $\ell_0$ la longueur à vide (position de repos).
\textbf{Forme vectorielle (complément) :} si l'on choisit un axe orienté dans le sens de l'étirement, la tension s'écrit $\vec{T} = -k\,x\,\vec{e}_x$.
\end{definition}

\begin{attention}[La tension n'est pas une force constante]
Contrairement au poids, la tension d'un ressort \emph{dépend de la position} : $T = k|x|$ croît avec l'allongement. Il est donc \textbf{interdit} d'appliquer directement la formule $W = F \cdot d \cdot \cos\alpha$ avec la valeur finale de la tension. C'est précisément le piège classique du Probatoire que nous dénoncerons plus bas. Il faut une méthode adaptée aux forces variables : la méthode graphique.
\end{attention}

\begin{methode}[Bien poser le problème : allongement, pas longueur]
Pour tout exercice sur les ressorts, adopte systématiquement les réflexes suivants :
\begin{enumerate}
\item Repérer la \textbf{position de repos} (longueur à vide $\ell_0$) : c'est l'origine des allongements.
\item Exprimer chaque allongement $x = \ell - \ell_0$ \textbf{en mètres} (convertir les cm !), jamais la longueur totale $\ell$ du ressort.
\item Vérifier les signes : $x > 0$ pour un étirement, $x < 0$ pour une compression ; le travail ne dépendra que des carrés $x^2$.
\end{enumerate}
Une erreur de 10 cm au lieu de \SI{0,10}{m} fausse le résultat d'un facteur $10^4$, car le travail dépend de $x^2$ !
\end{methode}

\courssub{Méthode graphique : le travail comme aire sous la courbe}

\textbf{Intuition.} Étirons le ressort très lentement, par petites étapes. Pendant un tout petit déplacement $\Delta x$ autour d'un allongement $x$, la tension est \emph{presque constante} : on peut alors lui appliquer la formule du travail élémentaire. Le travail de la tension sur ce petit intervalle vaut, en valeur absolue, $T \cdot \Delta x = k|x|\,\Delta x$ : c'est l'aire d'un petit rectangle de hauteur $k|x|$ et de largeur $\Delta x$ sous la courbe $T = k|x|$. En additionnant tous ces petits rectangles entre l'allongement initial $x_A$ et l'allongement final $x_B$, on obtient l'\textbf{aire totale sous la courbe} $T = f(x)$ entre $x_A$ et $x_B$.

\begin{propriete}[Travail de la tension : aire sous la courbe]
Le travail de la tension d'un ressort lorsque son extrémité libre passe de l'allongement $x_A$ à l'allongement $x_B$ est égal, au signe près, à l'\cle{aire sous la courbe} $T = k|x|$ comprise entre $x_A$ et $x_B$. Comme la tension s'oppose au déplacement lors d'un éloignement de la position de repos (force de rappel), son travail est \textbf{résistant} si $|x_B| > |x_A|$ :
\[
W_{AB}(\vec{T}) = \frac{1}{2}\,k\left(x_A^2 - x_B^2\right)
\]
\end{propriete}

\begin{experience}[Démonstration guidée par le calcul d'aire (exigible au Probatoire)]
Considérons un ressort horizontal dont l'extrémité libre passe de l'allongement $x_A$ à l'allongement $x_B$, avec $0 \le x_A < x_B$ (étirement). La généralisation aux autres cas (compression, passage par le repos) conduit à la même formule algébrique.

\textbf{Étape 1 — Tracé.} La courbe $T = f(x)$ est la droite $T = kx$ passant par l'origine (pour $x \ge 0$). Aux allongements $x_A$ et $x_B$ correspondent les tensions $T_A = kx_A$ et $T_B = kx_B$.

\textbf{Étape 2 — Aire.} L'aire sous la droite entre $x_A$ et $x_B$ est un \textbf{trapèze} de bases $T_A$ et $T_B$ et de hauteur $(x_B - x_A)$ :
\[
\mathcal{A} = \frac{(T_A + T_B)}{2}\,(x_B - x_A) = \frac{(kx_A + kx_B)}{2}\,(x_B - x_A)
\]

\textbf{Étape 3 — Factorisation.} En utilisant l'identité remarquable $(a+b)(a-b) = a^2 - b^2$ :
\[
\mathcal{A} = \frac{k}{2}\,(x_A + x_B)(x_B - x_A) = \frac{k}{2}\left(x_B^2 - x_A^2\right)
\]

\textbf{Étape 4 — Signe.} Lors d'un étirement, la tension $\vec{T}$ est dirigée vers la position de repos, c'est-à-dire en sens \emph{opposé} au déplacement : son travail est résistant, donc négatif. D'où :
\[
W_{AB}(\vec{T}) = -\mathcal{A} = \frac{1}{2}\,k\left(x_A^2 - x_B^2\right) \qquad \text{c.q.f.d.}
\]

\textbf{Vérification dimensionnelle :} $[k] = \SI{}{N.m^{-1}}$ et $[x^2] = \SI{}{m^2}$, donc $[kx^2] = \SI{}{N.m} = \SI{}{J}$. La formule est homogène.
\end{experience}

\begin{popfigure}
\begin{center}
\begin{tikzpicture}
\begin{axis}[
  width=0.85\linewidth, height=6.5cm,
  axis lines=left,
  xlabel={Allongement $x$ (m)}, ylabel={Tension $T$ (N)},
  xmin=0, xmax=0.16, ymin=0, ymax=35,
  xtick={0.05,0.10}, xticklabels={$x_A$,$x_B$},
  ytick={10,20}, yticklabels={$T_A$,$T_B$},
  every axis x label/.style={at={(ticklabel* cs:1)},anchor=west},
  every axis y label/.style={at={(ticklabel* cs:1)},anchor=south},
]
\addplot[popcurve,domain=0:0.15,samples=50]{200*x};
\addplot[fill=poporange!40,draw=none] coordinates {(0.05,0) (0.05,10) (0.10,20) (0.10,0)} \closedcycle;
\draw[dashed,popmuted] (axis cs:0.05,0) -- (axis cs:0.05,10) -- (axis cs:0,10);
\draw[dashed,popmuted] (axis cs:0.10,0) -- (axis cs:0.10,20) -- (axis cs:0,20);
\node[popdark] at (axis cs:0.075,6) {$\mathcal{A} = |W|$};
\node[popblue,anchor=south west] at (axis cs:0.105,24) {$T = kx$};
\end{axis}
\end{tikzpicture}
\legende{Courbe $T = kx$ (ici $k = \SI{200}{N.m^{-1}}$). L'aire du trapèze hachuré entre $x_A$ et $x_B$ représente la valeur absolue du travail de la tension.}
\end{center}
\end{popfigure}

\begin{popfigure}
\begin{center}
\begin{tikzpicture}[scale=0.9]
% État 1 : repos
\draw[thick,popmuted] (0,0) -- (0,1.2);
\draw[thick,popblue,decorate,decoration={coil,aspect=0.4,segment length=4pt,amplitude=4pt}] (0,0.6) -- (2,0.6);
\draw[fill=popblue!60] (2,0.4) rectangle (2.5,0.8);
\draw[<->,popdark] (0,-0.3) -- (2.25,-0.3) node[midway,below]{$\ell_0$ (repos)};
\node[popdark] at (1.25,1.5) {\textbf{Repos : } $x = 0$};
% État 2 : allongé x_A
\draw[thick,popmuted] (4.5,0) -- (4.5,1.2);
\draw[thick,poporange,decorate,decoration={coil,aspect=0.4,segment length=6pt,amplitude=4pt}] (4.5,0.6) -- (7.2,0.6);
\draw[fill=poporange!60] (7.2,0.4) rectangle (7.7,0.8);
\draw[<->,popdark] (4.5,-0.3) -- (7.45,-0.3) node[midway,below]{$\ell_A = \ell_0 + x_A$};
\draw[->,thick,popgreen] (7.45,0.6) -- (6.6,0.6) node[midway,above]{$\vec{T}_A$};
\node[popdark] at (6.1,1.5) {\textbf{Allongé : } $x = x_A$};
% État 3 : allongé x_B
\draw[thick,popmuted] (9.5,0) -- (9.5,1.2);
\draw[thick,poppurple,decorate,decoration={coil,aspect=0.4,segment length=8pt,amplitude=4pt}] (9.5,0.6) -- (13,0.6);
\draw[fill=poppurple!60] (13,0.4) rectangle (13.5,0.8);
\draw[<->,popdark] (9.5,-0.3) -- (13.25,-0.3) node[midway,below]{$\ell_B = \ell_0 + x_B$};
\draw[->,thick,popgreen] (13.25,0.6) -- (12.2,0.6) node[midway,above]{$\vec{T}_B$};
\node[popdark] at (11.5,1.5) {\textbf{Allongé : } $x = x_B$};
\end{tikzpicture}
\legende{Ressort horizontal dans trois états : au repos, allongé de $x_A$, allongé de $x_B$. La tension de rappel $\vec{T}$ croît avec l'allongement et s'oppose à l'étirement.}
\end{center}
\end{popfigure}

\begin{aretenir}[Formule du travail de la tension]
Pour un ressort à réponse linéaire de raideur $k$, dont l'extrémité passe de l'allongement $x_A$ à l'allongement $x_B$ :
\[
\boxed{\;W_{AB}(\vec{T}) = \frac{1}{2}\,k\left(x_A^2 - x_B^2\right)\;}
\]
avec $W$ en joules (J), $k$ en \SI{}{N.m^{-1}}, $x_A$ et $x_B$ en mètres, mesurés \textbf{depuis la position de repos}.
\begin{itemize}
\item Si $|x_B| > |x_A|$ (on éloigne de la position de repos) : $W < 0$, travail \textbf{résistant}.
\item Si $|x_B| < |x_A|$ (le ressort se rapproche du repos) : $W > 0$, travail \textbf{moteur}.
\end{itemize}
\end{aretenir}

\courssub{Cas particulier : étirement depuis la position de repos}

Situation très fréquente au Probatoire : le ressort part de sa position de repos ($x_A = 0$) et on l'étire jusqu'à l'allongement $x$ ($x_B = x$). La formule générale donne immédiatement :
\[
W(\vec{T}) = \frac{1}{2}k\left(0 - x^2\right) = -\frac{1}{2}\,k\,x^2
\]
Le travail de la tension est négatif : le ressort \emph{résiste} à l'étirement. Remarque que ce travail vaut l'opposé de l'aire du \emph{triangle} sous la droite $T = kx$ entre $0$ et $x$ : $\mathcal{A} = \frac{1}{2} \times \text{base} \times \text{hauteur} = \frac{1}{2} x \cdot kx = \frac{1}{2}kx^2$. On retrouve bien le même résultat, le trapèze étant ici dégénéré en triangle.

\begin{attention}[Erreur fréquente n°1 : $W = T \times x$ avec la tension finale]
Beaucoup d'élèves écrivent $W = T_{\text{final}} \times x = kx \times x = kx^2$. C'est \textbf{faux} : la tension n'a pas valu $kx$ pendant tout l'étirement ! Elle est passée continûment de $0$ à $kx$. La bonne valeur à utiliser dans un produit « force $\times$ déplacement » serait la \textbf{valeur moyenne} de la tension :
\[
T_{\text{moy}} = \frac{0 + kx}{2} = \frac{1}{2}kx \quad\Longrightarrow\quad |W| = T_{\text{moy}} \times x = \frac{1}{2}kx^2
\]
Le facteur $\frac{1}{2}$ n'est pas négociable : il vient de la variation linéaire de la force. Retenir : \emph{force variable $\Rightarrow$ force moyenne $\Rightarrow$ facteur $\frac{1}{2}$}.
\end{attention}

\begin{exemplebox}[Étirement d'un ressort de raideur $k = \SI{200}{N.m^{-1}}$]
Un opérateur étire lentement un ressort horizontal de raideur $k = \SI{200}{N.m^{-1}}$ depuis sa position de repos jusqu'à l'allongement $x = \SI{10}{cm} = \SI{0,10}{m}$.

\textbf{Travail de la tension :}
\[
W(\vec{T}) = -\frac{1}{2}kx^2 = -\frac{1}{2}\times 200 \times (0,10)^2 = -\SI{1,0}{J}
\]

\textbf{Travail de l'opérateur :} l'étirement étant lent (quasi-statique), la force $\vec{F}$ de l'opérateur est à chaque instant opposée à la tension : $W(\vec{F}) = -W(\vec{T}) = +\SI{1,0}{J}$.

L'opérateur doit donc fournir un travail moteur de \SI{1,0}{J}. Vérification par la force moyenne : $F_{\text{moy}} = \frac{1}{2}\times 200 \times 0,10 = \SI{10}{N}$, d'où $W = 10 \times 0,10 = \SI{1,0}{J}$. Cohérent !
\end{exemplebox}

\courssub{La tension du ressort est une force conservative}

Regarde attentivement la formule $W_{AB}(\vec{T}) = \frac{1}{2}k(x_A^2 - x_B^2)$ : le travail ne fait intervenir \textbf{que les allongements initial et final}. Peu importe le chemin suivi par l'extrémité du ressort — qu'on l'ait étirée directement de $x_A$ à $x_B$, ou qu'on ait fait des allers-retours intermédiaires — le travail total est le même. C'est exactement le critère défini à la section précédente.

\begin{propriete}[Conservativité de la tension]
La tension d'un ressort à réponse linéaire est une \cle{force conservative} : son travail entre deux états ne dépend que des allongements initial $x_A$ et final $x_B$, et non du chemin suivi. En particulier, sur tout cycle fermé (retour à l'allongement de départ), le travail total de la tension est nul :
\[
W_{\text{cycle}}(\vec{T}) = \frac{1}{2}k\left(x_A^2 - x_A^2\right) = 0
\]
\end{propriete}

\textbf{Interprétation énergétique (complément Probatoire).} Quand l'opérateur étire lentement le ressort de $0$ à $x$, il fournit le travail $W_{\text{op}} = +\frac{1}{2}kx^2$. Ce travail n'est pas perdu : il est \emph{stocké} dans le ressort déformé, sous forme d'\cle{énergie potentielle élastique} :
\[
E_{pe} = \frac{1}{2}\,k\,x^2
\]
Cette énergie pourra être restituée intégralement quand le ressort reviendra à sa position de repos (c'est le principe du lance-pierre, du ressort de montre ou du trampoline). On retrouve ici l'idée générale associée aux forces conservatives : le travail d'une force conservative s'accompagne d'une variation d'énergie potentielle, $\Delta E_p = -W(\vec{F}_{\text{cons}})$, exactement comme pour l'énergie potentielle de pesanteur $E_{pp} = mgz$ associée au poids.

\begin{savaistu}[Des ressorts partout autour de nous]
Les suspensions des motos-taxis — les célèbres \emph{benskin} de Douala et Yaoundé — utilisent des ressorts hélicoïdaux : chaque nid-de-poule comprime le ressort, qui stocke l'énergie $\frac{1}{2}kx^2$ puis la restitue en douceur, protégeant le passager. De même, les balances à ressort que l'on voit dans les marchés (pour peser le poisson fumé ou les arachides) exploitent directement la loi de Hooke : l'allongement $x = \dfrac{mg}{k}$ étant proportionnel à la masse, une simple graduation linéaire suffit. C'est la réponse linéaire du ressort qui rend ces instruments si simples et si fiables !
\end{savaistu}

\begin{exoresolu}[Ressort de suspension d'une moto]
La suspension arrière d'une moto-taxi est modélisée par un ressort de raideur $k = \SI{5000}{N.m^{-1}}$. Sous l'effet d'un choc, le ressort passe d'une compression $x_A = -\SI{4,0}{cm}$ à une compression $x_B = -\SI{10,0}{cm}$.

\begin{enumerate}
\item Calculer le travail de la tension du ressort au cours de cette compression. Commenter son signe.
\item Calculer le travail de la tension lorsque le ressort revient ensuite de $x_B$ à $x_A$.
\item Conclure sur le travail total au cours du cycle.
\end{enumerate}

\tcblower

\textbf{1.} On convertit : $x_A = -\SI{0,040}{m}$ et $x_B = -\SI{0,100}{m}$. Application de la formule :
\begin{align*}
W_{AB}(\vec{T}) &= \frac{1}{2}k\left(x_A^2 - x_B^2\right) = \frac{1}{2}\times 5000 \times \left[(0,040)^2 - (0,100)^2\right] \\
&= 2500 \times (0,0016 - 0,0100) = 2500 \times (-0,0084) = -\SI{21}{J}
\end{align*}
Le travail est \textbf{négatif} : le ressort s'oppose à la compression (travail résistant), ce qui est normal puisque la tension pousse vers la position de repos alors que le déplacement s'en éloigne. Notons que seuls les \emph{carrés} des allongements interviennent : les signes négatifs de $x_A$ et $x_B$ disparaissent.

\textbf{2.} Pour le retour de $x_B$ vers $x_A$ :
\[
W_{BA}(\vec{T}) = \frac{1}{2}k\left(x_B^2 - x_A^2\right) = -W_{AB}(\vec{T}) = +\SI{21}{J}
\]
Le travail est \textbf{moteur} : le ressort restitue l'énergie stockée en repoussant la suspension.

\textbf{3.} Travail total sur le cycle :
\[
W_{\text{cycle}} = W_{AB} + W_{BA} = -21 + 21 = \SI{0}{J}
\]
On vérifie ainsi concrètement que la tension est une \textbf{force conservative} : son travail sur un chemin fermé est nul. L'énergie absorbée lors du choc est intégralement restituée au retour.
\end{exoresolu}

\begin{exercice}[Pour t'entraîner]
Un ressort de longueur à vide $\ell_0 = \SI{20}{cm}$ et de raideur $k = \SI{150}{N.m^{-1}}$ est étiré horizontalement. Son extrémité libre passe de la longueur $\ell_A = \SI{26}{cm}$ à la longueur $\ell_B = \SI{34}{cm}$.
\begin{enumerate}
\item Exprimer les allongements $x_A$ et $x_B$ en mètres.
\item Calculer le travail de la tension du ressort entre ces deux états.
\item Quel travail minimal l'opérateur doit-il fournir pour réaliser cet étirement lentement ?
\end{enumerate}
\end{exercice}

\begin{corrige}[Exercice : ressort étiré entre deux longueurs]
\textbf{1.} Les allongements se mesurent depuis la position de repos :
\[
x_A = \ell_A - \ell_0 = 26 - 20 = \SI{6}{cm} = \SI{0,06}{m} \qquad x_B = \ell_B - \ell_0 = \SI{14}{cm} = \SI{0,14}{m}
\]

\textbf{2.} Travail de la tension :
\begin{align*}
W_{AB}(\vec{T}) &= \frac{1}{2}k\left(x_A^2 - x_B^2\right) = \frac{1}{2}\times 150 \times \left[(0,06)^2 - (0,14)^2\right] \\
&= 75 \times (0,0036 - 0,0196) = 75 \times (-0,016) = -\SI{1,2}{J}
\end{align*}
Travail résistant, comme attendu pour un étirement supplémentaire.

\textbf{3.} Étirement lent : la force de l'opérateur compense la tension à chaque instant, donc
\[
W_{\text{op}} = -W_{AB}(\vec{T}) = +\SI{1,2}{J}
\]
L'opérateur fournit \SI{1,2}{J}, énergie qui vient s'ajouter à l'énergie potentielle élastique déjà stockée dans le ressort.
\end{corrige}

\begin{aretenir}[L'essentiel de la section]
\begin{itemize}
\item La tension d'un ressort à réponse linéaire obéit à la loi de Hooke : $T = k|x|$, avec $x$ allongement algébrique mesuré \textbf{depuis la position de repos}.
\item Force variable $\Rightarrow$ méthode graphique : le travail est l'aire (trapèze ou triangle) sous la droite $T = k|x|$.
\item Formule à connaître par cœur : $W_{AB}(\vec{T}) = \dfrac{1}{2}k\left(x_A^2 - x_B^2\right)$.
\item Depuis le repos : $W(\vec{T}) = -\dfrac{1}{2}kx^2$ (travail résistant) ; l'opérateur fournit $+\dfrac{1}{2}kx^2$.
\item La tension est \textbf{conservative} : travail indépendant du chemin, nul sur un cycle, et associée à l'énergie potentielle élastique $E_{pe} = \dfrac{1}{2}kx^2$.
\end{itemize}
\end{aretenir}

\coursec{Travail et moment d'une force (rotation)}

Dans la section précédente, nous avons étudié le travail de la tension d'un ressort, une force dont l'intensité varie au cours du déplacement. Nous avons vu que le travail se calcule alors par une sommation (ou une intégrale) le long de la trajectoire. Dans tous les cas rencontrés jusqu'ici, le point d'application de la force se déplaçait en \emph{translation} : le long d'une droite ou d'une courbe quelconque.

Or, dans la vie courante et dans les machines, de nombreuses forces agissent sur des solides en \emph{rotation} autour d'un axe fixe : la main qui actionne la manivelle d'un treuil au chantier, la clé qui serre un écrou dans un garage de Douala, la pédale d'une bicyclette, le volant d'une machine. Comment exprimer le travail d'une force dans ce cas ? C'est l'objet de cette section, qui aboutira à une relation remarquablement simple : $W = \mathcal{M}(\vec{F}) \cdot \Delta\theta$.

\courssub{Situation d'étude : solide en rotation autour d'un axe fixe}

Considérons un solide (disque, poulie, treuil) mobile autour d'un axe fixe $(\Delta)$ passant par un point $O$ et perpendiculaire au plan du mouvement. Une force $\vec{F}$ est appliquée en un point $M$ du solide, situé à la distance $R = OM$ de l'axe.

Lorsque le solide tourne d'un angle $\Delta\theta$, le point $M$ décrit un \cle{arc de cercle} de rayon $R$ et d'angle au centre $\Delta\theta$. La longueur de cet arc est :
\[
s = R \, \Delta\theta \qquad (\Delta\theta \text{ en radians}).
\]

\begin{popfigure}
\begin{center}
\begin{tikzpicture}[scale=1.5]
% disque
\fill[popblueL] (0,0) circle (2);
\draw[popdark, thick] (0,0) circle (2);
% axe
\fill[popdark] (0,0) circle (0.06);
\node[below left, popdark] at (0,0) {$O$};
\draw[popdark] (0,-0.18) -- (0,0.18);
\draw[popdark] (-0.18,0) -- (0.18,0);
\node[above right=2pt, popdark] at (0.05,0.1) {$(\Delta)$};
% rayon initial et final
\draw[popblue, thick, -{Stealth}] (0,0) -- (30:2) node[midway, above left] {$R$};
\draw[popblue, dashed, thick] (0,0) -- (75:2);
% points M et M'
\fill[popink] (30:2) circle (0.07);
\node[right, popink] at (30:2) {$M$};
\fill[popink] (75:2) circle (0.07);
\node[above, popink] at (75:2) {$M'$};
% arc parcouru
\draw[poporange, very thick, -{Stealth}] (30:2) arc (30:75:2);
\node[poporange] at (52:2.45) {$s = R\,\Delta\theta$};
% angle
\draw[popgreen, thick, -{Stealth}] (30:0.9) arc (30:75:0.9);
\node[popgreen] at (52:0.62) {$\Delta\theta$};
% force tangente en M
\draw[popink, very thick, -{Stealth}] (30:2) -- ++(120:1.3) node[above right] {$\vec{F}$};
% angle alpha entre F et tangente (ici tangente) : montrer alpha générique
\draw[popmuted, dashed] (30:2) -- ++(120:0.4);
\end{tikzpicture}
\end{center}
\legende{Solide en rotation autour de l'axe $(\Delta)$ : le point d'application $M$ de la force $\vec{F}$ décrit l'arc de cercle $\wideparen{MM'}$ de longueur $s = R\,\Delta\theta$. Ici la force est tangente au cercle ($\alpha = 0$).}
\end{popfigure}

\courssub{Établissement de la relation $dW = \mathcal{M}(\vec{F}) \cdot d\theta$}

\textbf{Travail élémentaire.} Pendant une rotation infiniment petite $d\theta$, le point $M$ se déplace de $ds = R\,d\theta$ le long de la tangente au cercle. Notons $\alpha$ l'angle entre la force $\vec{F}$ et la direction du déplacement (la tangente au cercle en $M$, orientée dans le sens du mouvement). Le travail élémentaire de $\vec{F}$ est :
\[
dW = F \cdot ds \cdot \cos\alpha = F \cdot R\, d\theta \cdot \cos\alpha.
\]

\textbf{Apparition du moment.} Rappelons l'expression du \cle{moment} de la force $\vec{F}$ par rapport à l'axe $(\Delta)$ :
\[
\mathcal{M}_{\Delta}(\vec{F}) = F \cdot d,
\]
où $d$ est le \emph{bras de levier}, distance de l'axe à la ligne d'action de $\vec{F}$. Or, géométriquement, le bras de levier vaut $d = R\cos\alpha$ lorsque $\alpha$ est l'angle entre $\vec{F}$ et la tangente (la tangente est perpendiculaire au rayon $OM$). En effet, la composante de $\vec{F}$ utile pour la rotation est sa composante tangentielle $F\cos\alpha$, et le moment s'écrit :
\[
\mathcal{M}_{\Delta}(\vec{F}) = F \cos\alpha \cdot R.
\]

En reportant dans l'expression du travail élémentaire, il vient :
\[
\boxed{\,dW = \mathcal{M}_{\Delta}(\vec{F}) \cdot d\theta\,}
\]

Cette relation est fondamentale : \textbf{le travail élémentaire d'une force appliquée à un solide en rotation est le produit du moment de la force par l'angle de rotation élémentaire}, exprimé en radians.

\begin{attention}[Le radian est obligatoire]
La relation $dW = \mathcal{M} \cdot d\theta$ découle de $ds = R\,d\theta$, qui n'est vraie que si l'angle est exprimé en \cle{radians}. Si tu exprimes $\Delta\theta$ en degrés, le résultat sera faux ! Conversions à connaître par cœur :
\[
180^{\circ} = \pi \ \text{rad} \ ; \qquad 1 \ \text{tour} = 360^{\circ} = 2\pi \ \text{rad} \ ; \qquad 1 \ \text{rad} \approx 57,3^{\circ}.
\]
Pour convertir un nombre de tours $n$ en radians : $\Delta\theta = 2\pi n$.
\end{attention}

\courssub{Relation fondamentale : travail d'une force de moment constant}

\begin{propriete}[Travail d'une force de moment constant — démonstration exigible]
Soit un solide en rotation autour d'un axe fixe $(\Delta)$, soumis à une force $\vec{F}$ dont le moment par rapport à $(\Delta)$ est \textbf{constant} : $\mathcal{M}_{\Delta}(\vec{F}) = \mathcal{M}$. Lorsque le solide tourne d'un angle $\Delta\theta$ (en radians), le travail de $\vec{F}$ est :
\[
W = \mathcal{M}_{\Delta}(\vec{F}) \cdot \Delta\theta
\]
avec $\mathcal{M}$ en newton-mètre (\si{N.m}), $\Delta\theta$ en radian (\si{rad}) et $W$ en joule (\si{J}).

\textbf{Démonstration.} Le travail total est la somme des travaux élémentaires le long de la rotation :
\[
W = \int_{\theta_1}^{\theta_2} \mathcal{M}_{\Delta}(\vec{F}) \, d\theta.
\]
Le moment étant constant, il sort de l'intégrale :
\[
W = \mathcal{M} \int_{\theta_1}^{\theta_2} d\theta = \mathcal{M} \, (\theta_2 - \theta_1) = \mathcal{M} \cdot \Delta\theta. \qquad \square
\]
\end{propriete}

\begin{aretenir}[Conventions de signe et cohérence des unités]
\begin{itemize}
\item Le travail est \textbf{moteur} ($W > 0$) si le moment et la rotation sont de même sens ; il est \textbf{résistant} ($W < 0$) si le moment s'oppose à la rotation (frottements, charge).
\item Vérification dimensionnelle : $[\mathcal{M}] = \si{N.m} = \si{kg.m^2.s^{-2}}$ et le radian est sans dimension, donc $[\mathcal{M} \cdot \Delta\theta] = \si{kg.m^2.s^{-2}} = \si{J}$. La relation est homogène.
\item $\mathcal{M} \cdot \Delta\theta$ n'est autre que la généralisation à la rotation de $F \cdot s$ pour la translation : le \emph{moment} joue le rôle de la force, l'\emph{angle} celui du déplacement.
\end{itemize}
\end{aretenir}

\begin{exemplebox}[Serrage d'un écrou à la clé]
Un mécanicien exerce une force $F = \SI{50}{N}$ perpendiculairement à une clé de longueur $L = \SI{30}{cm}$, et fait tourner l'écrou d'un quart de tour.
\begin{itemize}
\item Moment exercé : $\mathcal{M} = F \times L = 50 \times 0,30 = \SI{15}{N.m}$.
\item Angle : $\Delta\theta = \dfrac{1}{4} \times 2\pi = \dfrac{\pi}{2} \approx \SI{1,57}{rad}$.
\item Travail : $W = \mathcal{M} \cdot \Delta\theta = 15 \times 1,57 \approx \SI{23,6}{J}$.
\end{itemize}
\end{exemplebox}

\courssub{Travail d'un couple de forces}

\begin{definition}[Couple de forces]
Un \cle{couple} est un ensemble de deux forces $\vec{F}_1$ et $\vec{F}_2$ de même intensité $F$, de directions parallèles, de sens opposés, n'ayant pas la même ligne d'action. Le moment du couple est indépendant de l'axe :
\[
\mathcal{M}_{\text{couple}} = F \cdot d,
\]
où $d$ est la distance entre les deux lignes d'action.
\end{definition}

Un couple ne provoque qu'une rotation (sa résultante est nulle : pas de translation). Son travail lors d'une rotation $\Delta\theta$ s'obtient en sommant les travaux des deux forces ; on montre que :

\begin{propriete}[Travail d'un couple]
Pour un couple de moment constant $\mathcal{M}_{\text{couple}}$ appliqué à un solide tournant de $\Delta\theta$ (rad) :
\[
W = \mathcal{M}_{\text{couple}} \cdot \Delta\theta.
\]
C'est la même relation que pour une force unique, avec le moment du couple. Exemple : les deux mains sur un volant de voiture ou sur le guidon d'une mototaxi forment un couple.
\end{propriete}

\courssub{Application : le treuil et la manivelle}

\begin{popfigure}
\begin{center}
\begin{tikzpicture}[scale=1.15]
% tambour
\fill[popgoldL] (0,0) circle (0.55);
\draw[popdark, thick] (0,0) circle (0.55);
\fill[popdark] (0,0) circle (0.05);
\node[below left=1pt, popdark] at (-0.05,-0.05) {$O$};
% manivelle
\draw[popdark, very thick] (0,0) -- (40:2.2);
\draw[popdark, thick] (40:2.2) -- ++(-50:0.5);
\fill[popink] (40:2.2) circle (0.06);
% force sur la poignée
\draw[popink, very thick, -{Stealth}] (40:2.2) ++(-50:0.5) -- ++(-50:1.1) node[right] {$\vec{F}$};
% cote L
\draw[popblue, -{Stealth}] (0.15,0.05) -- (40:2.05) node[midway, above left] {$L$};
% sens de rotation
\draw[popgreen, thick, -{Stealth}] (140:0.9) arc (140:60:0.9);
\node[popgreen] at (100:1.15) {$\Delta\theta$};
% corde et charge
\draw[popmuted, thick] (-0.55,0) -- (-0.55,-2.2);
\draw[popmuted, fill=popcream] (-0.95,-2.7) rectangle (-0.15,-2.2);
\node at (-0.55,-2.45) {$m$};
\draw[popink, thick, -{Stealth}] (-0.55,-2.7) -- (-0.55,-3.3) node[below] {$\vec{P}$};
% rayon du tambour
\draw[popblue, -{Stealth}] (0,0) -- (-0.55,0) node[midway, above] {$r$};
\end{tikzpicture}
\end{center}
\legende{Treuil : l'opérateur exerce la force $\vec{F}$ au bout d'une manivelle de longueur $L$ ; le tambour de rayon $r$ enroule la corde qui soulève la charge de poids $\vec{P}$.}
\end{popfigure}

\begin{methode}[Calculer le travail fourni à un treuil ou une manivelle]
\begin{enumerate}
\item \textbf{Calculer le moment} de la force motrice par rapport à l'axe : $\mathcal{M} = F \times L$ (si $\vec{F}$ est perpendiculaire à la manivelle ; sinon $\mathcal{M} = F L \sin\beta$).
\item \textbf{Convertir la rotation en radians} : pour $n$ tours, $\Delta\theta = 2\pi n$ ; pour un angle en degrés, multiplier par $\pi/180$.
\item \textbf{Appliquer la relation fondamentale} : $W = \mathcal{M} \cdot \Delta\theta$.
\item \textbf{Conclure avec le signe} : travail moteur ($W>0$) pour la force appliquée ; le poids de la charge, lui, effectue un travail résistant $W(\vec{P}) = -mgh$ où $h = r\,\Delta\theta$ est la hauteur de levée.
\end{enumerate}
\end{methode}

\begin{exoresolu}[Travail fourni à la manivelle d'un treuil]
\textbf{Énoncé.} Sur un chantier à Yaoundé, un ouvrier actionne la manivelle d'un treuil en exerçant une force constante $F = \SI{40}{N}$, perpendiculaire à la manivelle de longueur $L = \SI{25}{cm}$. Il effectue $n = 10$ tours complets.
\begin{enumerate}
\item Calculer le moment de la force exercée.
\item Calculer le travail fourni par l'ouvrier.
\item Le tambour a pour rayon $r = \SI{10}{cm}$. Quelle hauteur $h$ la charge a-t-elle parcourue ? En déduire la masse maximale théorique soulevée si tout le travail sert à lever la charge (on prendra $g = \SI{10}{N/kg}$).
\end{enumerate}
\tcblower
\textbf{Solution.}
\begin{enumerate}
\item La force est perpendiculaire à la manivelle, donc le bras de levier est $L$ :
\[
\mathcal{M} = F \times L = 40 \times 0,25 = \SI{10}{N.m}.
\]
\item L'angle total de rotation en radians :
\[
\Delta\theta = 2\pi n = 2\pi \times 10 = 20\pi \approx \SI{62,8}{rad}.
\]
Travail moteur :
\[
W = \mathcal{M} \cdot \Delta\theta = 10 \times 20\pi = 200\pi \approx \SI{628}{J}.
\]
\item La corde s'enroule sur le tambour : la longueur enroulée est $\ell = r\,\Delta\theta = 0,10 \times 62,8 = \SI{6,28}{m}$, donc $h = \SI{6,28}{m}$.

Le travail du poids lors de la montée est résistant : $W(\vec{P}) = -mgh$. Si tout le travail de l'ouvrier compense le travail du poids (rendement parfait) :
\[
W = mgh \quad \Longrightarrow \quad m = \dfrac{W}{g\,h} = \dfrac{628}{10 \times 6,28} = \SI{10}{kg}.
\]
\textbf{Remarque.} On retrouve ce résultat directement par le rapport des bras : $\dfrac{F \cdot L}{r} = \dfrac{40 \times 0,25}{0,10} = \SI{100}{N}$, soit une masse de \SI{10}{kg}. Le treuil est bien une \emph{machine simple} : il multiplie la force par le rapport $L/r = 2,5$, au prix d'un déplacement plus grand — le travail, lui, se conserve.
\end{enumerate}
\end{exoresolu}

\begin{attention}[Pièges fréquents au Probatoire]
\begin{itemize}
\item \textbf{Angle en degrés} : c'est l'erreur n°1. $W = \mathcal{M}\,\Delta\theta$ exige $\Delta\theta$ en \textbf{radians}. Un élève qui écrit $W = 10 \times 3600$ pour 10 tours obtient un résultat 57 fois trop grand !
\item \textbf{Bras de levier mal identifié} : le moment n'est égal à $F \times L$ que si $\vec{F}$ est \emph{perpendiculaire} à la manivelle. Sinon, $\mathcal{M} = F L \sin\beta$.
\item \textbf{Confusion moment / travail} : le moment s'exprime en \si{N.m} et le travail en joule. Bien que $\si{N.m}$ et \si{J} soient dimensionnellement identiques, on réserve le \si{N.m} au moment et le joule au travail (et à l'énergie) pour éviter toute confusion entre ces deux grandeurs de nature différente.
\item \textbf{Signe du travail} : ne pas oublier que le moment résistant (frottements, charge) donne un travail \emph{négatif}.
\end{itemize}
\end{attention}

\begin{savaistu}[Du moulin au moteur électrique]
La relation $W = \mathcal{M}\,\Delta\theta$ est au cœur de toutes les machines tournantes. Dans les barrages hydroélectriques d'Edéa, de Song Loulou ou de Nachtigal, l'eau exerce sur les pales de la turbine un couple de moment $\mathcal{M}$ ; le travail récupéré par tour est $\mathcal{M} \times 2\pi$, et c'est ce travail que l'alternateur convertit en énergie électrique distribuée par ENEO. À l'inverse, dans un moteur électrique (ventilateur, pompe de forage en zone rurale), c'est le courant qui crée un couple moteur : le moteur fournit le travail $\mathcal{M}\,\Delta\theta$ à l'arbre. Nous verrons dans la section suivante que la \emph{puissance} de ces machines s'exprime alors très simplement : $P = \mathcal{M}\,\omega$, où $\omega$ est la vitesse angulaire.
\end{savaistu}

\begin{exercice}[Volant d'inertie et couple de frottement]
Un volant de machine tourne autour de son axe. Un mécanicien exerce sur sa jante, de rayon $R = \SI{40}{cm}$, une force tangentielle $F = \SI{60}{N}$ pendant $5$ tours. Simultanément, les frottements exercent un couple résistant constant de moment $\mathcal{M}_f = \SI{4}{N.m}$.
\begin{enumerate}
\item Calculer le travail moteur fourni par le mécanicien.
\item Calculer le travail du couple de frottement sur la même rotation.
\item En déduire le travail total reçu par le volant. Que devient ce travail si le volant tourne à vitesse constante ?
\end{enumerate}
\end{exercice}

\begin{corrige}[Exercice : volant d'inertie]
\begin{enumerate}
\item Moment moteur : $\mathcal{M} = F \times R = 60 \times 0,40 = \SI{24}{N.m}$. Angle : $\Delta\theta = 5 \times 2\pi = 10\pi \approx \SI{31,4}{rad}$. Travail moteur :
\[
W_m = \mathcal{M}\,\Delta\theta = 24 \times 31,4 \approx \SI{754}{J}.
\]
\item Le couple de frottement s'oppose à la rotation :
\[
W_f = -\mathcal{M}_f \, \Delta\theta = -4 \times 31,4 \approx -\SI{126}{J}.
\]
\item Travail total : $W = W_m + W_f = 754 - 126 \approx \SI{628}{J}$. Si le volant tourne à vitesse angulaire constante, son énergie cinétique ne varie pas : ce travail net est intégralement transmis à la machine entraînée (ou dissipé ailleurs) ; en régime permanent, le travail moteur compense exactement l'ensemble des travaux résistants.
\end{enumerate}
\end{corrige}

\begin{aretenir}[L'essentiel de la section]
\begin{itemize}
\item Pour un solide en rotation autour d'un axe fixe, le point d'application de la force décrit un arc de cercle : $ds = R\,d\theta$.
\item Le travail élémentaire s'écrit $dW = \mathcal{M}_{\Delta}(\vec{F}) \cdot d\theta$.
\item Pour un moment \textbf{constant} : $W = \mathcal{M}_{\Delta}(\vec{F}) \cdot \Delta\theta$, avec $\mathcal{M}$ en \si{N.m}, $\Delta\theta$ en \textbf{radians}, $W$ en joules.
\item Cette relation s'applique aussi au \textbf{couple} : $W = \mathcal{M}_{\text{couple}} \cdot \Delta\theta$.
\item Conversions indispensables : $1$ tour $= 2\pi$ rad ; degrés $\times \dfrac{\pi}{180} =$ radians.
\item Le travail est moteur si moment et rotation sont de même sens, résistant sinon.
\end{itemize}
\end{aretenir}

\coursec{Puissance d'une force}

Dans la section précédente, nous avons vu que le travail d'une force de moment constant appliquée à un solide en rotation s'écrit $W = \mathcal{M}(\vec{F})\,\Delta\theta$. Nous savons donc désormais \emph{quantifier} l'effet énergétique d'une force, en translation comme en rotation. Mais une question pratique demeure, que tout chantier de construction à Yaoundé ou à Douala pose quotidiennement : \textbf{en combien de temps} ce travail est-il fourni ? Deux grues peuvent accomplir exactement le même travail, l'une en dix secondes, l'autre en dix minutes. Pour l'ingénieur, elles ne se valent pas. Il nous faut une nouvelle grandeur qui mesure la \emph{rapidité} avec laquelle un travail est effectué : c'est la \cle{puissance}.

\courssub{Le besoin d'une grandeur de rapidité}

\begin{experience}[Deux grues sur un chantier]
Sur le chantier d'un immeuble, deux grues soulèvent chacune une palette de ciment de masse $m = 500$ kg d'une hauteur $h = 12$ m.
\begin{itemize}
\item La grue A effectue ce travail en $\Delta t_1 = 30$ s ;
\item la grue B effectue le même travail en $\Delta t_2 = 120$ s.
\end{itemize}
Le travail \textbf{fourni par la grue} (force motrice compensant le poids à vitesse constante) est le même dans les deux cas : $W = mgh = 500 \times 9{,}8 \times 12 = \SI{58800}{J}$. Pourtant, la grue A est manifestement plus « efficace » : elle fournit la même énergie quatre fois plus vite. Le travail seul ne suffit donc pas à caractériser une machine ; il faut le rapporter à la durée.
\end{experience}

\courssub{Puissance moyenne d'une force}

\begin{definition}[Puissance moyenne d'une force]
La \cle{puissance moyenne} d'une force dont le point d'application effectue un travail $W$ pendant la durée $\Delta t$ est :
\[ P_m = \dfrac{W}{\Delta t} \]
\begin{itemize}
\item $W$ : travail de la force, en joules (J) ;
\item $\Delta t$ : durée, en secondes (s) ;
\item $P_m$ : puissance moyenne, en \cle{watts} (W).
\end{itemize}
L'unité SI de puissance est le watt : $\SI{1}{W} = \SI{1}{J/s}$.
\end{definition}

Vérifions l'homogénéité : $[P] = \dfrac{[W]}{[t]} = \dfrac{\text{kg}\cdot\text{m}^2\cdot\text{s}^{-2}}{\text{s}} = \text{kg}\cdot\text{m}^2\cdot\text{s}^{-3}$, ce qui correspond bien au watt.

\begin{aretenir}[Unités et multiples usuels]
\begin{itemize}
\item Le kilowatt : $\SI{1}{kW} = \SI{e3}{W}$ (moteurs, appareils électriques) ;
\item le mégawatt : $\SI{1}{MW} = \SI{e6}{W}$ (centrales électriques) ;
\item le \cle{cheval-vapeur} : $\SI{1}{ch} = \SI{736}{W}$, encore utilisé pour caractériser les moteurs d'automobiles et de motos.
\end{itemize}
\end{aretenir}

\begin{exemplebox}[Puissance d'une grue]
Reprenons la grue A : elle soulève $m = 500$ kg de $h = 12$ m en $\Delta t = 30$ s.
\begin{align*}
W &= mgh = 500 \times 9{,}8 \times 12 = \SI{58800}{J} \\
P &= \dfrac{W}{\Delta t} = \dfrac{58800}{30} = \SI{1960}{W} \approx \SI{2}{kW}
\end{align*}
La grue B, qui met 120 s, ne développe que $P = \dfrac{58800}{120} = \SI{490}{W}$, soit quatre fois moins.
\end{exemplebox}

\begin{popfigure}
\begin{center}
\begin{tikzpicture}
\begin{axis}[
width=0.85\linewidth, height=6cm,
xlabel={Temps $t$ (s)}, ylabel={Travail $W$ (J)},
xmin=0, xmax=130, ymin=0, ymax=65000,
xtick={0,30,60,90,120},
grid=major, grid style={popline!40},
legend pos=north west, legend style={font=\small}
]
\addplot[popcurve, popblue, very thick, domain=0:30, samples=10]{1960*x};
\addplot[popcurve, poporange, very thick, domain=0:120, samples=10]{490*x};
\addplot[dashed, popmuted, domain=0:130]{58800};
\legend{Grue A : $P_1 = \SI{1960}{W}$, Grue B : $P_2 = \SI{490}{W}$}
\node[font=\small, popmuted] at (axis cs:75,61000) {$W = \SI{58800}{J}$};
\end{axis}
\end{tikzpicture}
\end{center}
\legende{Travail en fonction du temps pour les deux grues. Le même travail $W$ est atteint, mais la pente de la droite $W(t)$ — c'est-à-dire la puissance — est quatre fois plus grande pour la grue A.}
\end{popfigure}

\courssub{Puissance instantanée}

La puissance moyenne ne renseigne que sur une durée finie. Pour connaître la puissance à un instant précis, on fait tendre l'intervalle de temps vers zéro.

\begin{propriete}[Puissance instantanée d'une force constante (admise)]
Pour une force constante $\vec{F}$ dont le point d'application se déplace à la vitesse $\vec{v}$, la puissance instantanée est :
\[ P = F\,v\,\cos\alpha \]
où $\alpha$ est l'angle entre le vecteur force $\vec{F}$ et le vecteur vitesse $\vec{v}$.
\end{propriete}

\textbf{Établissement (formateur).} Pendant une petite durée $\Delta t$, le point d'application parcourt la distance $\Delta \ell = v\,\Delta t$. Le travail de la force constante vaut alors $W = F\,\Delta\ell\,\cos\alpha = F\,v\,\Delta t\,\cos\alpha$. En divisant par $\Delta t$ :
\[ P = \dfrac{W}{\Delta t} = F\,v\,\cos\alpha. \]

\begin{attention}[Signe de la puissance]
Comme le travail, la puissance est \emph{algébrique} :
\begin{itemize}
\item si $\alpha < 90°$ : $P > 0$, la force est motrice (elle accélère le mouvement) ;
\item si $\alpha > 90°$ : $P < 0$, la force est résistante (ex. : frottements, $\alpha = 180°$, $P = -Fv$) ;
\item si $\alpha = 90°$ : $P = 0$, la force ne travaille pas (ex. : réaction normale d'un support).
\end{itemize}
\end{attention}

\begin{exemplebox}[Mototaxi en côte]
Une mototaxi développe une force motrice $F = \SI{800}{N}$ et roule à $v = \SI{36}{km/h} = \SI{10}{m/s}$ dans la direction de la force ($\alpha = 0$). Sa puissance instantanée est :
\[ P = Fv = 800 \times 10 = \SI{8000}{W} = \SI{8}{kW} \approx \SI{10,9}{ch}. \]
\end{exemplebox}

\courssub{Lien entre travail, puissance et énergie : le kilowattheure}

De la définition $P = \dfrac{W}{\Delta t}$, on tire immédiatement :
\[ W = P\,\Delta t \]
Cette relation montre qu'une énergie (un travail) peut s'exprimer comme le produit d'une puissance par une durée. C'est l'origine d'une unité d'énergie très utilisée par ENEO sur les factures d'électricité :

\begin{definition}[Le kilowattheure]
Le \cle{kilowattheure} (kWh) est l'énergie correspondant à une puissance de $\SI{1}{kW}$ délivrée pendant une durée d'une heure :
\[ \SI{1}{kWh} = \SI{e3}{W} \times \SI{3600}{s} = \SI{3,6e6}{J}. \]
\end{definition}

\begin{attention}[Piège classique : puissance ou énergie ?]
Ne confonds jamais les deux grandeurs :
\begin{itemize}
\item la \textbf{puissance} s'exprime en watts (W, kW, MW, ch) : c'est un \emph{débit} d'énergie ;
\item le \textbf{travail} (énergie) s'exprime en joules (J) ou en kilowattheures (kWh).
\end{itemize}
Le kWh n'est \textbf{pas} une puissance, malgré le « kW » dans son nom : c'est une énergie. Dire « ce moteur consomme 5 kW par heure » est une erreur ; il faut dire « ce moteur de puissance 5 kW consomme 5 kWh en une heure de fonctionnement ».
\end{attention}

\courssub{Ordres de grandeur de puissances}

\begin{center}
\begin{tabularx}{\linewidth}{|l|X|}\hline
\rowcolor{popblueL} \textbf{Système} & \textbf{Puissance typique} \\ \hline
Porteur transportant une charge & $\approx \SI{100}{W}$ \\ \hline
Motocyclette (mototaxi) & $\approx \SI{10}{kW}$ \\ \hline
Voiture particulière & $\approx \SI{50}{kW}$ \\ \hline
Centrale hydroélectrique de Song Loulou & $\approx \SI{384}{MW}$ \\ \hline
\end{tabularx}
\end{center}

\begin{popfigure}
\begin{center}
\begin{tikzpicture}
\begin{axis}[
width=0.9\linewidth, height=6cm,
ybar, bar width=0.8cm,
ylabel={Puissance (W)}, ymode=log, log origin=infty,
ymin=10, ymax=1000000000,
symbolic x coords={Porteur, Moto, Voiture, Song Loulou},
xtick=data, x tick label style={font=\small},
grid=major, grid style={popline!40},
nodes near coords, every node near coord/.append style={font=\scriptsize, rotate=90, anchor=west}
]
\addplot[fill=popblue, draw=popdark] coordinates {(Porteur,100) (Moto,10000) (Voiture,50000) (Song Loulou,384000000)};
\end{axis}
\end{tikzpicture}
\end{center}
\legende{Ordres de grandeur de puissances (échelle logarithmique) : du porteur ($\sim \SI{e2}{W}$) à la centrale de Song Loulou ($\sim \SI{4e8}{W}$), soit six ordres de grandeur d'écart.}
\end{popfigure}

\begin{savaistu}[Pourquoi parle-t-on encore de « chevaux » ?]
À la fin du XVIII\textsuperscript{e} siècle, James Watt cherchait à vendre ses machines à vapeur aux propriétaires de mines, qui utilisaient des chevaux pour remonter l'eau et le minerai. Pour comparer, il mesura le travail qu'un cheval pouvait fournir durablement et définit le \emph{horsepower} : une machine de « 10 chevaux » remplaçait dix chevaux. L'unité est restée : en Europe, $\SI{1}{ch} = \SI{736}{W}$. C'est pourquoi la puissance « fiscale » des véhicules, qui sert notamment au calcul des taxes, s'exprime encore en chevaux, plus de deux siècles après l'invention de la machine à vapeur !
\end{savaistu}

\begin{exoresolu}[Pompage d'eau en zone rurale]
\textbf{Énoncé.} Dans un village en cours d'électrification rurale, une pompe remonte un volume $V = \SI{2}{m^3}$ d'eau d'un puits de profondeur $h = \SI{15}{m}$ en $\Delta t = 10$ min. On donne $\rho_{\text{eau}} = \SI{1000}{kg/m^3}$ et $g = \SI{9,8}{N/kg}$.
\begin{enumerate}
\item Calculer le travail fourni par la pompe.
\item En déduire sa puissance utile en watts puis en chevaux.
\item Quelle énergie (en kWh) la pompe consomme-t-elle si elle fonctionne 3 h par jour à cette puissance ?
\end{enumerate}
\tcblower
\textbf{Solution.}
\begin{enumerate}
\item La masse d'eau pompée est $m = \rho V = 1000 \times 2 = \SI{2000}{kg}$. Le travail fourni (opposé au travail du poids) :
\[ W = mgh = 2000 \times 9{,}8 \times 15 = \SI{294000}{J} = \SI{294}{kJ}. \]
\item La durée est $\Delta t = 10 \times 60 = \SI{600}{s}$. D'où :
\[ P = \dfrac{W}{\Delta t} = \dfrac{294000}{600} = \SI{490}{W} = \SI{0,49}{kW}, \]
soit $P = \dfrac{490}{736} \approx \SI{0,67}{ch}$.
\item Énergie quotidienne : $W' = P\,\Delta t' = 0{,}49 \times 3 = \SI{1,47}{kWh}$ par jour.
\end{enumerate}
\end{exoresolu}

\begin{methode}[Calculer une puissance : la démarche]
\begin{enumerate}
\item Identifier la force qui travaille et calculer (ou exprimer) son travail $W$ en joules.
\item Convertir la durée $\Delta t$ en secondes.
\item Appliquer $P = \dfrac{W}{\Delta t}$ et exprimer le résultat en watts (avec un multiple adapté : kW, MW).
\item Si la vitesse du point d'application est connue et la force constante, préférer $P = Fv\cos\alpha$.
\item Pour obtenir une énergie à partir d'une puissance, utiliser $W = P\,\Delta t$ (et convertir en kWh si demandé).
\end{enumerate}
\end{methode}

\begin{exercice}[Pour t'entraîner]
Un ascenseur de masse totale $m = 800$ kg monte de 6 étages, soit $h = \SI{18}{m}$, en $\Delta t = \SI{20}{s}$ à vitesse constante. On prend $g = \SI{9,8}{N/kg}$.
\begin{enumerate}
\item Quel est le travail de la force motrice ?
\item Quelle est la puissance du moteur ?
\item Vérifier le résultat en utilisant la relation $P = Fv\cos\alpha$.
\end{enumerate}
\end{exercice}

\begin{corrige}[Exercice : ascenseur]
\begin{enumerate}
\item À vitesse constante, la force motrice compense le poids : $F = mg = 800 \times 9{,}8 = \SI{7840}{N}$. Le travail vaut $W = Fh = 7840 \times 18 = \SI{141120}{J} \approx \SI{141}{kJ}$.
\item $P = \dfrac{W}{\Delta t} = \dfrac{141120}{20} = \SI{7056}{W} \approx \SI{7,1}{kW}$.
\item La vitesse est $v = \dfrac{h}{\Delta t} = \dfrac{18}{20} = \SI{0,9}{m/s}$. Force et vitesse sont colinéaires ($\alpha = 0$) : $P = Fv = 7840 \times 0{,}9 = \SI{7056}{W}$. Les deux méthodes concordent.
\end{enumerate}
\end{corrige}

\begin{aretenir}[L'essentiel de la section]
\begin{itemize}
\item Puissance moyenne : $P = \dfrac{W}{\Delta t}$, en watts ($\SI{1}{W} = \SI{1}{J/s}$).
\item Puissance instantanée d'une force constante : $P = Fv\cos\alpha$.
\item Énergie et puissance sont liées par $W = P\,\Delta t$ ; $\SI{1}{kWh} = \SI{3,6e6}{J}$ est une \emph{énergie}, pas une puissance.
\item Unités usuelles : kW, MW, et le cheval-vapeur ($\SI{1}{ch} = \SI{736}{W}$).
\end{itemize}
Dans la section suivante, nous transposerons ces idées au cas de la rotation : la puissance d'une force de moment constant.
\end{aretenir}

\coursec{Puissance d'une force de moment constant}

Dans la section précédente, nous avons défini la \cle{puissance} d'une force comme la rapidité avec laquelle cette force travaille : $P = \dfrac{W}{\Delta t}$, et nous avons montré que pour une force constante dont le point d'application se déplace à la vitesse $\vec{v}$, on obtient la relation remarquable $P = \vec{F}\cdot\vec{v}$, soit $P = F\,v\cos\alpha$ lorsque la force fait un angle $\alpha$ avec la vitesse. Mais toutes les machines qui nous entourent ne fonctionnent pas en translation : le moteur d'une mototaxe, la perceuse du menuisier, le ventilateur de plafond, la roue du moulin à manioc... tous tournent autour d'un axe. Comment exprimer la puissance d'une force qui fait \emph{tourner} un solide ? C'est l'objet de cette dernière section, qui va nous conduire à l'une des relations les plus utiles de la mécanique des machines : $P = M\,\omega$.

\courssub{De la translation à la rotation : l'intuition}

Considérons une situation concrète : un mécanicien serre un écrou avec une clé. Il exerce une force $\vec{F}$ à l'extrémité de la clé, à la distance $d$ de l'axe de l'écrou. La clé tourne d'un angle $\Delta\theta$ autour de l'axe. Nous avons établi dans la section 4 que le travail de cette force s'écrit :
\[ W = M_{\Delta}(\vec{F})\times\Delta\theta \]
où $M_{\Delta}(\vec{F}) = F\,d$ est le \cle{moment} de la force par rapport à l'axe de rotation $\Delta$ (en $\mathrm{N\cdot m}$) et $\Delta\theta$ l'angle de rotation, exprimé en \cle{radians}.

Remarque essentielle : il ne faut surtout pas écrire $W = M\,\Delta t$ ! Le moment se multiplie par un \emph{angle}, pas par une durée. Le travail est le produit d'un moment par un déplacement angulaire, exactement comme, en translation, le travail est le produit d'une force par un déplacement linéaire. Cette analogie va nous guider tout au long de la section.

\begin{attention}[L'angle doit être en radians]
La relation $W = M\,\Delta\theta$ n'est valable que si $\Delta\theta$ est exprimé en \textbf{radians}. Un tour complet correspond à $\Delta\theta = 2\pi\ \mathrm{rad}$, et non à $360$ ! C'est la même exigence que pour la vitesse angulaire $\omega$.
\end{attention}

\courssub{Établissement de la relation $P = M\,\omega$}

Partons de la définition de la puissance moyenne établie dans la section précédente :
\[ P = \frac{W}{\Delta t} \]
et remplaçons le travail par son expression en rotation, $W = M\,\Delta\theta$, valable pour un moment constant :
\[ P = \frac{M\,\Delta\theta}{\Delta t} = M\times\frac{\Delta\theta}{\Delta t} \]
Or le quotient $\dfrac{\Delta\theta}{\Delta t}$ n'est autre que la \cle{vitesse angulaire} $\omega$ du solide, c'est-à-dire l'angle balayé par unité de temps. Nous obtenons ainsi le résultat fondamental de cette section.

\begin{propriete}[Puissance d'une force de moment constant]
Pour un système de forces de \textbf{moment résultant constant} $M$ par rapport à un axe fixe $\Delta$, appliqué à un solide tournant autour de cet axe à la vitesse angulaire $\omega$, la puissance développée est :
\[ \boxed{P = M\,\omega} \]
avec :
\begin{itemize}
\item $P$ en \textbf{watts} ($\mathrm{W}$) ;
\item $M$ en \textbf{newtons-mètres} ($\mathrm{N\cdot m}$) ;
\item $\omega$ en \textbf{radians par seconde} ($\mathrm{rad/s}$).
\end{itemize}
\emph{Démonstration (exigible) :} pour un moment constant, $W = M\,\Delta\theta$ pendant la durée $\Delta t$. En divisant par $\Delta t$ : $P = \dfrac{W}{\Delta t} = M\,\dfrac{\Delta\theta}{\Delta t} = M\,\omega$, puisque $\omega = \dfrac{\Delta\theta}{\Delta t}$ est la vitesse angulaire. \textbf{Vérification dimensionnelle :} $[M\,\omega] = \mathrm{N\cdot m \cdot rad/s} = \mathrm{J/s} = \mathrm{W}$ (le radian est sans dimension). La relation est homogène.
\end{propriete}

\begin{aretenir}[La grande analogie]
En translation : $P = F\,v$ (force $\times$ vitesse linéaire).
En rotation : $P = M\,\omega$ (moment $\times$ vitesse angulaire).
Le moment joue en rotation le rôle que joue la force en translation ; la vitesse angulaire joue le rôle de la vitesse linéaire. Retiens ce parallélisme : il structure toute la mécanique.
\end{aretenir}

\courssub{Lien avec la fréquence de rotation}

Dans la pratique industrielle, la vitesse de rotation d'un moteur n'est presque jamais donnée en $\mathrm{rad/s}$ : on l'exprime en \textbf{tours par minute} ($\mathrm{tr/min}$, noté $N$) ou parfois en \textbf{tours par seconde} (fréquence $n$, en $\mathrm{Hz}$). Il faut donc savoir convertir.

Un tour correspond à un angle de $2\pi\ \mathrm{rad}$. Par conséquent :
\begin{itemize}
\item si le solide effectue $n$ tours par seconde : $\boxed{\omega = 2\pi n}$ avec $n$ en $\mathrm{tr/s}$ (ou $\mathrm{Hz}$) ;
\item si le solide effectue $N$ tours par minute : il effectue $\dfrac{N}{60}$ tours par seconde, donc :
\[ \boxed{\omega = \frac{2\pi N}{60}} \quad \text{avec } N \text{ en } \mathrm{tr/min}. \]
\end{itemize}

\begin{methode}[Convertir systématiquement en rad/s avant tout calcul]
Face à tout exercice donnant une vitesse de rotation, applique rigoureusement les étapes suivantes :
\begin{enumerate}
\item \textbf{Repérer} l'unité de la vitesse de rotation dans l'énoncé ($\mathrm{tr/min}$, $\mathrm{tr/s}$ ou $\mathrm{rad/s}$).
\item \textbf{Convertir} en $\mathrm{rad/s}$ : $\omega = \dfrac{2\pi N}{60}$ si $N$ est en $\mathrm{tr/min}$ ; $\omega = 2\pi n$ si $n$ est en $\mathrm{tr/s}$.
\item \textbf{Convertir} le moment en $\mathrm{N\cdot m}$ si nécessaire (attention aux $\mathrm{daN\cdot m}$, aux $\mathrm{mN\cdot m}$...).
\item \textbf{Appliquer} $P = M\,\omega$ ; le résultat est en watts.
\item \textbf{Conclure} avec l'unité et un nombre raisonnable de chiffres significatifs.
\end{enumerate}
\end{methode}

\begin{attention}[L'erreur classique du Probatoire]
L'erreur la plus fréquente consiste à injecter directement la valeur en $\mathrm{tr/min}$ dans la formule $P = M\,\omega$. C'est \textbf{faux} : $\omega$ doit être en $\mathrm{rad/s}$. Comme $\omega = \dfrac{2\pi}{60}\,N$, utiliser $N$ à la place de $\omega$ donne un résultat $\dfrac{60}{2\pi} \approx 9,5$ fois trop grand ! Un moteur de $\SI{8}{\kilo\watt}$ devient miraculeusement un moteur de $\SI{76}{\kilo\watt}$... Le correcteur, lui, ne sera pas dupe. Convertis toujours, et écris la conversion explicitement sur ta copie.
\end{attention}

\courssub{Schéma d'un moteur entraînant un arbre}

\begin{popfigure}
\begin{center}
\begin{tikzpicture}[scale=1.0, >=Stealth]
% Moteur (cylindre vu de côté)
\draw[fill=popblueL, draw=popblue, thick] (-3.2,-1.0) rectangle (-1.2,1.0);
\node[popdark] at (-2.2,0) {\small Moteur};
\draw[fill=popblue!40, draw=popblue, thick] (-1.2,-1.0) ellipse (0.18 and 1.0);
% Arbre
\draw[fill=popmuted!40, draw=popdark, thick] (-1.2,-0.22) rectangle (3.4,0.22);
\draw[fill=popmuted!60, draw=popdark, thick] (3.4,-0.22) ellipse (0.10 and 0.22);
% Axe de rotation
\draw[dashed, popdark] (-3.6,0) -- (4.1,0);
\node[popdark, right] at (4.1,0) {\small $\Delta$};
% Poulie / charge en bout d'arbre
\draw[fill=poporangeL, draw=poporange, thick] (2.6,-0.9) rectangle (3.0,0.9);
\draw[fill=poporange!40, draw=poporange, thick] (3.0,-0.9) ellipse (0.15 and 0.9);
% Vitesse angulaire : flèche courbe
\draw[->, popgreen, very thick] (1.2,0.75) arc (120:40:0.85);
\node[popgreen, above] at (1.35,1.15) {\small $\omega$ (rad/s)};
% Moment : double flèche courbe près du moteur
\draw[->, poppink, very thick] (-0.7,0.95) arc (110:20:1.05);
\node[poppink, above] at (-0.35,1.45) {\small $M$ (N$\cdot$m)};
% Puissance
\node[popdark, align=center] at (1.0,-1.6) {\small Puissance transmise à l'arbre : \quad $\boxed{P = M\,\omega}$ (W)};
\draw[->, popdark] (0.2,-1.25) -- (-0.4,-0.35);
\draw[->, popdark] (1.9,-1.25) -- (2.3,-0.35);
\end{tikzpicture}
\end{center}
\legende{Un moteur exerce sur l'arbre de rotation (axe $\Delta$) un couple de moment constant $M$ ; l'arbre tourne à la vitesse angulaire $\omega$. La puissance mécanique transmise à la charge (poulie, roue, hélice...) est $P = M\,\omega$.}
\end{popfigure}

\courssub{Exemple chiffré : puissance d'un moteur}

\begin{exemplebox}[Moteur développant un couple de $\SI{50}{\newton\meter}$ à $\SI{1500}{\tr\per\minute}$]
Un moteur électrique entraîne une machine en exerçant un couple de moment constant $M = \SI{50}{\newton\meter}$. L'arbre tourne à $N = \SI{1500}{\tr\per\minute}$. Calculons la puissance développée.

\textbf{Étape 1 : conversion de la vitesse de rotation.}
\[ \omega = \frac{2\pi N}{60} = \frac{2\pi \times 1500}{60} = 50\pi \approx \SI{157}{\radian\per\second}. \]

\textbf{Étape 2 : application de $P = M\,\omega$.}
\[ P = M\,\omega = 50 \times 157 \approx \SI{7850}{\watt} \approx \SI{7,9}{\kilo\watt}. \]

\textbf{Étape 3 : conversion en chevaux-vapeur (unité encore courante chez les garagistes).}
Sachant que $1\ \mathrm{ch} = \SI{736}{\watt}$ :
\[ P = \frac{7850}{736} \approx 10,7\ \mathrm{ch}. \]
Le moteur développe donc une puissance d'environ $\SI{7,9}{\kilo\watt}$, soit environ $10,7$ chevaux.
\end{exemplebox}

\courssub{Conséquence fondamentale : couple et vitesse varient en sens inverse}

Réécrivons la relation sous la forme :
\[ M = \frac{P}{\omega} \]
Cette écriture révèle un phénomène capital pour comprendre les machines : \textbf{à puissance donnée, le moment disponible diminue lorsque la vitesse de rotation augmente, et augmente lorsque la vitesse diminue}. Moment et vitesse angulaire sont inversement proportionnels.

C'est exactement le principe de la \cle{boîte de vitesses} d'un véhicule. Le moteur thermique fournit une puissance $P$ à peu près imposée par l'accélérateur. La boîte de vitesses, par un jeu d'engrenages, choisit le rapport entre la vitesse de rotation du moteur et celle des roues :
\begin{itemize}
\item en \textbf{première vitesse}, les roues tournent lentement ($\omega$ faible) : le moment transmis aux roues est \textbf{énorme}, ce qui permet de démarrer ou de gravir une côte raide ;
\item en \textbf{quatrième ou cinquième}, les roues tournent vite ($\omega$ grande) : le moment disponible est \textbf{faible}, mais il suffit à entretenir la vitesse sur route plate.
\end{itemize}
On ne peut pas avoir à la fois un grand couple et une grande vitesse avec une puissance limitée : c'est la loi $P = M\,\omega$ qui l'interdit, et non une question de technologie.

\begin{savaistu}[La côte de Mbanga et la première vitesse]
C'est grâce à la relation $P = M\,\omega$ qu'un camion chargé gravissant la côte de Mbanga, sur l'axe Douala--N'Kongsamba, dispose d'un couple énorme malgré une vitesse très faible. Le chauffeur rétrograde en première : la vitesse de rotation des roues chute, et comme la puissance du moteur reste disponible, le moment aux roues est multiplié par le rapport de démultiplication de la boîte. Le même principe explique pourquoi le conducteur d'une mototaxe peinant dans une montée à Bafoussam ou à Dschang doit « descendre les rapports » : à puissance constante, sacrifier la vitesse, c'est gagner du couple. Inversement, sur l'autoroute Yaoundé--Douala, on roule en grand rapport : vitesse élevée, couple modeste.
\end{savaistu}

\courssub{Exercice résolu}

\begin{exoresolu}[Treuil de chantier à Ngaoundéré]
Un treuil électrique est utilisé pour hisser des sacs de ciment sur un chantier. Son tambour, de rayon $r = \SI{20}{\centi\meter}$, est entraîné par un moteur qui exerce un couple de moment constant $M = \SI{120}{\newton\meter}$. Le tambour tourne à $N = \SI{30}{\tr\per\minute}$.
\begin{enumerate}
\item Calculer la vitesse angulaire $\omega$ du tambour en $\mathrm{rad/s}$.
\item Calculer la puissance $P$ développée par le moteur.
\item En déduire la force de traction $F$ exercée sur le câble, puis la vitesse de montée $v$ de la charge. Vérifier que $P = F\,v$.
\end{enumerate}
\tcblower
\textbf{1. Vitesse angulaire.} On convertit la vitesse de rotation en $\mathrm{rad/s}$ :
\[ \omega = \frac{2\pi N}{60} = \frac{2\pi \times 30}{60} = \pi \approx \SI{3,14}{\radian\per\second}. \]

\textbf{2. Puissance du moteur.} Le moment est constant, donc :
\[ P = M\,\omega = 120 \times 3,14 \approx \SI{377}{\watt}. \]

\textbf{3. Force de traction et vitesse de montée.} Le câble s'enroule sur le tambour à la distance $r$ de l'axe ; le moment du couple moteur s'écrit $M = F\,r$, d'où :
\[ F = \frac{M}{r} = \frac{120}{0,20} = \SI{600}{\newton}. \]
Un point de la périphérie du tambour se déplace à la vitesse $v = r\,\omega$ :
\[ v = r\,\omega = 0,20 \times 3,14 \approx \SI{0,63}{\meter\per\second}. \]
Vérification par la formule de translation :
\[ F\,v = 600 \times 0,63 \approx \SI{377}{\watt} = P. \]
Les deux approches, rotation ($P = M\,\omega$) et translation ($P = F\,v$), donnent rigoureusement le même résultat : la puissance est un invariant, seul le point de vue change.
\end{exoresolu}

\courssub{Bilan de synthèse de la leçon : translation et rotation}

Toute la leçon « Travail et puissance d'une force » peut désormais se résumer dans le tableau suivant, qui met en regard les grandeurs de la translation et celles de la rotation. Mémorise-le : il condense l'essentiel de ce qui est exigible au Probatoire.

\begin{center}
\rowcolors{1}{popcream}{white}
\begin{tabularx}{\linewidth}{|>{\raggedright\arraybackslash}X|>{\centering\arraybackslash}X|>{\centering\arraybackslash}X|}
\hline
\rowcolor{popblueL} \textbf{Grandeur} & \textbf{Translation} & \textbf{Rotation} \\
\hline
Cause du mouvement & Force $\vec{F}$ (N) & Moment $M_{\Delta}(\vec{F}) = F \times \text{bras de levier}$ ($\mathrm{N\cdot m}$) \\
\hline
Déplacement & Distance $d$ (m) & Angle $\Delta\theta$ (rad) \\
\hline
Travail & $W = F\,d\cos\alpha$ (J) & $W = M\,\Delta\theta$ (J) \\
\hline
Vitesse & $v = \dfrac{d}{\Delta t}$ (m/s) & $\omega = \dfrac{\Delta\theta}{\Delta t} = \dfrac{2\pi N}{60}$ (rad/s) \\
\hline
Puissance & $P = \dfrac{W}{\Delta t} = F\,v\cos\alpha$ (W) & $P = \dfrac{W}{\Delta t} = M\,\omega$ (W) \\
\hline
\end{tabularx}
\end{center}

\begin{popfigure}
\begin{center}
\begin{tikzpicture}[scale=0.95, >=Stealth]
% Bloc translation
\draw[fill=popblueL, draw=popblue, thick, rounded corners] (-5.2,-1.5) rectangle (-0.4,1.5);
\node[popdark, font=\bfseries] at (-2.8,1.1) {\small TRANSLATION};
\draw[fill=white, draw=popdark] (-4.6,-0.5) rectangle (-3.4,0.3);
\draw[->, poppink, very thick] (-3.4,0.1) -- (-1.9,0.6) node[midway, above, poppink] {\small $\vec{F}$};
\draw[->, popgreen, very thick] (-4.4,-0.85) -- (-1.6,-0.85) node[midway, below, popgreen] {\small $\vec{v}$};
\node[popdark] at (-2.8,-1.2) {\small $W = F\,d\cos\alpha$ \ ; \ $P = F\,v\cos\alpha$};
% Bloc rotation
\draw[fill=poporangeL, draw=poporange, thick, rounded corners] (0.4,-1.5) rectangle (5.2,1.5);
\node[popdark, font=\bfseries] at (2.8,1.1) {\small ROTATION};
\draw[fill=white, draw=popdark, thick] (2.8,-0.1) circle (0.55);
\fill[popdark] (2.8,-0.1) circle (0.05);
\draw[->, poppink, very thick] (3.9,0.55) -- (3.35,0.42) node[pos=0, right, poppink] {\small $\vec{F}$};
\draw[->, popgreen, very thick] (2.15,0.75) arc (140:60:0.75);
\node[popgreen, above] at (2.8,0.95) {\small $\omega$};
\node[popdark] at (2.8,-1.2) {\small $W = M\,\Delta\theta$ \ ; \ $P = M\,\omega$};
% Double flèche d'analogie
\draw[<->, poppurple, very thick] (-0.3,0) -- (0.3,0);
\end{tikzpicture}
\end{center}
\legende{L'analogie translation--rotation : à la force $\vec{F}$ correspond le moment $M$, au déplacement $d$ correspond l'angle $\Delta\theta$, à la vitesse $v$ correspond la vitesse angulaire $\omega$. Les formules du travail et de la puissance se transposent terme à terme.}
\end{popfigure}

\begin{aretenir}[Analogies translation--rotation à mémoriser pour le Probatoire]
\begin{itemize}
\item Force $F$ (N) $\longleftrightarrow$ Moment $M$ ($\mathrm{N\cdot m}$) ;
\item Déplacement $d$ (m) $\longleftrightarrow$ Angle $\Delta\theta$ (rad) ;
\item Vitesse $v$ (m/s) $\longleftrightarrow$ Vitesse angulaire $\omega$ (rad/s) ;
\item Travail : $W = F\,d\cos\alpha$ $\longleftrightarrow$ $W = M\,\Delta\theta$ ;
\item Puissance : $P = F\,v\cos\alpha$ $\longleftrightarrow$ $P = M\,\omega$ ;
\item Dans tous les cas : $P = \dfrac{W}{\Delta t}$, avec $W$ en joules, $\Delta t$ en secondes, $P$ en watts.
\end{itemize}
Si tu connais les formules de la translation, tu connais celles de la rotation : remplace chaque grandeur par son analogue.
\end{aretenir}

\begin{exercice}[Alternateur du barrage de Nachtigal]
Sur le site hydroélectrique de Nachtigal, un alternateur est entraîné par une turbine qui lui transmet une puissance mécanique $P = \SI{10}{\mega\watt}$. L'arbre tourne à $N = \SI{214}{\tr\per\minute}$.
\begin{enumerate}
\item Convertir la vitesse de rotation en $\mathrm{rad/s}$.
\item Calculer le moment du couple exercé par la turbine sur l'arbre de l'alternateur.
\item Si la vitesse de rotation était deux fois plus faible à puissance égale, que deviendrait le moment ? Commenter.
\end{enumerate}
\end{exercice}

\begin{corrige}[Exercice : alternateur du barrage de Nachtigal]
\textbf{1.} $\omega = \dfrac{2\pi N}{60} = \dfrac{2\pi \times 214}{60} \approx \SI{22,4}{\radian\per\second}$.

\textbf{2.} De $P = M\,\omega$, on tire $M = \dfrac{P}{\omega} = \dfrac{10\times 10^{6}}{22,4} \approx \SI{4,5e5}{\newton\meter}$, soit environ $\SI{450}{\kilo\newton\meter}$ : un couple colossal, ce qui explique la taille massive des arbres d'alternateurs.

\textbf{3.} À puissance constante, $M$ est inversement proportionnel à $\omega$ : si $\omega$ est divisée par deux, le moment est \textbf{doublé}, soit $M' \approx \SI{9,0e5}{\newton\meter}$. C'est la même loi que pour la boîte de vitesses : tourner moins vite permet de transmettre un couple plus grand.
\end{corrige}

Tu disposes maintenant de tous les outils de la leçon : le travail d'une force constante, le travail du poids et la notion de force conservative, le travail de la tension d'un ressort, le lien entre travail et moment, et enfin la puissance en translation comme en rotation. Ces relations sont la charpente de la mécanique énergétique que nous développerons dans les leçons suivantes avec le théorème de l'énergie cinétique. Entraîne-toi sur les exercices : la conversion des unités et la rigueur des schémas de forces feront la différence le jour du Probatoire.

\newpage

\coursec{Activité d'intégration}

\coursec{Travail et puissance d'une force}

\courssub{Objectifs de l'activité}

À l'issue de cette activité d'intégration, tu seras capable de :
\begin{itemize}
\item calculer le \cle{travail du poids} d'un corps soulevé et vérifier (par le calcul) que ce travail \textbf{ne dépend pas du chemin suivi} (poids = force conservative) ;
\item exploiter la relation entre \cle{travail} et \cle{moment d'une force} en rotation : $W = \mathcal{M}\,\Delta\theta$ ;
\item exprimer la \cle{puissance moyenne} d'une force : $P = \dfrac{W}{\Delta t}$ ;
\item exprimer la \cle{puissance d'une force de moment constant} : $P = \mathcal{M}\,\omega$ ;
\item comparer deux dispositifs (manivelle humaine, moteur électrique) et justifier l'intérêt mécanique d'une \textbf{machine simple} (le treuil) ;
\item \textbf{(Complément)} calculer le travail de la tension d'un ressort à réponse linéaire et identifier son caractère conservative.
\end{itemize}

\courssub{Situation}

\textbf{Le treuil du chantier de Yaoundé.}\par
Dans le quartier Nkolbisson à Yaoundé, une entreprise de BTP construit un immeuble R+2. Pour monter le béton au premier étage, faute de grue, l'ouvrier Ahmadou utilise un \textbf{treuil à manivelle} fixé sur un échafaudage : une corde s'enroule sur un \textbf{tambour} cylindrique de rayon $r = \SI{10}{cm}$, mis en rotation par une \textbf{manivelle} de longueur $L = \SI{40}{cm}$ (distance entre l'axe du tambour et la poignée). Ahmadou exerce sur la poignée une force $\vec{F}$ supposée \textbf{constante} et \textbf{toujours perpendiculaire à la manivelle}.

Il monte un seau de béton de masse $m = \SI{30}{kg}$, initialement au sol (point A), jusqu'à une plateforme située à la hauteur $h = \SI{8,0}{m}$ (point B), à vitesse pratiquement constante, en une durée $\Delta t = \SI{50}{s}$.

On donne : $g = \SI{9,8}{N.kg^{-1}}$ ; on néglige les frottements dans les axes et la masse de la corde ; le treuil est supposé parfait (rendement de 100 \%).

Le chef de chantier envisage de remplacer l'effort humain par un \textbf{moteur électrique} fournissant \textbf{le même couple} que l'ouvrier sur l'axe du tambour, mais tournant à $N = \SI{300}{tr.min^{-1}}$. Ahmadou se demande si ce moteur sera plus puissant que lui, et pourquoi sa manivelle de 40 cm le fatigue moins qu'une manivelle courte.

\courssub{Consignes}

\begin{enumerate}
\item \textbf{Travail du poids.} Exprimer puis calculer le travail $W_{AB}(\vec{P})$ du poids du seau lors de la montée de A vers B. Ce travail est-il moteur ou résistant ? Justifier le signe.
\item \textbf{Force conservative.} Un second ouvrier propose de monter le seau en le tirant le long d'une rampe inclinée reliant A à B. Sans refaire tout le calcul, que vaut le travail du poids le long de ce nouveau chemin ? Quelle propriété du poids cela illustre-t-il ?
\item \textbf{Nombre de tours.} Calculer la longueur de corde enroulée par tour de tambour, puis le nombre $n$ de tours de manivelle nécessaires pour monter le seau de $h = \SI{8,0}{m}$. En déduire l'angle total $\Delta\theta$ balayé par la manivelle (en radians).
\item \textbf{Travail de l'ouvrier.} La montée s'effectuant à vitesse constante et le treuil étant parfait, montrer que le travail fourni par Ahmadou est $W(\vec{F}) = mgh$. Le calculer.
\item \textbf{Force sur la manivelle.} En utilisant la relation $W = \mathcal{M}\,\Delta\theta$ avec $\mathcal{M} = F \times L$, déterminer l'intensité $F$ de la force exercée par l'ouvrier. Comparer $F$ au poids du seau et commenter.
\item \textbf{Puissance moyenne.} Calculer la puissance moyenne $P = \dfrac{W}{\Delta t}$ développée par Ahmadou.
\item \textbf{Comparaison avec le moteur.} Le moteur applique le même moment $\mathcal{M}$ sur l'axe et tourne à $N = \SI{300}{tr.min^{-1}}$. Convertir cette vitesse en $\si{rad.s^{-1}}$, puis calculer la puissance $P' = \mathcal{M}\,\omega$ du moteur. Comparer à la puissance de l'ouvrier et conclure.
\item \textbf{Discussion.} Expliquer, à l'aide de la relation $W = F L\,\Delta\theta$, pourquoi une manivelle plus longue fatigue moins l'ouvrier pour le même travail fourni. Quel est le « prix à payer » ?
\end{enumerate}

\courssub{Ressources à mobiliser}

\begin{aretenir}[Boîte à outils de la leçon — Rappels de cours]
\textbf{Travail d'une force constante}
\begin{itemize}
\item Définition : $W_{AB}(\vec{F}) = \vec{F} \cdot \vec{AB} = F \times AB \times \cos(\vec{F}, \vec{AB})$.
\item Unité : le joule (\si{J}) ; $1\ \si{J} = 1\ \si{N.m} = 1\ \si{kg.m^2.s^{-2}}$.
\item Travail moteur si $W > 0$ (force favorise le déplacement), résistant si $W < 0$, nul si $\vec{F} \perp \vec{AB}$.
\end{itemize}

\textbf{Travail du poids et force conservative}
\begin{itemize}
\item $W_{AB}(\vec{P}) = mg(z_A - z_B) = -mg\Delta z$. Ne dépend \textbf{que de la dénivellation} $\Delta z = z_B - z_A$, pas du chemin.
\item Le poids est une \cle{force conservative} : son travail sur un parcours fermé est nul.
\end{itemize}

\textbf{Travail de la tension d'un ressort à réponse linéaire (rappel / complément)}
\begin{itemize}
\item Force de rappel : $\vec{F} = -k\,\vec{x}$ (loi de Hooke), $k$ en $\si{N.m^{-1}}$.
\item Travail entre $x_A$ et $x_B$ : $W_{AB}(\vec{F}) = \dfrac{1}{2}k(x_A^2 - x_B^2)$.
\item La force de rappel du ressort est \textbf{conservative} : son travail ne dépend que des positions initiale et finale (écarts à la position d'équilibre).
\end{itemize}

\textbf{Travail et moment d'une force (rotation)}
\begin{itemize}
\item Moment d'une force $\vec{F}$ par rapport à un axe $\Delta$ : $\mathcal{M} = F \times d$ ($d$ = bras de levier).
\item Si $\mathcal{M}$ constant : $W = \mathcal{M}\,\Delta\theta$ ($\Delta\theta$ en \textbf{radians}).
\item Relation géométrique treuil (roulement sans glissement) : $\Delta x = r\,\Delta\theta$ (longueur enroulée).
\end{itemize}

\textbf{Puissance}
\begin{itemize}
\item Puissance moyenne : $P = \dfrac{W}{\Delta t}$ (unité : watt \si{W} = \si{J.s^{-1}}).
\item Puissance instantanée : $P = \vec{F} \cdot \vec{v} = F v \cos(\vec{F}, \vec{v})$.
\item Puissance d'une force de moment constant : $P = \mathcal{M}\,\omega$ ($\omega$ en $\si{rad.s^{-1}}$).
\item Conversion : $\omega = \dfrac{2\pi N}{60}$ si $N$ en $\si{tr.min^{-1}}$ ; $1\ \text{tour} = 2\pi\ \si{rad}$.
\end{itemize}

\textbf{Conversions utiles}
\begin{itemize}
\item $r = \SI{10}{cm} = \SI{0,10}{m}$ ; $L = \SI{40}{cm} = \SI{0,40}{m}$.
\item Circonférence du tambour : $\mathcal{C} = 2\pi r$.
\end{itemize}
\end{aretenir}

\courssub{Démarche et corrigé commenté}

\begin{exoresolu}[Le treuil du chantier de Yaoundé]
Énoncé : $m = \SI{30}{kg}$ ; $h = \SI{8,0}{m}$ ; $\Delta t = \SI{50}{s}$ ; $r = \SI{10}{cm} = \SI{0,10}{m}$ ; $L = \SI{40}{cm} = \SI{0,40}{m}$ ; $g = \SI{9,8}{N.kg^{-1}}$ ; moteur : $N = \SI{300}{tr.min^{-1}}$, même couple que l'ouvrier. Répondre aux huit questions des consignes.
\tcblower
\textbf{1. Travail du poids.}\par
L'axe $Oz$ étant orienté vers le haut, $z_A = 0$ et $z_B = h$ :
\[ W_{AB}(\vec{P}) = mg(z_A - z_B) = -mgh = -30 \times 9,8 \times 8,0 \]
\[ \boxed{W_{AB}(\vec{P}) \approx \SI{-2,35e3}{J} = \SI{-2,35}{kJ}} \]
Le travail est négatif : le poids s'oppose à la montée (angle $(\vec{P}, \vec{AB}) = 180^\circ$), c'est un travail \textbf{résistant}. 
Vérification dimensionnelle : $[\,mgh\,] = \si{kg} \times \si{m.s^{-2}} \times \si{m} = \si{kg.m^2.s^{-2}} = \si{J}$. \checkmark

\textbf{2. Force conservative.}\par
Le travail du poids ne dépend que de l'altitude des points de départ et d'arrivée : $W_{AB}(\vec{P}) = mg(z_A - z_B)$, quelle que soit la trajectoire (verticale, rampe, escalier, etc.). Le long de la rampe inclinée, on trouvera donc \textbf{exactement la même valeur} $W_{AB}(\vec{P}) = \SI{-2,35}{kJ}$. C'est précisément la définition d'une \cle{force conservative} : le poids est conservative. 
\begin{attention}[Piège fréquent]
La rampe ne « supprime » pas le travail contre le poids ; elle permet seulement d'exercer une force plus faible (parallèle à la pente) sur une distance plus longue. Le travail total à fournir reste $mgh$.
\end{attention}

\textbf{3. Nombre de tours.}\par
À chaque tour de tambour, la corde s'enroule d'une longueur égale à la circonférence :
\[ \mathcal{C} = 2\pi r = 2\pi \times 0,10 \approx \SI{0,628}{m} \]
Nombre de tours nécessaires :
\[ n = \dfrac{h}{\mathcal{C}} = \dfrac{8,0}{0,628} \approx 12,7 \text{ tours} \]
La manivelle est solidaire du tambour : elle balaie le même angle. L'angle total est :
\[ \Delta\theta = n \times 2\pi = \dfrac{h}{r} = \dfrac{8,0}{0,10} = \boxed{\SI{80}{rad}} \]
(soit $12,7 \times 2\pi \approx 80\ \si{rad}$). Remarque : la relation $\Delta\theta = h/r$ traduit simplement le roulement sans glissement de la corde sur le tambour.

\textbf{4. Travail de l'ouvrier.}\par
La montée se fait à vitesse constante ($\Delta E_c = 0$) et le treuil est parfait (pas de pertes par frottement) : tout le travail fourni par Ahmadou sert à compenser le travail résistant du poids. Donc :
\[ W(\vec{F}) = -W_{AB}(\vec{P}) = mgh = 30 \times 9,8 \times 8,0 \]
\[ \boxed{W(\vec{F}) \approx \SI{2,35e3}{J} = \SI{2,35}{kJ}} \]
Ce travail est \textbf{moteur} (positif) : la force $\vec{F}$ favorise la rotation de la manivelle (angle $(\vec{F}, \text{déplacement poignée}) = 0^\circ$).

\textbf{5. Force sur la manivelle.}\par
La force $\vec{F}$ est constante et toujours perpendiculaire à la manivelle : son moment par rapport à l'axe est $\mathcal{M} = F \times L$, constant. La relation travail--moment s'applique :
\[ W(\vec{F}) = \mathcal{M}\,\Delta\theta = F \times L \times \Delta\theta \]
d'où :
\[ F = \dfrac{W(\vec{F})}{L\,\Delta\theta} = \dfrac{2352}{0,40 \times 80} = \boxed{F \approx \SI{73,5}{N}} \]
Comparaison : le poids du seau vaut $P = mg = 30 \times 9,8 = \SI{294}{N}$. La force à exercer est \textbf{4 fois plus faible} que le poids ! C'est l'\textbf{avantage mécanique} du treuil : $\dfrac{P}{F} = \dfrac{L}{r} = \dfrac{0,40}{0,10} = 4$. 
On peut vérifier autrement : à vitesse constante, les moments s'équilibrent sur l'axe, $F \times L = mg \times r$, d'où $F = mg\,\dfrac{r}{L} = 294 \times \dfrac{0,10}{0,40} = \SI{73,5}{N}$. \checkmark

\textbf{6. Puissance moyenne d'Ahmadou.}\par
\[ P = \dfrac{W(\vec{F})}{\Delta t} = \dfrac{2352}{50} \approx \boxed{\SI{47}{W}} \]
C'est une puissance modeste, tout à fait compatible avec un effort humain prolongé (un adulte en bonne santé peut développer environ 50 à 100 W en régime durable ; un athlète jusqu'à 300-400 W sur de courtes durées).

\textbf{7. Comparaison avec le moteur.}\par
Le moteur applique le même moment que l'ouvrier sur l'axe du tambour :
\[ \mathcal{M} = F \times L = 73,5 \times 0,40 = \SI{29,4}{N.m} \]
(ou directement $\mathcal{M} = mgr = 294 \times 0,10 = \SI{29,4}{N.m}$). 
Sa vitesse angulaire :
\[ \omega = \dfrac{2\pi N}{60} = \dfrac{2\pi \times 300}{60} = 10\pi \approx \SI{31,4}{rad.s^{-1}} \]
Puissance du moteur (force de moment constant) :
\[ P' = \mathcal{M}\,\omega = 29,4 \times 31,4 \approx \boxed{\SI{9,2e2}{W} \approx \SI{0,92}{kW}} \]
Vérification dimensionnelle : $[\,\mathcal{M}\omega\,] = \si{N.m} \times \si{rad.s^{-1}} = \si{kg.m^2.s^{-3}} = \si{W}$. \checkmark
Conclusion : $P' \approx 20\,P$. À \textbf{couple égal}, le moteur est environ \textbf{vingt fois plus puissant} que l'ouvrier, car il tourne beaucoup plus vite : la puissance en rotation dépend du produit $\mathcal{M}\,\omega$, pas du moment seul. Le même seau serait monté en $\Delta t' = \dfrac{W}{P'} \approx \dfrac{2352}{923} \approx \SI{2,5}{s}$ !

\textbf{8. Discussion : pourquoi la manivelle longue fatigue-t-elle moins ?}\par
Le travail à fournir est imposé par la charge : $W = mgh$, quelle que soit la machine. Or $W = F\,L\,\Delta\theta$ avec $\Delta\theta = h/r$ fixé par le tambour. Pour un même travail, $F$ et $L$ sont donc \textbf{inversement proportionnels} :
\[ F = \dfrac{mgh}{L\,\Delta\theta} = mg\,\dfrac{r}{L} \]
Allonger la manivelle (grand bras de levier $L$) réduit la force à exercer : l'ouvrier se fatigue moins à chaque instant (contrainte musculaire moindre). Le « prix à payer » : la poignée parcourt un cercle plus grand ($2\pi L$ par tour), donc l'ouvrier déplace sa main sur une \textbf{distance plus longue} (et le nombre de tours reste le même). C'est le principe général des machines simples (treuil, levier, poulies, plan incliné) : \textbf{on ne gagne jamais en travail, on échange force contre distance} (ou temps contre puissance).
\end{exoresolu}

\courssub{Exemple d'application : Travail de la tension d'un ressort}

\begin{exemplebox}[Bloc sur ressort horizontal — Force conservative]
Un bloc de masse $m = \SI{0,50}{kg}$ repose sur un plan horizontal sans frottement, attaché à un ressort de raideur $k = \SI{200}{N.m^{-1}}$ (position d'équilibre $x=0$). On étire le ressort de $x_A = 0$ à $x_B = \SI{10}{cm} = \SI{0,10}{m}$ en tirant doucement sur le bloc (vitesse quasi nulle).
\begin{enumerate}
\item Calculer le travail $W_{AB}(\vec{F}_{\text{rappel}})$ de la force de rappel du ressort.
\item Calculer le travail $W_{AB}(\vec{F}_{\text{ext}})$ de la force extérieure qui étire le ressort.
\item Que vaut le travail de la force de rappel si on ramène le bloc de $x_B$ à $x_A$ ? Et sur un aller-retour $A \to B \to A$ ?
\end{enumerate}
\tcblower
\textbf{Solution.}
\begin{enumerate}
\item La force de rappel est $\vec{F} = -k x\,\vec{u}_x$. Son travail entre $x_A$ et $x_B$ :
\[ W_{AB}(\vec{F}) = \int_{x_A}^{x_B} (-k x)\,dx = \left[ -\frac{1}{2}k x^2 \right]_{0}^{0,10} = -\frac{1}{2} \times 200 \times (0,10)^2 = \boxed{\SI{-1,0}{J}} \]
(Travail résistant : la force s'oppose à l'étirement).
\item Comme la vitesse est quasi nulle, $\Delta E_c \approx 0$. Le théorème de l'énergie cinétique donne $W_{AB}(\vec{F}_{\text{ext}}) + W_{AB}(\vec{F}) = 0$, d'où $W_{AB}(\vec{F}_{\text{ext}}) = -W_{AB}(\vec{F}) = \boxed{\SI{+1,0}{J}}$. C'est l'énergie potentielle élastique stockée : $E_p = \frac{1}{2}k x^2$.
\item De $x_B$ vers $x_A$ : $W_{BA}(\vec{F}) = \frac{1}{2}k(x_B^2 - x_A^2) = +\SI{1,0}{J}$ (travail moteur). Sur l'aller-retour $A \to B \to A$ : $W_{\text{cycle}} = W_{AB} + W_{BA} = \SI{-1,0}{J} + \SI{+1,0}{J} = \boxed{\SI{0}{J}}$. La force de rappel du ressort est bien \textbf{conservative}.
\end{enumerate}
\end{exemplebox}

\courssub{Attention : Pièges fréquents dans ce type de problème}

\begin{attention}[Pièges à éviter absolument]
\begin{itemize}
\item \textbf{Unités :} Oublier de convertir les longueurs en mètres ($r = \SI{0,10}{m}$, $L = \SI{0,40}{m}$) $\rightarrow$ facteur 10 ou 100 d'erreur sur $F$ et $\mathcal{M}$ !
\item \textbf{Angle en radians :} Utiliser $\Delta\theta$ en degrés dans $W = \mathcal{M}\,\Delta\theta$ ou $P = \mathcal{M}\,\omega$. L'angle \textbf{doit} être en radians (adimensionnel).
\item \textbf{Vitesse angulaire :} Confondre $N$ (en $\si{tr.min^{-1}}$) et $\omega$ (en $\si{rad.s^{-1}}$). Ne jamais oublier $\omega = \dfrac{2\pi N}{60}$.
\item \textbf{Signe du travail :} Ne pas confondre travail moteur ($W>0$, force dans le sens du déplacement) et travail résistant ($W<0$, force opposée). Pour le poids à la montée : $W<0$.
\item \textbf{Machine parfaite :} Croire que le treuil « réduit le travail ». Il réduit la \textbf{force} (avantage mécanique), jamais le travail (principe de l'énergie : $W_{\text{entrée}} = W_{\text{sortie}}$ si rendement 100 \%). 
\item \textbf{Ressort :} Ne pas utiliser $W = F \times \Delta x$ avec $F$ constante (erreur !). La force du ressort varie : il faut intégrer ou utiliser $W = \frac{1}{2}k(x_A^2 - x_B^2)$.
\item \textbf{Chiffres significatifs :} Les données ($m=30$, $h=8,0$, $g=9,8$) imposent 2 chiffres significatifs pour les résultats finaux ($W \approx 2,4\ \si{kJ}$, $F \approx 74\ \si{N}$, $P \approx 47\ \si{W}$, $P' \approx 920\ \si{W}$). Conserver 3 chiffres ($2,35$) pour les calculs intermédiaires.
\end{itemize}
\end{attention}

\courssub{Bilan de l'activité}

Cette activité t'a permis de mobiliser l'ensemble des savoir-faire de la leçon dans une situation unique et concrète :
\begin{itemize}
\item le \textbf{travail du poids} $W_{AB}(\vec{P}) = mg(z_A - z_B)$, négatif à la montée, et l'idée essentielle que le poids est une \textbf{force conservative} : son travail ne dépend que de la dénivellation, pas du chemin (vertical ou rampe) ;
\item le travail de la \textbf{tension d'un ressort} $W = \frac{1}{2}k(x_A^2 - x_B^2)$, lui aussi conservative (énergie potentielle élastique) ;
\item le lien géométrique entre translation de la charge et rotation du tambour : $h = r\,\Delta\theta$ (roulement sans glissement) ;
\item la relation \textbf{travail--moment} $W = \mathcal{M}\,\Delta\theta$, qui a permis de remonter du travail à la force sur la manivelle et de mettre en évidence l'avantage mécanique $\dfrac{L}{r}$ du treuil ;
\item la \textbf{puissance moyenne} $P = \dfrac{W}{\Delta t}$ et la \textbf{puissance en rotation} $P = \mathcal{M}\,\omega$, dont la comparaison montre qu'à couple égal, c'est la vitesse angulaire qui fait la différence de puissance entre l'homme et le moteur ;
\item enfin, le principe général des \textbf{machines simples} : elles ne créent pas de travail, elles le « répartissent » autrement — moins de force, mais plus de distance (ou plus de temps pour une même puissance).
\end{itemize}
Tu es maintenant prêt à analyser n'importe quel dispositif de levage ou de manutention — treuil de chantier, poulie de puits, cric de mototaxi, vérin hydraulique — avec les outils rigoureux du travail et de la puissance.

\begin{savaistu}[Le watt, James Watt et la motorisation au Cameroun]
L'unité de puissance, le \textbf{watt} (symbole \si{W}), rend hommage à l'ingénieur écossais James Watt (1736--1819), qui a perfectionné la machine à vapeur. Il a défini le « cheval-vapeur » (\si{ch} ou \si{hp}) pour vendre ses machines aux mineurs : $1\ \si{ch} = 735,5\ \si{W}$ (puissance d'un cheval tirant une charge). Aujourd'hui, au Cameroun, la puissance des groupes électrogènes (si fréquents face aux délestages d'ENEO) s'exprime en \si{kVA} (kilovoltampère, puissance apparente) ou en \si{kW} (puissance active). Un groupe de \SI{10}{kVA} peut alimenter une maison (éclairage, TV, frigo, ventilateurs). Les moteurs des pompes à eau (forages) ou des broyeuses agricoles sont dimensionnés en \si{kW} : un moteur de \SI{2,2}{kW} (3 ch) tourne souvent à 1500 ou 3000 \si{tr.min^{-1}} (50 Hz). La relation $P = \mathcal{M}\,\omega$ est au cœur du dimensionnement : pour une puissance donnée, un moteur rapide (grande $\omega$) a un couple faible (arbre fin), un moteur lent (petite $\omega$) a un couple fort (arbre gros, réducteur nécessaire). C'est exactement ce que tu as comparé entre Ahmadou ($\omega \approx 1,6\ \si{rad.s^{-1}}$) et le moteur ($\omega \approx 31\ \si{rad.s^{-1}}$) !
\end{savaistu}

\courssub{Exercices d'entraînement}

\begin{exercice}[Treuil de forage — Application directe]
Un forage villageois utilise un treuil manuel pour remonter un seau d'eau de masse $m = \SI{15}{kg}$ depuis une profondeur $h = \SI{20}{m}$. Le tambour a un rayon $r = \SI{12}{cm}$, la manivelle une longueur $L = \SI{50}{cm}$. L'opérateur tourne à une vitesse constante de $12\ \si{tr.min^{-1}}$.
\begin{enumerate}
\item Calculer le travail du poids et le travail fourni par l'opérateur (treuil parfait).
\item Déterminer la force tangente que l'opérateur doit exercer sur la manivelle.
\item Calculer la puissance moyenne développée par l'opérateur.
\item Si l'on remplace la manivelle par une de $L' = \SI{30}{cm}$, que deviennent la force et la puissance ? (La vitesse de rotation en \si{tr.min^{-1}} reste la même).
\end{enumerate}
\end{exercice}

\begin{exercice}[Ressort vertical — Énergie potentielle]
Une masse $m = \SI{200}{g}$ est suspendue à un ressort vertical de raideur $k = \SI{50}{N.m^{-1}}$. On tire doucement la masse vers le bas pour l'écarter de $\SI{5,0}{cm}$ de sa position d'équilibre, puis on la lâche sans vitesse initiale.
\begin{enumerate}
\item Calculer le travail de la force de rappel du ressort pendant cette extension (de l'équilibre à $x = -\SI{5,0}{cm}$).
\item Calculer le travail du poids pendant ce même déplacement.
\item En déduire, par le théorème de l'énergie cinétique, la vitesse de la masse lorsqu'elle repasse par la position d'équilibre.
\end{enumerate}
\end{exercice}

\begin{corrige}[Exercice 1 : Treuil de forage]
\begin{enumerate}
\item $W(P) = -mgh = -15 \times 9,8 \times 20 = \SI{-2940}{J} \approx \SI{-2,9}{kJ}$. $W(\text{opérateur}) = +\SI{2,9}{kJ}$.
\item $\Delta\theta = h/r = 20 / 0,12 \approx 166,7\ \si{rad}$. $\mathcal{M} = W/\Delta\theta = 2940 / 166,7 \approx \SI{17,6}{N.m}$. $F = \mathcal{M}/L = 17,6 / 0,50 = \SI{35,3}{N}$.
\item $\omega = 2\pi \times 12 / 60 = 0,4\pi \approx \SI{1,26}{rad.s^{-1}}$. $P = \mathcal{M}\omega = 17,6 \times 1,26 \approx \SI{22}{W}$. (Ou $P = W/\Delta t$, $\Delta t = h / (v_{\text{corde}}) = 20 / (0,12 \times 1,26) \approx 132\ \si{s}$, $P = 2940/132 \approx 22\ \si{W}$).
\item $L' = 0,30\ \si{m}$. $F' = \mathcal{M}/L' = 17,6 / 0,30 \approx \SI{59}{N}$ (force augmente). $\omega$ inchangé $\rightarrow$ $P$ inchangée ($\approx \SI{22}{W}$). L'opérateur fatigue plus (force plus grande) pour la même puissance.
\end{enumerate}
\end{corrige}

\begin{corrige}[Exercice 2 : Ressort vertical]
\begin{enumerate}
\item $W_{\text{ressort}} = \frac{1}{2}k(0^2 - x^2) = -\frac{1}{2} \times 50 \times (0,05)^2 = \SI{-0,0625}{J}$.
\item $W_{\text{poids}} = mg\Delta z = 0,20 \times 9,8 \times 0,05 = \SI{+0,098}{J}$ (poids moteur, même sens que déplacement vers le bas).
\item Au passage par l'équilibre ($x=0$), $E_c = \frac{1}{2}mv^2 = W_{\text{ressort}} + W_{\text{poids}} = -0,0625 + 0,098 = \SI{0,0355}{J}$. $v = \sqrt{2 \times 0,0355 / 0,20} \approx \SI{0,60}{m.s^{-1}}$.
\end{enumerate}
\end{corrige}

\newpage

\coursec{Méthodes et exercices résolus}

\coursec{Travail et puissance d'une force}

\courssub{Travail d'une force constante}

\begin{definition}[Travail d'une force constante]
Soit une force \cle{constante} $\vec{F}$ (norme, direction, sens invariables) dont le point d'application se déplace du point $A$ au point $B$. Le \cle{travail} de cette force au cours de ce déplacement est le scalaire :
\[ W_{AB}(\vec{F}) = \vec{F}\cdot\overrightarrow{AB} = F \times AB \times \cos\alpha \]
où $\alpha = (\vec{F}, \overrightarrow{AB})$ est l'angle orienté entre la force et le déplacement.
L'unité SI du travail est le \cle{joule} (symbole \textbf{J}) : $\SI{1}{J} = \SI{1}{N.m}$.
\end{definition}

\begin{propriete}[Nature du travail]
\begin{itemize}
\item $0 \leq \alpha < 90°$ $\Rightarrow$ $\cos\alpha > 0$ $\Rightarrow$ $W > 0$ : travail \cle{moteur} (la force favorise le déplacement).
\item $\alpha = 90°$ $\Rightarrow$ $\cos\alpha = 0$ $\Rightarrow$ $W = 0$ : travail \cle{nul} (force perpendiculaire au déplacement).
\item $90° < \alpha \leq 180°$ $\Rightarrow$ $\cos\alpha < 0$ $\Rightarrow$ $W < 0$ : travail \cle{résistant} (la force s'oppose au déplacement).
\end{itemize}
Le travail ne dépend que des points $A$ et $B$, pas de la trajectoire suivie (pour une force constante).
\end{propriete}

\begin{methode}[Calculer le travail d'une force constante]
\begin{enumerate}
\item \textbf{Vérifier les conditions} : force constante (norme, direction, sens fixes). Le déplacement peut être curviligne.
\item \textbf{Schéma} : représenter $\vec{F}$, $\overrightarrow{AB}$, repérer l'angle $\alpha = (\vec{F}, \overrightarrow{AB})$.
\item \textbf{Formule} : $W_{AB}(\vec{F}) = F \times AB \times \cos\alpha$ ($F$ en N, $AB$ en m, $W$ en J).
\item \textbf{Convertir les unités} : distances en \textbf{mètres}, angles en \textbf{degrés} (pour $\cos$) ou radians.
\item \textbf{Conclure} : valeur numérique, unité (J), nature (moteur/nul/résistant).
\end{enumerate}
\end{methode}

\begin{exoresolu}[Le porteur de marchandises à la gare de Yaoundé]
Un porteur tire une caisse de masse $m = \SI{40}{kg}$ sur un quai horizontal, sur une distance $AB = \SI{25}{m}$, au moyen d'une corde faisant un angle $\alpha = 30°$ avec l'horizontale. La tension de la corde est constante : $T = \SI{120}{N}$. Les frottements du sol sont modélisés par une force constante $\vec{f}$ horizontale, de norme $f = \SI{60}{N}$, opposée au déplacement.
\begin{enumerate}
\item Calculer le travail de la tension $\vec{T}$ sur le trajet $AB$. Préciser sa nature.
\item Calculer le travail de la force de frottement $\vec{f}$. Préciser sa nature.
\item Calculer le travail du poids et celui de la réaction normale du sol.
\end{enumerate}
\tcblower
\textbf{1. Travail de la tension.}

La force $\vec{T}$ est constante et le déplacement est rectiligne. L'angle entre $\vec{T}$ et $\overrightarrow{AB}$ est $\alpha = 30°$.
\begin{align*}
W_{AB}(\vec{T}) &= T \times AB \times \cos\alpha \\
&= 120 \times 25 \times \cos 30° \\
&= 3000 \times \dfrac{\sqrt{3}}{2} = 1500\sqrt{3} \approx \SI{2598}{J} \approx \SI{2.6e3}{J}
\end{align*}
Comme $\alpha < 90°$, $\cos\alpha > 0$ : le travail est \textbf{moteur}. C'est cohérent : la tension aide la caisse à avancer.

\textbf{2. Travail des frottements.}

$\vec{f}$ est constante, horizontale et opposée au déplacement : l'angle entre $\vec{f}$ et $\overrightarrow{AB}$ vaut $180°$, donc $\cos 180° = -1$.
\[ W_{AB}(\vec{f}) = f \times AB \times \cos 180° = 60 \times 25 \times (-1) = \SI{-1500}{J} = \SI{-1.5e3}{J} \]
Le travail est \textbf{résistant} : les frottements s'opposent au mouvement.

\textbf{3. Travail du poids et de la réaction normale.}

Le poids $\vec{P}$ et la réaction normale $\vec{R}_N$ sont tous deux verticaux, alors que le déplacement $\overrightarrow{AB}$ est horizontal : l'angle vaut $90°$ dans les deux cas, et $\cos 90° = 0$.
\[ W_{AB}(\vec{P}) = 0 \quad ; \quad W_{AB}(\vec{R}_N) = 0 \]
\textbf{Conclusion} : $W(\vec{T}) \approx \SI{2.6e3}{J}$ (moteur), $W(\vec{f}) = \SI{-1.5e3}{J}$ (résistant), et les forces perpendiculaires au déplacement ne travaillent pas.
\end{exoresolu}

\courssub{Travail du poids et forces conservatives}

\begin{definition}[Force conservative]
Une force est dite \cle{conservative} si son travail entre deux points $A$ et $B$ ne dépend \emph{que} des positions de $A$ et $B$, et non du chemin suivi pour aller de $A$ à $B$.
Le poids $\vec{P} = m\vec{g}$ est une force conservative (dans le champ de gravité uniforme).
\end{definition}

\begin{propriete}[Travail du poids]
Dans un référentiel terrestre, en choisissant un axe vertical $Oz$ orienté vers le haut ($z$ = altitude), le travail du poids $\vec{P} = m\vec{g}$ entre $A(z_A)$ et $B(z_B)$ vaut :
\[ W_{AB}(\vec{P}) = mg\,(z_A - z_B) \]
\begin{itemize}
\item Si le corps \emph{descend} ($z_A > z_B$) : $W > 0$ (travail moteur).
\item Si le corps \emph{monte} ($z_A < z_B$) : $W < 0$ (travail résistant).
\item Si $z_A = z_B$ : $W = 0$.
\end{itemize}
\emph{Attention} : on n'utilise \textbf{jamais} la longueur du trajet $L$, seule la dénivellation $h = |z_A - z_B|$ intervient.
\end{propriete}

\begin{methode}[Utiliser le fait que le poids est une force conservative]
\begin{enumerate}
\item \textbf{Repérer un axe vertical} $Oz$ orienté vers le haut, relever les altitudes $z_A$ et $z_B$.
\item \textbf{Écrire} : $W_{AB}(\vec{P}) = mg\,(z_A - z_B)$.
\item \textbf{Analyser le signe} : descente $\Rightarrow$ moteur ($+$) ; montée $\Rightarrow$ résistant ($-$).
\item \textbf{Ne jamais utiliser la longueur du trajet} $L$.
\end{enumerate}
\end{methode}

\begin{exoresolu}[Descente du mont Mbappé]
Un randonneur de masse $m = \SI{70}{kg}$ (équipement compris) descend du sommet d'une colline d'altitude $z_A = \SI{850}{m}$ jusqu'à un village d'altitude $z_B = \SI{400}{m}$, en suivant un sentier sinueux de longueur $L = \SI{6.0}{km}$. On prend $g = \SI{9.8}{N.kg^{-1}}$.
\begin{enumerate}
\item Calculer le travail du poids du randonneur entre $A$ et $B$.
\item Ce travail dépend-il de la longueur du sentier ? Justifier.
\item Que vaudrait ce travail si le randonneur remontait du village au sommet ?
\end{enumerate}
\tcblower
\textbf{1. Travail du poids.}

Le poids est une force conservative : son travail ne dépend que des altitudes de départ et d'arrivée :
\begin{align*}
W_{AB}(\vec{P}) &= mg\,(z_A - z_B) \\
&= 70 \times 9.8 \times (850 - 400) \\
&= 686 \times 450 = \SI{308700}{J} \approx \SI{3.1e5}{J}
\end{align*}
Le randonneur descend ($z_A > z_B$) : le travail est \textbf{moteur}, ce qui est cohérent.

\textbf{2. Indépendance vis-à-vis du chemin.}

Non, ce travail ne dépend pas de la longueur $L = \SI{6.0}{km}$ du sentier. Le poids étant une force \textbf{conservative}, son travail ne dépend que des positions initiale et finale (ici des altitudes $z_A$ et $z_B$), et non du chemin suivi. Un sentier plus long ou plus direct donnerait exactement le même travail.

\textbf{3. Remontée de $B$ vers $A$.}

\[ W_{BA}(\vec{P}) = mg\,(z_B - z_A) = 70 \times 9.8 \times (400 - 850) = \SI{-3.1e5}{J} \]
Le travail est \textbf{résistant} : lors de la montée, le poids s'oppose au déplacement. On remarque que $W_{BA}(\vec{P}) = -W_{AB}(\vec{P})$.

\textbf{Conclusion} : $W_{AB}(\vec{P}) \approx \SI{+3.1e5}{J}$ (moteur à la descente), indépendant du chemin ; $W_{BA}(\vec{P}) \approx \SI{-3.1e5}{J}$ (résistant à la montée).
\end{exoresolu}

\courssub{Travail de la tension d'un ressort}

\begin{definition}[Ressort à réponse linéaire -- Loi de Hooke]
Un ressort de raideur $k$ (en $\si{N.m^{-1}}$), au repos pour $x=0$ (position à vide), exerce une force de rappel (tension) $\vec{T}$ telle que :
\[ T = k\,|x| \quad \text{et} \quad \vec{T} \text{ est dirigée vers la position d'équilibre.} \]
La tension n'est \textbf{pas constante} lors d'un étirement ou d'un relâchement.
\end{definition}

\begin{propriete}[Travail de la tension d'un ressort]
Le travail de la tension $\vec{T}$ lorsque l'allongement passe de $x_1$ à $x_2$ (algorithme : $x$ algébrique, $x>0$ étirement, $x<0$ compression) est :
\[ W_{1\to 2}(\vec{T}) = \dfrac{1}{2}\,k\,\left(x_1^{2} - x_2^{2}\right) \]
La tension d'un ressort est une force \cle{conservative}.
\begin{itemize}
\item Si $|x_2| > |x_1|$ (le ressort s'étire ou se compresse davantage) : $W < 0$ (résistant).
\item Si $|x_2| < |x_1|$ (le ressort se détend) : $W > 0$ (moteur).
\end{itemize}
\end{propriete}

\begin{methode}[Calculer le travail de la tension d'un ressort à réponse linéaire]
\begin{enumerate}
\item \textbf{Données} : raideur $k$ ($\si{N.m^{-1}}$), allongements $x_1$ et $x_2$ (en \textbf{mètres}).
\item \textbf{Ne pas utiliser} $W = F\cdot AB\cdot\cos\alpha$ (force non constante).
\item \textbf{Appliquer} : $W_{1\to 2}(\vec{T}) = \dfrac{1}{2}k(x_1^2 - x_2^2)$.
\item \textbf{Interpréter le signe} : étirement $\Rightarrow$ résistant ($-$) ; détente $\Rightarrow$ moteur ($+$).
\item \textbf{Force extérieure} : si une force $\vec{F}_{\text{ext}}$ étire le ressort \emph{lentement} (quasi-statique), $\vec{F}_{\text{ext}} = -\vec{T}$, donc $W(\vec{F}_{\text{ext}}) = -W(\vec{T}) = \dfrac{1}{2}k(x_2^2 - x_1^2)$.
\end{enumerate}
\end{methode}

\begin{exoresolu}[Le dynamomètre de l'atelier de physique]
Un ressort à spires non jointives, de raideur $k = \SI{200}{N.m^{-1}}$, est fixé par une extrémité à un support. Un opérateur tire lentement sur l'extrémité libre et fait passer son allongement de $x_1 = \SI{2.0}{cm}$ à $x_2 = \SI{8.0}{cm}$.
\begin{enumerate}
\item Calculer le travail de la tension du ressort au cours de cette opération. Commenter son signe.
\item En déduire le travail de la force exercée par l'opérateur, supposée à chaque instant opposée à la tension.
\end{enumerate}
\tcblower
\textbf{1. Travail de la tension.}

Conversion des allongements en mètres : $x_1 = \SI{0.020}{m}$ et $x_2 = \SI{0.080}{m}$.

La tension du ressort n'est pas constante pendant l'étirement : on applique la formule du travail de la tension d'un ressort à réponse linéaire :
\begin{align*}
W_{1\to 2}(\vec{T}) &= \dfrac{1}{2}\,k\,\left(x_1^{2} - x_2^{2}\right) \\
&= \dfrac{1}{2} \times 200 \times \left(0.020^2 - 0.080^2\right) \\
&= 100 \times \left(4.0\times 10^{-4} - 6.4\times 10^{-3}\right) \\
&= 100 \times (-6.0\times 10^{-3}) = \SI{-0.60}{J}
\end{align*}
Le travail est \textbf{négatif} : c'est normal, car le ressort s'étire et sa tension (force de rappel dirigée vers sa position d'équilibre) s'oppose à l'allongement. Le travail de la tension est résistant.

\textbf{2. Travail de la force de l'opérateur.}

L'opérateur tire \emph{lentement} : à chaque instant, sa force $\vec{F}$ est opposée à la tension $\vec{T}$, donc $\vec{F} = -\vec{T}$. Les travaux sont alors opposés :
\[ W_{1\to 2}(\vec{F}) = -W_{1\to 2}(\vec{T}) = \dfrac{1}{2}\,k\,\left(x_2^{2} - x_1^{2}\right) = \SI{+0.60}{J} \]
\textbf{Conclusion} : $W(\vec{T}) = \SI{-0.60}{J}$ (résistant) et $W(\vec{F}) = \SI{+0.60}{J}$ (moteur) : l'opérateur a fourni \SI{0.60}{J} pour étirer le ressort.
\end{exoresolu}

\courssub{Travail et moment d'une force (rotation)}

\begin{definition}[Moment d'une force par rapport à un axe]
Pour une force $\vec{F}$ appliquée à un solide en rotation autour d'un axe fixe $(\Delta)$, le \cle{moment} de $\vec{F}$ par rapport à $(\Delta)$ est :
\[ \mathcal{M}_{\Delta}(\vec{F}) = \pm F \times d \]
où $d$ est le bras de levier (distance de la ligne d'action de $\vec{F}$ à l'axe) et le signe $\pm$ dépend du sens de rotation choisi comme positif (généralement trigonométrique). Unité : $\si{N.m}$.
\end{definition}

\begin{propriete}[Travail d'une force de moment constant en rotation]
Si le moment $\mathcal{M}_{\Delta}(\vec{F})$ est constant pendant une rotation d'angle $\theta$ (en \textbf{radians}), le travail de la force vaut :
\[ W(\vec{F}) = \mathcal{M}_{\Delta}(\vec{F}) \times \theta \]
\emph{Analyse dimensionnelle} : $\si{N.m} \times \si{rad} = \si{J}$ (le radian est sans dimension).
Le travail est moteur si le moment et la rotation ont même signe, résistant sinon.
\end{propriete}

\begin{methode}[Travail d'une force appliquée à un solide en rotation]
\begin{enumerate}
\item \textbf{Calculer le moment} $\mathcal{M}_{\Delta}(\vec{F}) = \pm F \times d$ (en $\si{N.m}$).
\item \textbf{Convertir l'angle de rotation en radians} : $\theta\,(\text{rad}) = \theta\,(\text{deg}) \times \dfrac{\pi}{180}$ ; $1$ tour $= 2\pi$ rad.
\item \textbf{Appliquer} : $W(\vec{F}) = \mathcal{M}_{\Delta}(\vec{F}) \times \theta$.
\item \textbf{Vérifier les signes} : moment et rotation même sens $\Rightarrow$ $W>0$ (moteur).
\end{enumerate}
\end{methode}

\begin{exoresolu}[Le treuil du chantier de Nachtigal]
Sur le chantier du barrage de Nachtigal, un ouvrier actionne un treuil pour soulever une charge. Il exerce sur la manivelle une force $\vec{F}$ constante, de norme $F = \SI{150}{N}$, perpendiculaire à la manivelle de longueur $d = \SI{40}{cm}$. La manivelle effectue $N = 15$ tours.
\begin{enumerate}
\item Calculer le moment de $\vec{F}$ par rapport à l'axe du treuil.
\item Calculer le travail de $\vec{F}$ au cours des 15 tours.
\end{enumerate}
\tcblower
\textbf{1. Moment de la force.}

La force est perpendiculaire à la manivelle : le bras de levier est exactement $d = \SI{0.40}{m}$. La force fait tourner la manivelle : son moment est moteur (on le compte positivement) :
\[ \mathcal{M}_{\Delta}(\vec{F}) = F \times d = 150 \times 0.40 = \SI{60}{N.m} \]

\textbf{2. Travail au cours des 15 tours.}

L'angle de rotation total, en radians :
\[ \theta = N \times 2\pi = 15 \times 2\pi = 30\pi \approx \SI{94.2}{rad} \]
Le moment est constant pendant toute la rotation, donc :
\begin{align*}
W(\vec{F}) &= \mathcal{M}_{\Delta}(\vec{F}) \times \theta \\
&= 60 \times 30\pi = 1800\,\pi \approx \SI{5655}{J} \approx \SI{5.7e3}{J}
\end{align*}
Le moment et la rotation sont dans le même sens : le travail est \textbf{moteur}, ce qui est cohérent avec le fait que l'ouvrier fournit un effort pour soulever la charge.

\textbf{Conclusion} : $\mathcal{M}_{\Delta}(\vec{F}) = \SI{60}{N.m}$ et $W(\vec{F}) \approx \SI{5.7e3}{J}$ (travail moteur).
\end{exoresolu}

\courssub{Puissance d'une force}

\begin{definition}[Puissance d'une force]
La \cle{puissance} d'une force est la quantité de travail fournie par unité de temps.
\begin{itemize}
\item \textbf{Puissance moyenne} sur une durée $\Delta t$ : $\mathcal{P}_{\text{moy}} = \dfrac{W}{\Delta t}$.
\item \textbf{Puissance instantanée} pour une force constante dont le point d'application a une vitesse $\vec{v}$ : $\mathcal{P} = \vec{F}\cdot\vec{v} = F\,v\,\cos\alpha$ ($\alpha = (\vec{F}, \vec{v})$).
\end{itemize}
Unité SI : le \cle{watt} (symbole \textbf{W}) : $\SI{1}{W} = \SI{1}{J.s^{-1}}$.
Multiples usuels : $\SI{1}{kW} = \SI{e3}{W}$, $\SI{1}{MW} = \SI{e6}{W}$, $1\,\text{ch} = \SI{736}{W}$.
\emph{Attention} : le kilowattheure ($\si{kWh}$) est une unité d'\textbf{énergie} ($\SI{1}{kWh} = \SI{3.6e6}{J}$), pas de puissance.
\end{definition}

\begin{methode}[Calculer une puissance moyenne ou instantanée]
\begin{enumerate}
\item \textbf{Puissance moyenne} : calculer $W$ (en J), $\Delta t$ (en s), puis $\mathcal{P}_{\text{moy}} = W/\Delta t$ (en W).
\item \textbf{Convertir les durées} en secondes (min $\times 60$, h $\times 3600$).
\item \textbf{Puissance instantanée} : $\mathcal{P} = \vec{F}\cdot\vec{v} = F\,v\,\cos\alpha$. Si $\vec{F}\parallel\vec{v}$ : $\mathcal{P} = \pm F\,v$.
\item \textbf{Ne pas confondre} puissance (W) et énergie (J ou kWh).
\end{enumerate}
\end{methode}

\begin{exoresolu}[La pompe du forage villageois]
Une pompe électrique remonte l'eau d'un forage de profondeur $h = \SI{12}{m}$ dans un village de la région de l'Extrême-Nord. Elle refoule un volume $V = \SI{3000}{L}$ d'eau en $\Delta t = \SI{25}{min}$. On donne : masse volumique de l'eau $\rho = \SI{1.0e3}{kg.m^{-3}}$, $g = \SI{9.8}{N.kg^{-1}}$.
\begin{enumerate}
\item Calculer le travail minimal fourni par la pompe pour remonter cette eau.
\item En déduire la puissance moyenne de la pompe en watts, puis en kilowatts.
\end{enumerate}
\tcblower
\textbf{1. Travail minimal de la pompe.}

Masse d'eau pompée :
\[ m = \rho \times V = 1.0\times 10^{3} \times 3.000 = \SI{3000}{kg} \quad (V = \SI{3000}{L} = \SI{3.0}{m^3}) \]
Pour remonter l'eau de la hauteur $h$, la force de pompage doit au minimum compenser le poids ; le travail minimal est l'opposé du travail du poids :
\begin{align*}
W &= mgh \\
&= 3000 \times 9.8 \times 12 \\
&= \SI{352800}{J} \approx \SI{3.5e5}{J}
\end{align*}

\textbf{2. Puissance moyenne.}

Conversion de la durée : $\Delta t = 25 \times 60 = \SI{1500}{s}$.
\begin{align*}
\mathcal{P}_{\text{moy}} &= \dfrac{W}{\Delta t} = \dfrac{352800}{1500} \approx \SI{235}{W} \approx \SI{0.24}{kW}
\end{align*}
\textbf{Conclusion} : la pompe fournit un travail minimal $W \approx \SI{3.5e5}{J}$ et développe une puissance moyenne $\mathcal{P}_{\text{moy}} \approx \SI{235}{W}$, soit environ $\SI{0.24}{kW}$. En pratique, à cause des frottements et du rendement, la puissance réelle installée est supérieure.
\end{exoresolu}

\courssub{Puissance d'une force de moment constant}

\begin{propriete}[Puissance en rotation -- Moment constant]
Pour une force de moment constant $\mathcal{M}_{\Delta}(\vec{F})$ appliquée à un solide tournant à la vitesse angulaire $\omega$ (en $\si{rad.s^{-1}}$), la puissance instantanée développée est :
\[ \mathcal{P} = \mathcal{M}_{\Delta}(\vec{F}) \times \omega \]
Si la fréquence de rotation $n$ est donnée en tours par minute (tr/min) :
\[ \omega = \dfrac{2\pi\,n}{60} \]
\emph{Cohérence} : comme $W = \mathcal{M}_{\Delta}\,\theta$ et $\theta = \omega\,\Delta t$, on retrouve $\mathcal{P} = \dfrac{W}{\Delta t} = \mathcal{M}_{\Delta}\,\omega$.
Le signe de $\mathcal{P}$ indique si la puissance est motrice ($+$) ou résistante ($-$).
\end{propriete}

\begin{methode}[Puissance d'une force appliquée à un solide en rotation]
\begin{enumerate}
\item \textbf{Identifier} $\mathcal{M}_{\Delta}(\vec{F})$ (en $\si{N.m}$) et $\omega$ (en $\si{rad.s^{-1}}$).
\item \textbf{Convertir $n$ (tr/min) en $\omega$} : $\omega = \dfrac{2\pi n}{60}$.
\item \textbf{Appliquer} : $\mathcal{P} = \mathcal{M}_{\Delta} \times \omega$.
\item \textbf{Travail sur $\Delta t$} : $W = \mathcal{P} \times \Delta t = \mathcal{M}_{\Delta} \times \omega \times \Delta t$.
\end{enumerate}
\end{methode}

\begin{exoresolu}[Le moteur du moulin à maïs]
Le moteur électrique d'un moulin à maïs tourne à la fréquence $n = \SI{1500}{tr.min^{-1}}$. Le couple moteur, c'est-à-dire le moment de l'ensemble des forces exercées par le moteur sur l'arbre, est constant et vaut $\mathcal{M} = \SI{12}{N.m}$.
\begin{enumerate}
\item Convertir la vitesse de rotation en $\si{rad.s^{-1}}$.
\item Calculer la puissance développée par le moteur.
\item Calculer le travail fourni par le moteur pendant $\Delta t = \SI{10}{min}$ de fonctionnement.
\end{enumerate}
\tcblower
\textbf{1. Conversion de la vitesse angulaire.}

Un tour correspond à $2\pi$ rad ; en une minute le moteur effectue 1500 tours, donc :
\[ \omega = \dfrac{2\pi\,n}{60} = \dfrac{2\pi \times 1500}{60} = 50\pi \approx \SI{157}{rad.s^{-1}} \]

\textbf{2. Puissance du moteur.}

Le moment est constant pendant la rotation : on applique la relation puissance--moment :
\[ \mathcal{P} = \mathcal{M} \times \omega = 12 \times 50\pi = 600\,\pi \approx \SI{1885}{W} \approx \SI{1.9}{kW} \]
Moment et rotation sont dans le même sens : la puissance est \textbf{motrice}.

\textbf{3. Travail pendant 10 minutes.}

Conversion : $\Delta t = 10 \times 60 = \SI{600}{s}$. D'après la relation $\mathcal{P} = \dfrac{W}{\Delta t}$ :
\[ W = \mathcal{P} \times \Delta t = 1885 \times 600 \approx \SI{1.13e6}{J} \approx \SI{1.1}{MJ} \]
Vérification par la rotation : l'angle balayé est $\theta = \omega\,\Delta t = 50\pi \times 600 = 30000\pi$ rad, d'où $W = \mathcal{M}\,\theta = 12 \times 30000\pi \approx \SI{1.13e6}{J}$ : même résultat, la cohérence est vérifiée.

\textbf{Conclusion} : $\omega \approx \SI{157}{rad.s^{-1}}$, $\mathcal{P} \approx \SI{1.9}{kW}$ et $W \approx \SI{1.1}{MJ}$ pour 10 minutes de mouture.
\end{exoresolu}

\begin{savaistu}[L'énergie hydraulique au Cameroun : de la force à la puissance]
Les barrages de \textbf{Lom Pangar} (réservoir de régulation), \textbf{Nachtigal} (au fil de l'eau, $\approx \SI{420}{MW}$) et \textbf{Edéa} exploitent le travail du poids de l'eau. L'eau chute d'une hauteur $h$ (dénivellation) : le travail moteur du poids par seconde est la \textbf{puissance hydraulique brute} $\mathcal{P}_h = \rho\,Q\,g\,h$ ($Q$ débit en $\si{m^3.s^{-1}}$). Les turbines transforment ce travail mécanique en travail électrique (alternateurs). Le rendement global (hydraulique $\to$ mécanique $\to$ électrique) dépasse souvent $90\%$. C'est la physique du travail et de la puissance qui éclaire nos foyers via ENEO/SONATREL !
\end{savaistu}

\begin{exercice}[Entraînement : Le freinage d'un vélo]
Un cycliste de masse totale $m = \SI{80}{kg}$ roule à $v_0 = \SI{20}{km.h^{-1}}$ sur une route plane. Il freine jusqu'à l'arrêt sur une distance $d = \SI{15}{m}$. On suppose la force de freinage $\vec{F}_f$ constante, horizontale et opposée au mouvement. Frottements autres négligés. $g = \SI{9.8}{N.kg^{-1}}$.
\begin{enumerate}
\item Exprimer le travail de la force de freinage $W(\vec{F}_f)$ en fonction de $F_f$ et $d$.
\item Appliquer le théorème de l'énergie cinétique (rappel : $\Delta E_c = W_{\text{ext}}$) pour déterminer $F_f$.
\item Calculer la puissance moyenne développée par la force de freinage si le freinage dure $\Delta t = \SI{4.0}{s}$.
\end{enumerate}
\end{exercice}

\begin{corrige}[Exercice : Le freinage d'un vélo]
\textbf{1.} $\vec{F}_f$ opposée au déplacement $\Rightarrow \alpha = 180°$, $\cos\alpha = -1$.
$W(\vec{F}_f) = -F_f \times d$.

\textbf{2.} $\Delta E_c = \dfrac{1}{2}m v_f^2 - \dfrac{1}{2}m v_0^2 = 0 - \dfrac{1}{2}m v_0^2$.
$v_0 = \SI{20}{km.h^{-1}} = \dfrac{20}{3.6} \approx \SI{5.56}{m.s^{-1}}$.
$\Delta E_c = -\dfrac{1}{2} \times 80 \times (5.56)^2 \approx \SI{-1235}{J}$.
Théorème : $W(\vec{F}_f) = \Delta E_c \Rightarrow -F_f \times 15 = -1235 \Rightarrow F_f \approx \SI{82.3}{N}$.

\textbf{3.} $\mathcal{P}_{\text{moy}} = \dfrac{W}{\Delta t} = \dfrac{-1235}{4.0} \approx \SI{-309}{W}$.
La puissance est négative (résistante), le frein dissipe l'énergie cinétique en chaleur.
\end{corrige}

\begin{attention}[Les 5 erreurs qui coûtent des points au Probatoire]
\begin{enumerate}
\item \textbf{Oublier le $\cos\alpha$ ou mal repérer l'angle.} L'angle $\alpha$ de la formule $W = F \cdot AB \cdot \cos\alpha$ est l'angle entre la \emph{force} et le \emph{déplacement}, pas forcément l'angle donné dans l'énoncé (par exemple angle avec la verticale). Fais toujours un schéma et justifie la valeur de $\alpha$.
\item \textbf{Négliger les conversions d'unités.} Les distances doivent être en \emph{mètres} (attention aux cm et km), les durées en \emph{secondes} (attention aux minutes et heures), les angles de rotation en \emph{radians} (un tour $= 2\pi$ rad, pas 360 !). Une valeur numérique juste avec une mauvaise unité est sanctionnée.
\item \textbf{Appliquer $W = F \cdot AB \cdot \cos\alpha$ à une force non constante.} Cette formule exige une force \emph{constante}. Pour la tension d'un ressort, qui varie avec l'allongement, il faut obligatoirement utiliser $W = \dfrac{1}{2}k\left(x_1^2 - x_2^2\right)$.
\item \textbf{Confondre travail du poids et longueur du trajet.} Le poids est conservatif : $W_{AB}(\vec{P}) = mg(z_A - z_B)$ ne dépend que de la \emph{dénivellation}. Écrire $W = mg \times L$ avec la longueur du chemin est faux, sauf si le déplacement est vertical. Attention aussi au signe : montée $\Rightarrow$ travail résistant.
\item \textbf{Confondre puissance et énergie, et mélanger les formules.} $\mathcal{P} = W/\Delta t$ donne des \emph{watts}, pas des joules ; le kilowattheure (kWh) est une \emph{énergie} ($\SI{1}{kWh} = \SI{3.6e6}{J}$), pas une puissance. De même, $\mathcal{P} = \mathcal{M}\,\omega$ exige $\omega$ en $\si{rad.s^{-1}}$ : ne jamais injecter directement une valeur en tr/min.
\end{enumerate}
\end{attention}

\newpage

\coursec{Exercices}

\courssub{Je vérifie mes connaissances}

\begin{exercice}[QCM : travail d'une force]
Pour chaque affirmation, choisir la (ou les) bonne(s) réponse(s) en justifiant brièvement.
\begin{enumerate}
\item Le travail d'une force constante $\vec{F}$ dont le point d'application se déplace de $A$ en $B$ s'exprime par :
\begin{enumerate}
\item $W_{AB}(\vec{F}) = F \times AB$ dans tous les cas ;
\item $W_{AB}(\vec{F}) = \vec{F} \cdot \overrightarrow{AB} = F \times AB \times \cos\alpha$ ;
\item $W_{AB}(\vec{F}) = F \times AB \times \sin\alpha$.
\end{enumerate}
\item L'unité SI du travail est :
\begin{enumerate}
\item le newton (N) ; \item le joule (J) ; \item le watt (W).
\end{enumerate}
\item Le travail d'une force est dit \cle{résistant} lorsque :
\begin{enumerate}
\item $\alpha < 90^{\circ}$ ; \item $\alpha = 90^{\circ}$ ; \item $90^{\circ} < \alpha \leq 180^{\circ}$.
\end{enumerate}
\item Une force dont le point d'application se déplace perpendiculairement à sa direction :
\begin{enumerate}
\item fournit un travail moteur ; \item ne travaille pas ; \item fournit un travail résistant.
\end{enumerate}
\end{enumerate}
\end{exercice}

\begin{exercice}[Vrai ou faux ? Justifier]
Répondre par vrai ou faux en justifiant chaque réponse par une propriété du cours.
\begin{enumerate}
\item Le travail du poids d'un corps dépend du chemin suivi entre le point de départ et le point d'arrivée.
\item Le poids est une force conservative.
\item Le travail de la tension d'un ressort à réponse linéaire, lorsque son allongement passe de $x_1$ à $x_2$, est $W = \dfrac{1}{2}k\left(x_2^2 - x_1^2\right)$.
\item La puissance moyenne d'une force s'exprime par $P = \dfrac{W}{\Delta t}$ et s'exprime en watts.
\item Une force de moment constant $\mathcal{M}$ par rapport à un axe $\Delta$, dont le point d'application tourne d'un angle $\theta$ (en radians) autour de $\Delta$, fournit un travail $W = \mathcal{M} \times \theta$.
\end{enumerate}
\end{exercice}

\begin{exercice}[Définitions et unités]
\begin{enumerate}
\item Définir le travail d'une force constante. Préciser la signification de chaque terme de la formule.
\item Qu'appelle-t-on force conservative ? Citer deux exemples de forces conservatives et un exemple de force non conservative.
\item Donner l'unité SI du travail et celle de la puissance, ainsi que leur expression en unités de base (kg, m, s).
\item Rappeler la relation entre la puissance, le moment d'une force et la vitesse angulaire de rotation. Préciser les unités de chaque grandeur.
\end{enumerate}
\end{exercice}

\begin{exercice}[Calculs directs]
\begin{enumerate}
\item Une force constante de norme $F = \SI{50}{N}$ déplace son point d'application de $d = \SI{8,0}{m}$ dans sa propre direction et dans son sens. Calculer son travail.
\item Un solide de masse $m = \SI{2,5}{kg}$ descend verticalement de $h = \SI{3,0}{m}$. Calculer le travail du poids. On donne $g = \SI{10}{N.kg^{-1}}$.
\item Un moteur fournit un travail $W = \SI{1,8e4}{J}$ en $\Delta t = \SI{30}{s}$. Calculer sa puissance moyenne.
\item Un arbre de machine tourne à $N = \SI{1500}{tr.min^{-1}}$ sous l'action d'un couple de moment $\mathcal{M} = \SI{12}{N.m}$. Calculer la puissance développée.
\end{enumerate}
\end{exercice}

\courssub{Je m'entraîne}

\begin{exercice}[La charrette du charretier]
Un charretier tire sa charrette sur une route horizontale de Mbalmayo en exerçant sur la brancard une force constante de norme $F = \SI{80}{N}$, faisant un angle $\alpha = 25^{\circ}$ avec l'horizontale. Le déplacement a une longueur $AB = \SI{200}{m}$.
\begin{enumerate}
\item Faire un schéma de la situation en représentant la force $\vec{F}$ et le déplacement $\overrightarrow{AB}$.
\item Calculer le travail de la force $\vec{F}$ sur ce trajet. Préciser sa nature (moteur, résistant ou nul).
\item Le trajet est parcouru en $\Delta t = 2~\text{min}~30~\text{s}$. Calculer la puissance moyenne développée par le charretier.
\item Les forces de frottement sur la charrette sont équivalentes à une force unique $\vec{f}$ horizontale, de norme $f = \SI{40}{N}$, opposée au déplacement. Calculer le travail de $\vec{f}$ et conclure sur sa nature.
\end{enumerate}
On donne : $\cos 25^{\circ} \approx 0,906$.
\end{exercice}

\begin{exercice}[L'ascenseur d'un immeuble de Yaoundé]
La cabine d'un ascenseur, de masse $m = \SI{800}{kg}$, monte de $h = \SI{15}{m}$ à vitesse constante. On donne $g = \SI{9,8}{N.kg^{-1}}$.
\begin{enumerate}
\item Faire le bilan des forces extérieures s'exerçant sur la cabine et les représenter sur un schéma.
\item Calculer le travail du poids de la cabine pendant la montée. Quelle est sa nature ?
\item Justifier que la force de traction $\vec{T}$ du câble a une norme égale à celle du poids, puis calculer le travail de $\vec{T}$.
\item La montée dure $\Delta t = \SI{20}{s}$. Calculer la puissance moyenne du moteur de l'ascenseur.
\end{enumerate}
\end{exercice}

\begin{exercice}[Compression d'un ressort]
Un ressort à spires non jointives, de raideur $k = \SI{400}{N.m^{-1}}$, est comprimé ; sa déformation (algébrique) passe de $x_1 = \SI{5,0}{cm}$ à $x_2 = \SI{12}{cm}$ en valeur absolue.
\begin{enumerate}
\item Rappeler l'expression du travail de la tension d'un ressort à réponse linéaire lorsque sa déformation passe de $x_1$ à $x_2$.
\item Calculer le travail de la tension du ressort lors de cette compression. Quelle est sa nature ? Justifier le signe.
\item En déduire le travail de la force exercée par l'opérateur qui comprime le ressort, en supposant la déformation quasi statique.
\item La tension d'un ressort est-elle une force conservative ? Justifier la réponse.
\end{enumerate}
\end{exercice}

\begin{exercice}[Serrage d'un écrou et couple d'une moto]
\textbf{Partie A.} Un ouvrier serre un écrou à l'aide d'une clé de longueur $d = \SI{30}{cm}$ en exerçant une force de norme $F = \SI{60}{N}$, perpendiculaire à la clé.
\begin{enumerate}
\item Calculer le moment de la force $\vec{F}$ par rapport à l'axe de l'écrou.
\item La clé tourne de $\dfrac{3}{4}$ de tour. Convertir cet angle en radians, puis calculer le travail de la force $\vec{F}$ pendant cette rotation.
\end{enumerate}
\textbf{Partie B.} Le moteur d'une moto-taxi développe une puissance $P = \SI{8,0}{kW}$ lorsque l'arbre tourne à $N = \SI{3000}{tr.min^{-1}}$.
\begin{enumerate}
\setcounter{enumi}{2}
\item Convertir la vitesse de rotation en $\si{rad.s^{-1}}$, puis calculer le moment $\mathcal{M}$ du couple moteur.
\item À puissance constante, que devient le moment du couple lorsque le régime atteint $N' = \SI{6000}{tr.min^{-1}}$ ? Commenter le résultat.
\end{enumerate}
\end{exercice}

\courssub{Je me prépare au Probatoire}

\begin{exercice}[Le treuil de chantier]
Sur un chantier de construction à Douala, un ouvrier utilise un treuil pour monter une charge. Le tambour du treuil a un rayon $r = \SI{15}{cm}$ ; il est actionné par une manivelle de longueur $L = \SI{50}{cm}$, sur l'extrémité de laquelle l'ouvrier exerce une force $\vec{F}$ constante, de norme $F = \SI{100}{N}$, perpendiculaire à la manivelle. La charge, de masse $m = \SI{40}{kg}$, est montée d'une hauteur $h = \SI{6,0}{m}$ à vitesse constante. On donne $g = \SI{10}{N.kg^{-1}}$ et $\pi \approx 3,14$.
\begin{enumerate}
\item Faire un schéma du dispositif (tambour, manivelle, câble, charge) en y reportant les données.
\item Montrer que le nombre $n$ de tours de manivelle nécessaires pour monter la charge de la hauteur $h$ vérifie $h = 2\pi r\, n$. Calculer $n$.
\item Calculer le moment de la force $\vec{F}$ par rapport à l'axe du tambour.
\item En déduire le travail $W_F$ fourni par l'ouvrier pendant la manœuvre.
\item Calculer le travail du poids de la charge pendant la montée, puis le travail utile $W_u$ fourni à la charge.
\item Définir le rendement $\eta$ de l'opération et le calculer. Commenter la valeur obtenue.
\item La manœuvre dure $\Delta t = \SI{1}{min}$. Calculer la puissance moyenne développée par l'ouvrier.
\end{enumerate}
\end{exercice}

\begin{exercice}[Remontée d'un skieur... version camerounaise : le télébenne d'une plantation]
Dans une plantation de bananes de la région du Littoral, une benne chargée de régimes, de masse totale $m = \SI{250}{kg}$, est tractée à vitesse constante le long d'un plan incliné d'un angle $\alpha = 20^{\circ}$ sur l'horizontale, sur une distance $AB = \SI{80}{m}$. Le câble de traction exerce une force $\vec{T}$ parallèle au plan incliné. Les frottements sont équivalents à une force unique $\vec{f}$, parallèle au plan, de norme $f = \SI{300}{N}$. On donne $g = \SI{10}{N.kg^{-1}}$, $\sin 20^{\circ} \approx 0,342$ et $\cos 20^{\circ} \approx 0,940$.
\begin{enumerate}
\item Faire le bilan des forces extérieures appliquées à la benne et les représenter sur un schéma soigné.
\item Calculer le travail du poids de la benne entre $A$ et $B$. Quelle est sa nature ?
\item Calculer le travail de la force de frottement $\vec{f}$ entre $A$ et $B$.
\item Le mouvement étant rectiligne et uniforme, montrer que $T = mg\sin\alpha + f$. Calculer $T$.
\item En déduire le travail de la force de traction $\vec{T}$ entre $A$ et $B$.
\item Vérifier que la somme des travaux de toutes les forces est nulle. Ce résultat était-il prévisible ? Justifier.
\item Le trajet dure $\Delta t = \SI{2}{min}$. Calculer la puissance moyenne du moteur qui actionne le câble.
\end{enumerate}
\end{exercice}

\begin{exercice}[Le forage et le treuil du village]
Pour équiper un village de l'Extrême-Nord en eau potable, une sonde de forage est remontée à l'aide d'un treuil électrique. La sonde et son équipement ont une masse totale $m = \SI{120}{kg}$. Le tambour du treuil, de rayon $r = \SI{10}{cm}$, est entraîné par un moteur électrique qui fournit un couple de moment constant $\mathcal{M} = \SI{150}{N.m}$. La sonde est remontée d'une profondeur $h = \SI{25}{m}$ à vitesse constante. On donne $g = \SI{10}{N.kg^{-1}}$ et $\pi \approx 3,14$.
\begin{enumerate}
\item Calculer le poids de la sonde et en déduire le travail utile nécessaire à sa remontée.
\item Calculer l'angle total $\theta$ (en radians) dont doit tourner le tambour pendant la remontée.
\item Calculer le travail $W_{\mathcal{M}}$ fourni par le couple moteur pendant cette rotation.
\item Comparer $W_{\mathcal{M}}$ au travail utile. Déterminer le rendement global du treuil et proposer une cause physique des pertes.
\item Le tambour tourne à la vitesse constante $N = \SI{60}{tr.min^{-1}}$. Convertir cette vitesse en $\si{rad.s^{-1}}$ et calculer la puissance du moteur.
\item En déduire la durée de la remontée, puis la vitesse verticale de la sonde en $\si{m.s^{-1}}$.
\item Le moteur est alimenté par un groupe électrogène. Sachant que le rendement du groupe est de $30\,\%$, calculer l'énergie thermique que doit fournir le carburant pour cette remontée.
\end{enumerate}
\end{exercice}

\newpage

\coursec{Corrigés des exercices}

\begin{corrige}[Exercice 1 — QCM : travail d'une force]
\begin{enumerate}
\item \textbf{Bonne réponse : b.} Le travail d'une force constante est le produit scalaire $W_{AB}(\vec{F}) = \vec{F}\cdot\overrightarrow{AB} = F\times AB\times\cos\alpha$, où $\alpha$ est l'angle orienté entre $\vec{F}$ et $\overrightarrow{AB}$. La réponse a n'est vraie que si $\vec{F}$ et $\overrightarrow{AB}$ sont colinéaires et de même sens ; la réponse c fait intervenir $\sin\alpha$, ce qui est faux.
\item \textbf{Bonne réponse : b.} L'unité SI du travail est le \cle{joule} (symbole J). Le newton (N) est l'unité de force, le watt (W) celle de la puissance.
\item \textbf{Bonne réponse : c.} Le travail est résistant lorsque $\cos\alpha < 0$, c'est-à-dire pour $90^{\circ} < \alpha \leq 180^{\circ}$ : la force s'oppose au déplacement.
\item \textbf{Bonne réponse : b.} Si le déplacement est perpendiculaire à la force, $\alpha = 90^{\circ}$ et $\cos 90^{\circ} = 0$ : la force \textbf{ne travaille pas} (exemple : la réaction normale du support sur un mobile en mouvement horizontal).
\end{enumerate}
\end{corrige}

\begin{corrige}[Exercice 2 — Vrai ou faux ?]
\begin{enumerate}
\item \textbf{Faux.} Le travail du poids ne dépend que de l'altitude des points de départ et d'arrivée : $W_{AB}(\vec{P}) = mg(z_A - z_B)$. Il est indépendant du chemin suivi.
\item \textbf{Vrai.} Le poids est une force conservative : son travail entre deux points ne dépend pas du chemin suivi (et il est nul sur tout trajet fermé).
\item \textbf{Faux.} Le travail de la tension d'un ressort à réponse linéaire est $W = \dfrac{1}{2}k\left(x_1^2 - x_2^2\right)$ : c'est la déformation \emph{initiale} qui porte l'indice 1 dans le premier terme. La formule proposée donne l'opposé du bon résultat.
\item \textbf{Vrai.} La puissance moyenne est $P_m = \dfrac{W}{\Delta t}$, avec $W$ en joules et $\Delta t$ en secondes ; $P$ s'exprime alors en watts (W).
\item \textbf{Vrai.} Pour une force de moment $\mathcal{M}$ constant par rapport à un axe $\Delta$, le travail lors d'une rotation d'angle $\theta$ (en radians) est $W = \mathcal{M}\,\theta$ (avec $\theta$ algébrique : le signe de $W$ dépend du sens de rotation).
\end{enumerate}
\end{corrige}

\begin{corrige}[Exercice 3 — Définitions et unités]
\begin{enumerate}
\item Le \cle{travail} d'une force constante $\vec{F}$ dont le point d'application se déplace de $A$ en $B$ est le produit scalaire :
\[ W_{AB}(\vec{F}) = \vec{F}\cdot\overrightarrow{AB} = F\times AB\times\cos\alpha \]
où $F$ est la norme de la force (en N), $AB$ la longueur du déplacement (en m) et $\alpha$ l'angle entre $\vec{F}$ et $\overrightarrow{AB}$.
\item Une force est dite \cle{conservative} lorsque son travail entre deux points $A$ et $B$ ne dépend pas du chemin suivi (de façon équivalente, son travail est nul sur tout trajet fermé). Exemples : le \textbf{poids}, la \textbf{tension d'un ressort} à réponse linéaire. Contre-exemple : les \textbf{forces de frottement}, dont le travail dépend de la longueur du trajet — ce sont des forces non conservatives.
\item Travail : le joule, $\SI{1}{J} = \SI{1}{N.m} = \SI{1}{kg.m^2.s^{-2}}$. Puissance : le watt, $\SI{1}{W} = \SI{1}{J.s^{-1}} = \SI{1}{kg.m^2.s^{-3}}$.
\item Pour une force de moment $\mathcal{M}$ par rapport à un axe de rotation tournant à la vitesse angulaire $\omega$ :
\[ P = \mathcal{M}\times\omega \]
avec $P$ en watts (W), $\mathcal{M}$ en newtons-mètres ($\si{N.m}$) et $\omega$ en radians par seconde ($\si{rad.s^{-1}}$).
\end{enumerate}
\end{corrige}

\begin{corrige}[Exercice 4 — Calculs directs]
\begin{enumerate}
\item Force et déplacement colinéaires et de même sens ($\alpha = 0$) :
\[ W = F\times d = 50\times 8,0 = \SI{4,0e2}{J} \quad (\text{travail moteur}). \]
\item Le poids est moteur lors d'une descente ($g = \SI{10}{N.kg^{-1}}$) :
\[ W(\vec{P}) = mgh = 2,5\times 10\times 3,0 = \SI{75}{J}. \]
\item Puissance moyenne :
\[ P = \dfrac{W}{\Delta t} = \dfrac{\num{1,8e4}}{30} = \SI{6,0e2}{W}. \]
\item Vitesse angulaire : $\omega = \dfrac{2\pi N}{60} = \dfrac{2\pi\times 1500}{60} = 50\pi \approx \SI{157}{rad.s^{-1}}$. D'où :
\[ P = \mathcal{M}\,\omega = 12\times 50\pi = 600\pi \approx \SI{1,88e3}{W} \approx \SI{1,9}{kW}. \]
\end{enumerate}
\end{corrige}

\begin{corrige}[Exercice 5 — La charrette du charretier]
\begin{enumerate}
\item Schéma : la route est horizontale, $\overrightarrow{AB}$ horizontal, $\vec{F}$ inclinée de $\alpha = 25^{\circ}$ au-dessus de l'horizontale, appliquée au brancard.
\begin{popfigure}
\begin{tikzpicture}[scale=1, >=Stealth]
\draw[thick] (-0.5,0) -- (6.5,0);
\draw[fill=popcream] (0.5,0.15) rectangle (2.5,1.0);
\draw[fill=popdark] (0.9,0.15) circle (0.18);
\draw[fill=popdark] (2.1,0.15) circle (0.18);
\draw[->, very thick, popblue] (2.5,0.7) -- ++(2.2,0.93) node[above right]{$\vec{F}$};
\draw[->, very thick, poporange] (0.5,-0.35) -- (5.5,-0.35) node[midway, below]{$\overrightarrow{AB}$};
\draw (4.7,0.7) arc[start angle=0, end angle=23, radius=2.2];
\node at (4.35,0.95) {$\alpha$};
\draw[dashed, popmuted] (2.5,0.7) -- (5.6,0.7);
\end{tikzpicture}
\legende{Force de traction $\vec{F}$ inclinée de $\alpha$ sur le déplacement horizontal $\overrightarrow{AB}$.}
\end{popfigure}
\item Travail de $\vec{F}$ :
\[ W_{AB}(\vec{F}) = F\times AB\times\cos\alpha = 80\times 200\times 0,906 \approx \SI{1,45e4}{J}. \]
Comme $\alpha < 90^{\circ}$, $\cos\alpha > 0$ : le travail est \textbf{moteur}.
\item $\Delta t = 2~\text{min}~30~\text{s} = \SI{150}{s}$. Puissance moyenne :
\[ P_m = \dfrac{W}{\Delta t} = \dfrac{\num{14496}}{150} \approx \SI{96,6}{W} \approx \SI{97}{W}. \]
\item $\vec{f}$ est opposée au déplacement : $\alpha = 180^{\circ}$, $\cos 180^{\circ} = -1$ :
\[ W_{AB}(\vec{f}) = -f\times AB = -40\times 200 = \SI{-8,0e3}{J}. \]
Le travail est \textbf{résistant} : les frottements s'opposent au mouvement.
\end{enumerate}
\end{corrige}

\begin{corrige}[Exercice 6 — L'ascenseur d'un immeuble de Yaoundé]
\begin{enumerate}
\item Bilan des forces extérieures sur la cabine : le \textbf{poids} $\vec{P}$ (vertical, vers le bas, appliqué au centre de gravité) et la \textbf{traction du câble} $\vec{T}$ (verticale, vers le haut). Sur le schéma, les deux vecteurs sont verticaux, de même longueur (vitesse constante), de sens opposés.
\begin{popfigure}
\begin{tikzpicture}[scale=0.8, >=Stealth]
\draw[thick, fill=popcream] (2,0) rectangle (4,3);
\draw[->, very thick, poppink] (3,1.5) -- ++(0,-2) node[below]{$\vec{P}=m\vec{g}$};
\draw[->, very thick, popblue] (3,1.5) -- ++(0,2) node[above]{$\vec{T}$};
\draw[dashed, popmuted] (3,3) -- (3,4) node[above]{Montée};
\end{tikzpicture}
\legende{Bilan des forces sur la cabine d'ascenseur en montée à vitesse constante.}
\end{popfigure}
\item Le poids est une force conservative ; la cabine \emph{monte} de $h$ ($g = \SI{9,81}{N.kg^{-1}}$) :
\[ W(\vec{P}) = -mgh = -800\times 9,81\times 15 = \SI{-1,1772e5}{J} \approx \SI{-1,18e5}{J}. \]
Le travail du poids est \textbf{résistant} (le poids s'oppose à la montée).
\item Le mouvement est rectiligne et uniforme : d'après la première loi de Newton (principe d'inertie), les forces se compensent : $\vec{T} + \vec{P} = \vec{0}$, donc $T = P = mg = 800\times 9,81 = \SI{7848}{N}$. La traction est dans le sens du déplacement :
\[ W(\vec{T}) = T\times h = 7848\times 15 = \SI{1,1772e5}{J} \approx \SI{1,18e5}{J} \quad (\text{travail moteur}). \]
On vérifie que $W(\vec{T}) = -W(\vec{P})$ : la somme des travaux est nulle, cohérent avec la vitesse constante.
\item Puissance moyenne du moteur :
\[ P_m = \dfrac{W(\vec{T})}{\Delta t} = \dfrac{\num{117720}}{20} = \SI{5886}{W} \approx \SI{5,9}{kW}. \]
\end{enumerate}
\end{corrige}

\begin{corrige}[Exercice 7 — Compression d'un ressort]
\begin{enumerate}
\item Pour un ressort à réponse linéaire de raideur $k$, lorsque la déformation algébrique passe de $x_1$ à $x_2$ :
\[ W = \dfrac{1}{2}k\left(x_1^2 - x_2^2\right). \]
\item Application numérique avec $x_1 = \SI{0,050}{m}$ et $x_2 = \SI{0,12}{m}$ :
\[ W = \dfrac{1}{2}\times 400\times\left(0,050^2 - 0,12^2\right) = 200\times\left(\num{2,5e-3} - \num{1,44e-2}\right) = 200\times(-\num{1,19e-2}) \approx \SI{-2,38}{J} \approx \SI{-2,4}{J}. \]
Le travail est \textbf{résistant} : lors de la compression, la tension du ressort est dirigée vers l'extérieur (elle repousse l'extrémité), donc \emph{opposée au déplacement} de son point d'application ; le signe négatif est cohérent.
\item En déformation quasi statique, la force de l'opérateur $\vec{F}_{\text{op}}$ compense à chaque instant la tension : $\vec{F}_{\text{op}} = -\vec{T}$. Donc :
\[ W(\vec{F}_{\text{op}}) = -W(\vec{T}) \approx \SI{+2,4}{J} \quad (\text{travail moteur}). \]
\item \textbf{Oui}, la tension d'un ressort est une force conservative : son travail ne dépend que des déformations initiale et finale ($x_1$ et $x_2$), et non du chemin suivi ni de la façon dont on passe de l'une à l'autre. En particulier, sur un cycle fermé ($x_2 = x_1$), le travail est nul.
\end{enumerate}
\end{corrige}

\begin{corrige}[Exercice 8 — Serrage d'un écrou et couple d'une moto]
\textbf{Partie A.}
\begin{enumerate}
\item La force est perpendiculaire à la clé ; le bras de levier est $d = \SI{0,30}{m}$ :
\[ \mathcal{M}(\vec{F}) = F\times d = 60\times 0,30 = \SI{18}{N.m}. \]
\item Trois quarts de tour : $\theta = \dfrac{3}{4}\times 2\pi = \dfrac{3\pi}{2} \approx \SI{4,71}{rad}$. Le moment est constant et la force reste perpendiculaire à la clé pendant la rotation :
\[ W = \mathcal{M}\times\theta = 18\times \dfrac{3\pi}{2} = 27\pi \approx \SI{84,8}{J} \approx \SI{85}{J}. \]
\end{enumerate}
\textbf{Partie B.}
\begin{enumerate}
\setcounter{enumi}{2}
\item $\omega = \dfrac{2\pi N}{60} = \dfrac{2\pi\times 3000}{60} = 100\pi \approx \SI{314}{rad.s^{-1}}$. De $P = \mathcal{M}\omega$ :
\[ \mathcal{M} = \dfrac{P}{\omega} = \dfrac{\num{8000}}{100\pi} \approx \SI{25,5}{N.m}. \]
\item À $N' = \SI{6000}{tr.min^{-1}}$ : $\omega' = 200\pi \approx \SI{628}{rad.s^{-1}}$, d'où
\[ \mathcal{M}' = \dfrac{P}{\omega'} = \dfrac{\num{8000}}{200\pi} \approx \SI{12,7}{N.m}. \]
\textbf{Commentaire :} à puissance constante, le moment du couple est divisé par deux lorsque la vitesse de rotation double ($\mathcal{M} = P/\omega$). C'est pourquoi on démultiplie (boîte de vitesses) pour obtenir un couple élevé à faible régime, par exemple au démarrage de la moto-taxi en côte.
\end{enumerate}
\end{corrige}

\begin{corrige}[Exercice 9 — Le treuil de chantier]
\begin{enumerate}
\item Schéma : tambour cylindrique de rayon $r = \SI{0,15}{m}$, manivelle de longueur $L = \SI{0,50}{m}$ solidaire du tambour, force $\vec{F}$ perpendiculaire à la manivelle, câble enroulé sur le tambour supportant la charge de masse $m = \SI{40}{kg}$.
\begin{popfigure}
\begin{tikzpicture}[scale=1, >=Stealth]
\draw[thick] (0,0) circle (1.5);
\draw[thick] (0,0) circle (0.5);
\draw[->, very thick, popblue] (0,1.5) -- ++(1.5,0) node[right]{$\vec{F}$};
\draw[thick] (0,0.5) -- (0,3.5);
\draw[thick, fill=popdark] (-0.3,3.5) rectangle (0.3,4.0) node[right]{$m$};
\draw[<->, popmuted] (0.6,0) -- (0.6,1.5) node[midway, right]{$r$};
\draw[<->, popmuted] (0,1.6) -- (1.5,1.6) node[midway, above]{$L$};
\end{tikzpicture}
\legende{Schéma du treuil : manivelle (longueur $L$) et tambour (rayon $r$).}
\end{popfigure}
\item À chaque tour de manivelle, le tambour fait un tour et une longueur de câble égale à sa circonférence $2\pi r$ s'enroule : la charge monte de $2\pi r$. Pour $n$ tours :
\[ h = 2\pi r\, n \quad\Longrightarrow\quad n = \dfrac{h}{2\pi r} = \dfrac{6,0}{2\times 3,14\times 0,15} = \dfrac{6,0}{0,942} \approx 6,37~\text{tours}. \]
\item La force est perpendiculaire à la manivelle ; le bras de levier est $L$ :
\[ \mathcal{M}(\vec{F}) = F\times L = 100\times 0,50 = \SI{50}{N.m}. \]
\item L'angle total de rotation est $\theta = 2\pi n = 2\times 3,14\times 6,37 \approx \SI{40,0}{rad}$. Le moment étant constant :
\[ W_F = \mathcal{M}\times\theta = 50\times 40,0 = \SI{2,00e3}{J}. \]
\item Travail du poids pendant la montée ($g = \SI{10}{N.kg^{-1}}$) :
\[ W(\vec{P}) = -mgh = -40\times 10\times 6,0 = \SI{-2,4e3}{J}. \]
Le travail utile fourni à la charge (pour la monter à vitesse constante) est :
\[ W_u = mgh = \SI{2,4e3}{J}. \]
\item Le rendement est le rapport du travail utile au travail fourni :
\[ \eta = \dfrac{W_u}{W_F} = \dfrac{2400}{2000} = 1,2. \]
\textbf{Commentaire :} on obtient $\eta > 1$, ce qui est \textbf{physiquement impossible} : un dispositif ne peut pas restituer plus d'énergie qu'il n'en reçoit. Les données de l'énoncé sont donc incohérentes : la force $F = \SI{100}{N}$ est \emph{insuffisante} pour monter cette charge à vitesse constante. La force minimale nécessaire vérifierait $F\times L\times 2\pi n = mgh$, soit $F = \dfrac{mgr}{L} = \dfrac{40\times 10\times 0,15}{0,50} = \SI{120}{N}$. (Avec $F = \SI{120}{N}$, le rendement serait de $100\,\%$ en l'absence de frottements.)
\item Puissance moyenne développée par l'ouvrier :
\[ P_m = \dfrac{W_F}{\Delta t} = \dfrac{2000}{60} \approx \SI{33,3}{W} \approx \SI{33}{W}. \]
\end{enumerate}
\textbf{Barème indicatif (sur 10) :}
\begin{itemize}
\item Schéma : 1 pt
\item Relation $h = 2\pi r n$ et calcul de $n$ : 1,5 pt
\item Moment $\mathcal{M}$ : 1 pt
\item Travail $W_F$ : 1,5 pt
\item $W(\vec{P})$ et $W_u$ : 1,5 pt
\item Rendement et commentaire critique : 2 pt
\item Puissance moyenne : 1,5 pt
\end{itemize}
\textbf{Points de vigilance :}
\begin{itemize}
\item Convertir $r$ et $L$ en mètres ;
\item L'angle $\theta$ doit être en \emph{radians} dans $W = \mathcal{M}\theta$ ;
\item Un rendement supérieur à 1 doit toujours alerter et être commenté — le correcteur attend cette analyse critique.
\end{itemize}
\end{corrige}

\begin{corrige}[Exercice 10 — Le télébenne d'une plantation]
\begin{enumerate}
\item Bilan des forces extérieures sur la benne : le \textbf{poids} $\vec{P}$ (vertical, vers le bas), la \textbf{réaction} $\vec{R}$ du plan (perpendiculaire au plan), la \textbf{traction} $\vec{T}$ du câble (parallèle au plan, vers le haut) et les \textbf{frottements} $\vec{f}$ (parallèles au plan, vers le bas).
\begin{popfigure}
\begin{tikzpicture}[scale=1, >=Stealth]
\draw[thick] (0,0) -- (7,0);
\draw[thick] (0,0) -- (6,2.19);
\draw (1.6,0) arc[start angle=0, end angle=20, radius=1.6];
\node at (1.95,0.28) {$\alpha$};
\begin{scope}[rotate around={20:(3,1.095)}]
\draw[fill=popcream] (2.6,1.095) rectangle (3.8,1.9);
\end{scope}
\coordinate (G) at (3.2,1.6);
\draw[->, very thick, popblue] (G) -- ++(20:2.0) node[above right]{$\vec{T}$};
\draw[->, very thick, poporange] (G) -- ++(200:1.2) node[below left]{$\vec{f}$};
\draw[->, very thick, popgreen] (G) -- ++(110:1.4) node[above left]{$\vec{R}$};
\draw[->, very thick, poppink] (G) -- ++(-90:1.8) node[below]{$\vec{P}$};
\end{tikzpicture}
\legende{Bilan des forces sur la benne tractée le long du plan incliné ($\alpha = 20^{\circ}$).}
\end{popfigure}
\item Le poids est conservatif ; la dénivellation entre $A$ et $B$ est $h = AB\sin\alpha = 80\times 0,342 = \SI{27,36}{m}$. La benne \emph{monte} ($g = \SI{10}{N.kg^{-1}}$) :
\[ W_{AB}(\vec{P}) = -mgh = -mg\,AB\sin\alpha = -250\times 10\times 27,36 \approx \SI{-6,84e4}{J}. \]
Travail \textbf{résistant} (le poids s'oppose à la montée).
\item $\vec{f}$ est opposée au déplacement :
\[ W_{AB}(\vec{f}) = -f\times AB = -300\times 80 = \SI{-2,4e4}{J}. \]
\item Mouvement rectiligne uniforme : d'après le principe d'inertie, la somme des forces est nulle. En projetant sur l'axe du plan incliné (orienté vers le haut) :
\[ T - mg\sin\alpha - f = 0 \quad\Longrightarrow\quad T = mg\sin\alpha + f = 250\times 10\times 0,342 + 300 = 855 + 300 = \SI{1155}{N}. \]
\item $\vec{T}$ est dans le sens du déplacement :
\[ W_{AB}(\vec{T}) = T\times AB = 1155\times 80 = \SI{9,24e4}{J}. \]
\item La réaction $\vec{R}$, perpendiculaire au déplacement, ne travaille pas. Somme des travaux :
\[ \sum W = \num{92400} - \num{68400} - \num{24000} + 0 = \SI{0}{J}. \]
Ce résultat était \textbf{prévisible} : le mouvement est rectiligne et uniforme, donc les forces se compensent ($\sum\vec{F} = \vec{0}$), et la somme de leurs travaux $\left(\sum\vec{F}\right)\cdot\overrightarrow{AB}$ est nulle.
\item $\Delta t = \SI{120}{s}$. La puissance du moteur correspond au travail de la traction :
\[ P_m = \dfrac{W_{AB}(\vec{T})}{\Delta t} = \dfrac{\num{92400}}{120} = \SI{770}{W}. \]
\end{enumerate}
\textbf{Barème indicatif (sur 10) :}
\begin{itemize}
\item Bilan + schéma : 1,5 pt
\item $W(\vec{P})$ : 1,5 pt
\item $W(\vec{f})$ : 1 pt
\item Démonstration de $T$ : 2 pt
\item $W(\vec{T})$ : 1 pt
\item Somme nulle justifiée : 1,5 pt
\item Puissance : 1,5 pt
\end{itemize}
\textbf{Points de vigilance :}
\begin{itemize}
\item Citer le principe d'inertie pour justifier $T = mg\sin\alpha + f$ ;
\item Ne pas oublier que $\vec{R}$ ne travaille pas ;
\item La dénivellation est $AB\sin\alpha$ et non $AB$ ;
\item Préciser la nature (moteur/résistant) de chaque travail.
\end{itemize}
\end{corrige}

\begin{corrige}[Exercice 11 — Le forage et le treuil du village]
\begin{enumerate}
\item Poids de la sonde : $P = mg = 120\times 10 = \SI{1200}{N}$ ($g = \SI{10}{N.kg^{-1}}$). La remontée s'effectuant à vitesse constante, le travail utile est l'opposé du travail du poids :
\[ W_u = mgh = 1200\times 25 = \SI{3,0e4}{J}. \]
\item Le câble s'enroule sans glisser : la longueur remontée vaut $h = r\theta$, d'où
\[ \theta = \dfrac{h}{r} = \dfrac{25}{0,10} = \SI{250}{rad} \quad (\approx 39,8~\text{tours}). \]
\item Le moment du couple est constant :
\[ W_{\mathcal{M}} = \mathcal{M}\times\theta = 150\times 250 = \SI{3,75e4}{J}. \]
\item $W_{\mathcal{M}} = \SI{37,5}{kJ} > W_u = \SI{30}{kJ}$ : le moteur fournit plus que le strict nécessaire. Rendement global :
\[ \eta = \dfrac{W_u}{W_{\mathcal{M}}} = \dfrac{\num{30000}}{\num{37500}} = 0,80 = 80\,\%. \]
Cause physique des pertes : les \textbf{frottements} dans les roulements du tambour, l'engrenage de réduction, la flexion du câble (et l'échauffement du moteur), qui dissipent de l'énergie.
\item $\omega = \dfrac{2\pi N}{60} = \dfrac{2\times 3,14\times 60}{60} = \SI{6,28}{rad.s^{-1}}$. Puissance du moteur :
\[ P = \mathcal{M}\,\omega = 150\times 6,28 = \SI{942}{W} \approx \SI{0,94}{kW}. \]
\item Durée de la remontée :
\[ \Delta t = \dfrac{W_{\mathcal{M}}}{P} = \dfrac{\num{37500}}{942} \approx \SI{39,8}{s} \approx \SI{40}{s}. \]
Vitesse verticale de la sonde : $v = \dfrac{h}{\Delta t} = \dfrac{25}{39,8} \approx \SI{0,628}{m.s^{-1}}$. (Vérification : $v = r\omega = 0,10\times 6,28 = \SI{0,628}{m.s^{-1}}$ — cohérent.)
\item Le groupe électrogène a un rendement de $30\,\%$ : l'énergie électrique fournie au moteur vaut $W_{\mathcal{M}} = 0,30\times E_{\text{th}}$, d'où
\[ E_{\text{th}} = \dfrac{W_{\mathcal{M}}}{0,30} = \dfrac{\num{37500}}{0,30} = \SI{1,25e5}{J} = \SI{125}{kJ}. \]
\end{enumerate}
\textbf{Barème indicatif (sur 10) :}
\begin{itemize}
\item Poids et $W_u$ : 1,5 pt
\item Angle $\theta$ : 1,5 pt
\item $W_{\mathcal{M}}$ : 1 pt
\item Comparaison, rendement et cause des pertes : 2 pt
\item Conversion $\omega$ et puissance : 1,5 pt
\item Durée et vitesse : 1,5 pt
\item Énergie thermique : 1 pt
\end{itemize}
\textbf{Points de vigilance :}
\begin{itemize}
\item La relation $h = r\theta$ exige $\theta$ en radians et $r$ en mètres ;
\item Bien distinguer le travail \emph{fourni} par le moteur ($W_{\mathcal{M}}$) du travail \emph{utile} ($mgh$) ;
\item Pour la question 7, c'est l'énergie thermique qui est la plus grande (on divise par le rendement, on ne multiplie pas) ;
\item Enchaîner les rendements en les multipliant si plusieurs étages sont en cascade.
\end{itemize}
\end{corrige}

\newpage

\coursec{Fiche bilan}

\begin{popfigure}
\begin{tikzpicture}[mindmap, grow cyclic, every node/.style={concept, font=\small},
  root concept/.append style={concept color=popblue, font=\bfseries},
  level 1/.append style={level distance=3.6cm, sibling angle=60},
  level 2/.append style={level distance=2.2cm, font=\scriptsize}]
\node[root concept]{Travail et puissance d'une force}
  child[concept color=poporange]{ node {Travail d'une force constante}
    child { node {$W_{AB}(\vec{F})=\vec{F}\cdot\overrightarrow{AB}$} }
    child { node {$W=F\,AB\cos\alpha$} }
  }
  child[concept color=popgreen]{ node {Poids : force conservative}
    child { node {$W_{AB}(\vec{P})=mg(z_A-z_B)$} }
    child { node {indépendant du chemin} }
  }
  child[concept color=poppurple]{ node {Tension d'un ressort}
    child { node {$W=\tfrac12 k(x_A^2-x_B^2)$} }
    child { node {conservative} }
  }
  child[concept color=popgold]{ node {Travail et moment}
    child { node {$W=\mathcal{M}_{\Delta}(\vec{F})\,\theta$} }
    child { node {$\theta$ en rad} }
  }
  child[concept color=poppink]{ node {Puissance}
    child { node {$P=\dfrac{W}{\Delta t}$} }
    child { node {$P=\vec{F}\cdot\vec{v}$} }
  }
  child[concept color=popblue!70]{ node {Puissance en rotation}
    child { node {$P=\mathcal{M}_{\Delta}\,\omega$} }
    child { node {$\omega$ en rad/s} }
  };
\end{tikzpicture}
\legende{Carte mentale de la leçon : les six notions clés autour du travail et de la puissance.}
\end{popfigure}

\begin{bilan}
\textbf{Définitions clés}
\begin{itemize}
\item Le \cle{travail} d'une force constante $\vec{F}$ dont le point d'application se déplace de $A$ vers $B$ est le produit scalaire $W_{AB}(\vec{F})=\vec{F}\cdot\overrightarrow{AB}$. Unité : le \cle{joule} (J).
\item Le travail est \cle{moteur} si $W>0$ ($\alpha<90^{\circ}$), \cle{résistant} si $W<0$ ($\alpha>90^{\circ}$), \cle{nul} si $\vec{F}\perp\overrightarrow{AB}$.
\item Une force est \cle{conservative} si son travail entre deux points ne dépend pas du chemin suivi (travail nul sur un trajet fermé).
\item La \cle{puissance} est le travail fourni par unité de temps : $P=\dfrac{W}{\Delta t}$. Unité : le \cle{watt} (W), $1~\mathrm{W}=1~\mathrm{J/s}$.
\end{itemize}

\textbf{Formules à connaître par c\oe ur}
\begin{center}
\rowcolors{2}{popblueL}{white}
\begin{tabularx}{\linewidth}{|l|X|l|}
\hline
\rowcolor{popblue!40} \textbf{Grandeur} & \textbf{Formule} & \textbf{Unités SI} \\
\hline
Travail d'une force constante & $W_{AB}(\vec{F})=F\cdot AB\cdot\cos\alpha$ & J (N, m) \\
\hline
Travail du poids & $W_{AB}(\vec{P})=mg\,(z_A-z_B)$ & J \\
\hline
Travail de la tension d'un ressort & $W_{A\to B}(\vec{T})=\dfrac{1}{2}k\left(x_A^{2}-x_B^{2}\right)$ & J ($k$ en N/m) \\
\hline
Travail en rotation & $W=\mathcal{M}_{\Delta}(\vec{F})\cdot\theta$ & J (N$\cdot$m, rad) \\
\hline
Puissance moyenne & $P=\dfrac{W}{\Delta t}$ & W \\
\hline
Puissance (translation) & $P=\vec{F}\cdot\vec{v}=F\,v\cos\alpha$ & W (m/s) \\
\hline
Puissance en rotation & $P=\mathcal{M}_{\Delta}(\vec{F})\cdot\omega$ & W (rad/s) \\
\hline
\end{tabularx}
\end{center}

\textbf{Méthodes express}
\begin{itemize}
\item \textbf{Calculer un travail} : identifier la force, vérifier qu'elle est constante, repérer l'angle $\alpha$ entre $\vec{F}$ et le déplacement, appliquer $W=F\,AB\cos\alpha$ puis conclure sur le signe (moteur/résistant).
\item \textbf{Travail du poids} : mesurer la dénivellation $h=z_A-z_B$ ; le poids travaille positivement en descente, négativement en montée ; jamais besoin de connaître le chemin.
\item \textbf{Ressort} : repérer les allongements $x_A$ et $x_B$ par rapport à la position à vide, puis $W=\tfrac12 k(x_A^2-x_B^2)$.
\item \textbf{Rotation} : convertir l'angle en radians ($\theta=\dfrac{\pi}{180}\times$ degrés), puis $W=\mathcal{M}\theta$ et $P=\mathcal{M}\omega$.
\item \textbf{Puissance} : si le travail est connu, $P=W/\Delta t$ ; si la vitesse est connue, $P=Fv\cos\alpha$ ; en rotation, $P=\mathcal{M}\omega$.
\end{itemize}
\end{bilan}

\begin{aretenir}[Checklist avant le Probatoire]
\begin{itemize}
\item[$\square$] Je sais écrire le travail d'une force constante comme un produit scalaire et donner sa signe (moteur, résistant, nul).
\item[$\square$] Je vérifie toujours l'homogénéité : $[W]=\mathrm{N}\times\mathrm{m}=\mathrm{J}$ et $[P]=\mathrm{J/s}=\mathrm{W}$.
\item[$\square$] Je sais que le travail du poids ne dépend que de la dénivellation : $W=mg(z_A-z_B)$.
\item[$\square$] Je sais citer la définition d'une force conservative et donner des exemples (poids, tension d'un ressort) et un contre-exemple (frottements).
\item[$\square$] Je sais exprimer le travail de la tension d'un ressort : $W=\tfrac12 k(x_A^2-x_B^2)$, avec $x$ mesuré depuis la position à vide.
\item[$\square$] Je convertis systématiquement les angles en radians avant d'utiliser $W=\mathcal{M}\theta$.
\item[$\square$] Je sais relier travail et moment : $W=\mathcal{M}_{\Delta}(\vec{F})\,\theta$ pour un solide en rotation autour d'un axe fixe.
\item[$\square$] Je sais calculer une puissance moyenne avec $P=W/\Delta t$ en convertissant le temps en secondes.
\item[$\square$] Je sais passer de la puissance au travail : $W=P\,\Delta t$ (utile pour l'énergie consommée, ex. facture ENEO).
\item[$\square$] Je sais exprimer la puissance d'une force de moment constant : $P=\mathcal{M}\,\omega$, avec $\omega$ en rad/s.
\item[$\square$] Je convertis correctement : $1~\mathrm{tr/min}=\dfrac{2\pi}{60}~\mathrm{rad/s}$ et $1~\mathrm{kWh}=3{,}6\times10^{6}~\mathrm{J}$.
\item[$\square$] Je rédige mes applications numériques avec unité SI et un nombre raisonnable de chiffres significatifs.
\end{itemize}
\end{aretenir}

\findecours

\end{document}
